% !TEX root=./adl.tex

% --- Sketch scene for vision intro (house, tree, person) ------------
% Drawn with bottom-left at (0,0), size 2x2.
\newcommand{\visionSketchScene}{%
  \fill[blue!12] (0, 0.5) rectangle (2, 2);
  \fill[green!40!brown!50] (0, 0) rectangle (2, 0.5);
  \fill[yellow!80!orange] (1.65, 1.7) circle (0.18);
  \fill[brown!60!black] (0.28, 0.5) rectangle (0.42, 0.95);
  \fill[green!55!black] (0.05, 0.85) -- (0.65, 0.85) -- (0.35, 1.55) -- cycle;
  \fill[orange!50!brown!70] (0.85, 0.5) rectangle (1.45, 1.05);
  \fill[red!60!black] (0.78, 1.05) -- (1.52, 1.05) -- (1.15, 1.5) -- cycle;
  \fill[brown!80!black] (1.05, 0.5) rectangle (1.25, 0.85);
  \fill[yellow!60!orange] (1.78, 0.92) circle (0.07);
  \draw[thick] (1.78, 0.85) -- (1.78, 0.65);
  \draw[thick] (1.78, 0.78) -- (1.7, 0.7);
  \draw[thick] (1.78, 0.78) -- (1.86, 0.7);
  \draw[thick] (1.78, 0.65) -- (1.71, 0.52);
  \draw[thick] (1.78, 0.65) -- (1.85, 0.52);
}
% Alternate sketch scene B (mountain + lake), 2x2 at origin.
\newcommand{\visionSketchSceneB}{%
  \fill[orange!20] (0, 0.6) rectangle (2, 2);
  \fill[blue!35] (0, 0) rectangle (2, 0.6);
  \fill[brown!50!gray] (0.25, 0.6) -- (1.45, 0.6) -- (0.85, 1.7) -- cycle;
  \fill[white!88!brown] (0.62, 1.28) -- (1.08, 1.28) -- (0.85, 1.7) -- cycle;
  \fill[orange!90!red] (1.55, 1.55) circle (0.17);
  \draw[white!85!blue, thick] (0.05, 0.4) -- (0.4, 0.45) -- (0.75, 0.4) -- (1.1, 0.45) -- (1.45, 0.4) -- (1.8, 0.45);
  \draw[white!85!blue, thick] (0.1, 0.2) -- (0.45, 0.25) -- (0.8, 0.2) -- (1.15, 0.25) -- (1.5, 0.2) -- (1.85, 0.25);
}
% --- Segmentation intro scene: sky, grass, tree, three people (same class) ---
% bottom-left at (0,0), size 2x2. \segPerson{x}{headcolor}{bodycolor}
\newcommand{\segPerson}[3]{%
  \fill[#2] (#1,1.02) circle (0.085);                        % head
  \fill[#3] (#1-0.085,0.55) rectangle (#1+0.085,0.95);       % torso
  \draw[#3,line width=1.4pt] (#1-0.05,0.55)--(#1-0.12,0.18); % leg
  \draw[#3,line width=1.4pt] (#1+0.05,0.55)--(#1+0.12,0.18); % leg
  \draw[#3,line width=1.4pt] (#1-0.085,0.85)--(#1-0.2,0.6);  % arm
  \draw[#3,line width=1.4pt] (#1+0.085,0.85)--(#1+0.2,0.6);  % arm
}
\newcommand{\segSceneRaw}{%
  \fill[blue!18] (0,0.9) rectangle (2,2);            % sky
  \fill[green!45!brown!55] (0,0) rectangle (2,0.9);  % grass
  \fill[brown!60!black] (1.55,0.9) rectangle (1.67,1.4);                      % trunk
  \fill[green!50!black] (1.3,1.25) -- (1.92,1.25) -- (1.61,1.95) -- cycle;    % foliage
  \segPerson{0.4}{brown!70!red}{gray!70}
  \segPerson{0.85}{brown!70!red}{gray!60}
  \segPerson{1.25}{brown!70!red}{gray!75}
}
% the "stuff" classes drawn as flat label regions (sky, grass, tree)
\newcommand{\segStuff}{%
  \fill[blue!45] (0,0.9) rectangle (2,2);
  \fill[green!55!black] (0,0) rectangle (2,0.9);
  \fill[teal!70!black] (1.55,0.9) rectangle (1.67,1.4);
  \fill[teal!70!black] (1.3,1.25) -- (1.92,1.25) -- (1.61,1.95) -- cycle;
}
% --- Tracking scene helpers (street with pedestrians), 2x2 at origin ----
% Used by the TrackFormer figures so the same identities can move across frames.
\newcommand{\trackBG}{%
  \fill[blue!12] (0,0.85) rectangle (2,2);            % sky
  \fill[gray!28] (0,0) rectangle (2,0.85);            % pavement
  \fill[gray!45] (0,0.80) rectangle (2,0.85);         % curb
  \foreach \x in {0.1,0.5,0.9,1.3,1.7} \fill[white!80!gray] (\x,0.36) rectangle (\x+0.22,0.44);
}
\newcommand{\trackPerson}[2]{% #1 x, #2 color
  \fill[#2] (#1,0.78) circle (0.10);
  \fill[#2] (#1-0.09,0.30) rectangle (#1+0.09,0.65);
  \draw[#2,line width=1.0pt] (#1-0.05,0.30)--(#1-0.11,0.06);
  \draw[#2,line width=1.0pt] (#1+0.05,0.30)--(#1+0.11,0.06);
}
\newcommand{\trackBox}[2]{% #1 x (person center), #2 color
  \draw[#2,very thick] (#1-0.17,0.02) rectangle (#1+0.17,0.92);
}

% --- MOT-style scene (same imagery as the "Multi-Object Tracking Problem" slide) ---
% Stick-figure person pic, shared by the MOT slide and the TrackFormer figures.
\tikzset{
  motperson/.pic={
    \begin{scope}[line cap=round, line join=round]
      \fill (0,0.85) circle (0.13);                    % head
      \draw[line width=0.9pt] (0,0.72) -- (0,0.32);    % torso
      \draw[line width=0.9pt] (0,0.62) -- (-0.18,0.40); % arms
      \draw[line width=0.9pt] (0,0.62) -- ( 0.18,0.40);
      \draw[line width=0.9pt] (0,0.32) -- (-0.15,0.0); % legs
      \draw[line width=0.9pt] (0,0.32) -- ( 0.15,0.0);
    \end{scope}}
}
% Street-scene background filling (0,0) rectangle (2,2): sky, sun, clouds, skyline, trees, ground.
\newcommand{\motBG}{%
  \fill[blue!12] (0,0) rectangle (2,2);                 % sky
  \fill[yellow!55!orange!80] (1.68,1.72) circle (0.13); % sun
  \fill[white] (0.35,1.60) circle (0.10) (0.50,1.64) circle (0.12) (0.66,1.60) circle (0.10);
  \fill[white] (1.12,1.80) circle (0.08) (1.24,1.82) circle (0.10) (1.36,1.80) circle (0.08);
  \fill[gray!22] (0.18,0.45) rectangle (0.55,1.02);     % skyline
  \fill[gray!16] (0.62,0.45) rectangle (0.95,0.82);
  \fill[gray!30] (1.00,0.45) rectangle (1.30,1.10);
  \fill[gray!28] (1.45,0.45) rectangle (1.82,1.18);
  \foreach \wy in {0.58,0.74,0.90} {                    % windows
    \fill[yellow!40!gray!40] (0.26,\wy) rectangle (0.32,\wy+0.05);
    \fill[yellow!40!gray!40] (0.41,\wy) rectangle (0.47,\wy+0.05);}
  \foreach \wy in {0.58,0.78,0.98} {
    \fill[yellow!40!gray!40] (1.53,\wy) rectangle (1.59,\wy+0.05);
    \fill[yellow!40!gray!40] (1.68,\wy) rectangle (1.74,\wy+0.05);}
  \fill[brown!60!black] (0.06,0.30) rectangle (0.10,0.50);  % trees
  \fill[green!45!black] (0.08,0.55) circle (0.11);
  \fill[brown!60!black] (1.90,0.30) rectangle (1.94,0.50);
  \fill[green!45!black] (1.92,0.55) circle (0.10);
  \fill[green!10!gray!18] (0,0) rectangle (2,0.45);         % ground
  \draw[gray!45, line width=0.4pt] (0,0.45) -- (2,0.45);    % horizon
}
% person standing at scene-x #1 (feet on the pavement) in colour #2
\newcommand{\motFig}[2]{\pic[#2, scale=0.72] at (#1,0.20) {motperson};}
% coloured detection box around the person at scene-x #1 in colour #2 (object-query detection)
\newcommand{\motFigBox}[2]{\draw[very thick,#2] (#1-0.22,0.16) rectangle (#1+0.22,0.96);}
% track-query detection box: identity colour with the same violet diagonal hatch as the
% track-query boxes (continued track), so the two continued detections read as track queries
\newcommand{\motFigBoxTrack}[2]{%
  \draw[very thick,#2, postaction={pattern=north east lines, pattern color=violet!65}]
       (#1-0.22,0.16) rectangle (#1+0.22,0.96);}

% Thumbnail of one 0.5x0.5 patch from the sketch scene.
% #1: thumb x, #2: thumb y, #3: scene col (0..3), #4: scene row (0..3 bottom up)
\newcommand{\visionPatchThumb}[4]{%
  \pgfmathsetmacro{\vptpx}{0.5*#3}
  \pgfmathsetmacro{\vptpy}{0.5*#4}
  \pgfmathsetmacro{\vptsx}{#1 - \vptpx}
  \pgfmathsetmacro{\vptsy}{#2 - \vptpy}
  \begin{scope}
    \clip (#1, #2) rectangle (#1+0.5, #2+0.5);
    \begin{scope}[shift={(\vptsx, \vptsy)}]
      \visionSketchScene
    \end{scope}
  \end{scope}
  \draw[thick] (#1, #2) rectangle (#1+0.5, #2+0.5);
}

% Draw scene A clipped to its 2x2 frame, with optional transform options.
% #1 = transform options applied to the scene before clipping
\newcommand{\sceneAframed}[1]{%
  \begin{scope}
    \clip (0,0) rectangle (2,2);
    \begin{scope}[#1]\visionSketchScene\end{scope}
  \end{scope}
  \draw[thick] (0,0) rectangle (2,2);
}

% 1x1 thumbnail of a scene with a transform (for repeated-augmentation figure).
% #1 = x, #2 = y, #3 = transform options, #4 = scene macro
\newcommand{\rathumb}[4]{%
  \begin{scope}[shift={(#1,#2)}]
    \begin{scope}\clip (0,0) rectangle (1,1);\fill[gray!20] (0,0) rectangle (1,1);
      \begin{scope}[scale=0.5]\begin{scope}[#3]#4\end{scope}\end{scope}
    \end{scope}
    \draw[thick] (0,0) rectangle (1,1);
  \end{scope}
}

% Distillation-token marker (DeiT), self-documented in the slide footnote.
\newcommand{\distill}{\ensuremath{^{\scriptstyle\triangle}}}

% -----------------------------------------------------------------------
\section{Vision Models}

% =======================================================================
\subsection{Overview}

% -----------------------------------------------------------------------
\begin{frame}[t]{Vision Models: From NLP Transformers to Images}

  \textbf{Goal}: Transformer-style architectures for image classification, detection, segmentation, tracking, $\ldots$

  \vspace{0.1em}
  \begin{center}
    \begin{tikzpicture}[
      scale=0.6, every node/.style={scale=0.6},
      tok/.style={draw, rounded corners=1pt, minimum width=0.7cm, minimum height=0.35cm, font=\tiny},
      arr/.style={-Stealth, thick},
    ]
      % --- Text side ---
      \node[font=\small\bfseries] at (-4.5, 2.4) {NLP};
      \node[font=\scriptsize] at (-4.5, 1.7) {``The cat sat on the mat''};
      \draw[arr] (-4.5, 1.4) -- (-4.5, 0.9);
      \node[tok, fill=blue!12] at (-6.3, 0.5) {The};
      \node[tok, fill=blue!12] at (-5.5, 0.5) {cat};
      \node[tok, fill=blue!12] at (-4.7, 0.5) {sat};
      \node[font=\scriptsize] at (-3.95, 0.5) {$\cdots$};
      \node[tok, fill=blue!12] at (-3.2, 0.5) {mat};
      \node[font=\scriptsize, text=gray, align=center] at (-4.5, -0.2) {1D sequence,\\discrete vocabulary};

      % --- Image side ---
      \node[font=\small\bfseries] at (1.5, 2.4) {Vision};

      % Scene (bottom-left at (0.5, 0), 2x2)
      \begin{scope}[shift={(0.5, 0)}]
        \begin{scope}
          \clip (0,0) rectangle (2,2);
          \visionSketchScene
        \end{scope}
      \end{scope}

      % 4x4 patch grid
      \draw[very thick] (0.5, 0) rectangle (2.5, 2);
      \foreach \i in {1,2,3} {
        \pgfmathsetmacro{\xpos}{0.5 + 0.5*\i}
        \pgfmathsetmacro{\ypos}{0.5*\i}
        \draw[thick, gray!70] (\xpos, 0) -- (\xpos, 2);
        \draw[thick, gray!70] (0.5, \ypos) -- (2.5, \ypos);
      }

      % Arrow from image to sequence
      \draw[arr] (2.7, 1.0) -- (3.4, 1.0);
      \node[font=\tiny, above] at (3.05, 1.0) {patchify};

      % Patch sequence on the right (thumbnails of actual scene patches)
      \def\seqx{3.7}
      \visionPatchThumb{\seqx}{1.85}{0}{3}
      \node[font=\tiny, left] at (\seqx, 2.1) {$p_1$};
      \visionPatchThumb{\seqx}{1.25}{2}{3}
      \node[font=\tiny, left] at (\seqx, 1.5) {$p_3$};
      \node[font=\tiny] at (3.95, 1.0) {$\vdots$};
      \visionPatchThumb{\seqx}{0.35}{2}{1}
      \node[font=\tiny, left] at (\seqx, 0.6) {$p_k$};
      \visionPatchThumb{\seqx}{-0.25}{3}{0}
      \node[font=\tiny, left] at (\seqx, 0.0) {$p_N$};

      \node[font=\scriptsize, text=gray, align=center] at (1.5, -0.6) {2D grid $\to$ patch tokens,\\continuous pixel values};
    \end{tikzpicture}
  \end{center}

  \vspace{-0.3em}
  \begin{columns}
    \begin{column}{0.48\textwidth}
      \textbf{Differences from NLP}
      \begin{itemize}
        \item \textbf{Continuous signal}: pixels are continuous values, not discrete symbols
        \item \textbf{No natural tokens}: patches are an arbitrary grid; text has words / BPE units of meaning
        \item \textbf{2D, no causal order}: spatial neighborhood replaces sequential order
        \item \textbf{High redundancy}: nearby patches are strongly correlated; word content is not
        \item \textbf{Weak inductive bias}: locality and translation equivariance are not built in
      \end{itemize}
    \end{column}
    \begin{column}{0.48\textwidth}
      \textbf{Key challenges}
      \begin{itemize}
        \item \textbf{Attention cost}: $N = HW/P^2$ tokens grows with resolution; self-attention is $\mathcal{O}(N^2)$
        \item \textbf{Sample efficiency}: weak prior, needs large pre-training data
        \item \textbf{Pre-training objective}: no next-token likelihood; need masked modeling or contrastive substitutes
        \item \textbf{Position encoding}: 2D layout must be reinjected into a sequence model
        \item \textbf{Multi-scale features}: objects span many spatial scales; hierarchical features needed
      \end{itemize}
    \end{column}
  \end{columns}

\end{frame}

% -----------------------------------------------------------------------
\begin{frame}{Vision Models: Recurring Design Themes}

  \begin{columns}[T]
    \begin{column}{0.48\textwidth}
      \begin{center}
        \begin{tikzpicture}[scale=0.65, every node/.style={scale=0.65}]
          \begin{scope}
            \clip (0,0) rectangle (2,2);
            \visionSketchScene
          \end{scope}
          \draw[thick] (0,0) rectangle (2,2);
          \foreach \i in {1,2,3} {
            \pgfmathsetmacro{\xpos}{0.5*\i}
            \draw[gray!60] (\xpos, 0) -- (\xpos, 2);
            \draw[gray!60] (0, \xpos) -- (2, \xpos);
          }
          \draw[-Stealth, thick] (2.2, 1) -- (2.9, 1);
          \visionPatchThumb{3.1}{1.6}{0}{3}
          \visionPatchThumb{3.1}{1.0}{2}{3}
          \node[font=\tiny] at (3.35, 0.7) {$\vdots$};
          \visionPatchThumb{3.1}{0.0}{3}{0}
        \end{tikzpicture}
      \end{center}
      \textbf{Patchification} \quad {\footnotesize (ViT, DeiT, MAE)}\\
      {\footnotesize Image $\to$ sequence of patch tokens.}

      \vspace{0.8em}
      \begin{center}
        \begin{tikzpicture}[scale=0.5, every node/.style={scale=0.5}]
          \draw[thick, fill=gray!8] (0,0) rectangle (5,5);
          \foreach \i in {1,2,3,4} {
            \draw[gray!40] (\i,0) -- (\i,5);
            \draw[gray!40] (0,\i) -- (5,\i);
          }
          \node[font=\scriptsize, text=gray, above] at (2.5, 5) {feature map};
          \fill[red] (2.5, 2.5) circle (5pt);
          \node[font=\scriptsize, text=red!60!black, right=2pt] at (2.5, 2.5) {$\mathbf{p}_q$};
          \draw[-Stealth, thick, blue!70!black] (2.5, 2.5) -- (0.7, 3.8);
          \fill[blue!70!black] (0.7, 3.8) circle (4pt);
          \node[font=\tiny, text=blue!70!black, left=1pt] at (0.7, 3.8) {$\Delta_1$};
          \draw[-Stealth, thick, blue!70!black] (2.5, 2.5) -- (4.3, 4.0);
          \fill[blue!70!black] (4.3, 4.0) circle (4pt);
          \node[font=\tiny, text=blue!70!black, right=1pt] at (4.3, 4.0) {$\Delta_2$};
          \draw[-Stealth, thick, blue!70!black] (2.5, 2.5) -- (4.1, 0.8);
          \fill[blue!70!black] (4.1, 0.8) circle (4pt);
          \node[font=\tiny, text=blue!70!black, right=1pt] at (4.1, 0.8) {$\Delta_3$};
          \draw[-Stealth, thick, blue!70!black] (2.5, 2.5) -- (0.9, 1.0);
          \fill[blue!70!black] (0.9, 1.0) circle (4pt);
          \node[font=\tiny, text=blue!70!black, left=1pt] at (0.9, 1.0) {$\Delta_4$};
        \end{tikzpicture}
      \end{center}
      \textbf{Restricted attention} \quad {\footnotesize (Swin, Def.\ DETR, Mask2Former)}\\
      {\footnotesize Query attends to $K$ learned sample points instead of $\mathcal{O}(HW)$ positions.}
    \end{column}
    \begin{column}{0.48\textwidth}
      \begin{center}
        \begin{tikzpicture}[scale=0.65, every node/.style={scale=0.65}]
          \begin{scope}
            \clip (0,0) rectangle (2,2);
            \visionSketchScene
          \end{scope}
          \draw[thick] (0,0) rectangle (2,2);
          \draw[very thick, green!50!black] (0.03, 0.49) rectangle (0.68, 1.58);
          \draw[very thick, red!70!black] (0.76, 0.49) rectangle (1.54, 1.52);
          \draw[very thick, violet] (1.68, 0.43) rectangle (1.9, 1.02);
          \node[font=\scriptsize, anchor=base] at (3.5, 2.1) {queries};
          \foreach \i/\y in {1/1.7, 2/1.2, 3/0.7, 4/0.2} {
            \node[draw, fill=violet!20, circle, font=\scriptsize, inner sep=1pt] (q\i) at (3.5, \y) {$q_{\i}$};
          }
          \draw[-Stealth, very thick, green!50!black] (q1.west) -- (0.4, 1.0);
          \draw[-Stealth, very thick, red!70!black]   (q2.west) -- (1.15, 1.0);
          \draw[-Stealth, very thick, gray, dashed]   (q3.west) -- ++(-0.6, 0) node[left, font=\tiny, text=gray] {$\varnothing$};
          \draw[-Stealth, very thick, violet]         (q4.west) -- (1.79, 0.7);
        \end{tikzpicture}
      \end{center}
      \textbf{Object queries} \quad {\footnotesize (DETR, MaskFormer family)}\\
      {\footnotesize Each query $\to$ (class, box) or $\varnothing$; a fixed set of slots.}

      \vspace{0.8em}
      \begin{center}
        \begin{tikzpicture}[scale=0.7, every node/.style={scale=0.7}]
          % Left: image with 75% of patches masked
          \begin{scope}
            \clip (0,0) rectangle (2,2);
            \visionSketchScene
          \end{scope}
          \foreach \r in {0,1,2,3} {
            \foreach \c in {0,1,2,3} {
              \pgfmathtruncatemacro{\visible}{(\r==2 && \c==1) ? 1 : ((\r==3 && \c==2) ? 1 : ((\r==1 && \c==3) ? 1 : ((\r==0 && \c==0) ? 1 : 0)))}
              \pgfmathsetmacro{\xx}{0.5*\c}
              \pgfmathsetmacro{\yy}{0.5*\r}
              \ifnum\visible=0
                \fill[gray!55] (\xx, \yy) rectangle (\xx+0.5, \yy+0.5);
                \node[font=\tiny, text=gray!90] at (\xx+0.25, \yy+0.25) {$\times$};
              \fi
            }
          }
          \draw[thick] (0,0) rectangle (2,2);
          \foreach \i in {1,2,3} {
            \pgfmathsetmacro{\xpos}{0.5*\i}
            \draw[gray!50] (\xpos, 0) -- (\xpos, 2);
            \draw[gray!50] (0, \xpos) -- (2, \xpos);
          }
          % Arrow
          \draw[-Stealth, thick] (2.2, 1) -- (3.1, 1);
          \node[font=\scriptsize, above] at (2.65, 1) {reconstruct};
          % Right: reconstructed scene
          \begin{scope}[shift={(3.3, 0)}]
            \clip (0,0) rectangle (2,2);
            \visionSketchScene
          \end{scope}
          \draw[thick] (3.3, 0) rectangle (5.3, 2);
          \foreach \i in {1,2,3} {
            \pgfmathsetmacro{\xpos}{3.3 + 0.5*\i}
            \pgfmathsetmacro{\ypos}{0.5*\i}
            \draw[gray!50] (\xpos, 0) -- (\xpos, 2);
            \draw[gray!50] (3.3, \ypos) -- (5.3, \ypos);
          }
        \end{tikzpicture}
      \end{center}
      \textbf{Masked image modeling} \quad {\footnotesize (MAE)}\\
      {\footnotesize Reconstruct $\sim$75\% masked patches for self-supervised pre-training.}
    \end{column}
  \end{columns}
\end{frame}

% =======================================================================
\subsection{Vision Transformer (ViT)}

% -----------------------------------------------------------------------
\begin{frame}{ViT: An Image is Worth 16$\times$16 Words}

  \cite{dosovitskiy2021vit}: Apply a \textbf{standard Transformer encoder} directly to sequences of image patches, no convolutions needed.

  \vspace{0.3em}
  \begin{center}
    \begin{tikzpicture}[
      scale=0.78, every node/.style={scale=0.78},
      box/.style={draw, rounded corners, minimum height=0.5cm, align=center, font=\small},
      arr/.style={-Stealth, thick},
    ]
      % --- Image (sketch scene) ---
      \node[inner sep=0] (img) at (-5.2, 0.44) {
        \begin{tikzpicture}
          \begin{scope}
            \clip (0,0) rectangle (2,2);
            \visionSketchScene
          \end{scope}
          \draw[thick] (0,0) rectangle (2,2);
          \foreach \i in {0.5, 1.0, 1.5} {
            \draw[gray!60] (\i, 0) -- (\i, 2);
            \draw[gray!60] (0, \i) -- (2, \i);
          }
        \end{tikzpicture}
      };
      \node[font=\scriptsize, below=0.1cm of img] {Image $H{\times}W{\times}C$};

      % --- Flatten + Linear Projection ---
      \node[box, fill=blue!10, minimum width=1.4cm] (proj) at (-1.5, 0.44) {Linear\\Proj.};
      \draw[arr] (-3.1, 0.44) -- (proj);

      % --- Patch embeddings + CLS ---
      % 7 tokens spanning y = -1.12 .. 2.0 (CLS at top)
      \foreach \i [count=\c from 0] in {0,...,6} {
        \pgfmathsetmacro{\yy}{2.0 - \c*0.52}
        \ifnum\i=0
          \node[draw, fill=violet!20, minimum width=0.85cm, minimum height=0.35cm,
                rounded corners=1pt, font=\tiny] (e\i) at (0.8, \yy) {[CLS]};
        \else
          \node[draw, fill=green!15, minimum width=0.85cm, minimum height=0.35cm,
                rounded corners=1pt, font=\tiny] (e\i) at (0.8, \yy) {$\mathbf{x}_{\i}$};
        \fi
      }
      \draw[arr] (proj) -- (e3);

      % --- Positional embeddings ---
      \node[font=\scriptsize, text=orange!70!black] at (0.8, 2.55) {$+\,E_\text{pos}$};

      % --- Transformer Encoder (sized to span the full token column) ---
      \node[box, fill=red!10, minimum width=1.6cm, minimum height=3.2cm] (enc) at (3.0, 0.44) {};
      \node[font=\small, text=red!60!black] at (3.0, 1.5) {\textbf{Transformer}};
      \node[font=\small, text=red!60!black] at (3.0, 1.1) {\textbf{Encoder}};
      \node[font=\scriptsize, text=gray] at (3.0, 0.6) {$L$ layers};
      \node[draw, rounded corners, fill=white, minimum width=1.2cm, font=\tiny, inner sep=2pt] at (3.0, 0.05) {MSA + MLP};
      \node[draw, rounded corners, fill=white, minimum width=1.2cm, font=\tiny, inner sep=2pt] at (3.0, -0.45) {MSA + MLP};

      \draw[arr] (1.3, 0.44) -- (2.1, 0.44);

      % --- MLP Head (aligned with [CLS] at top of encoder) ---
      \node[box, fill=yellow!20, minimum width=1.2cm] (head) at (5.5, 2.0) {MLP\\Head};
      \draw[arr] (3.85, 2.0) -- node[above, font=\tiny, text=violet] {[CLS] output} (head);

      % --- Class output ---
      \node[font=\small] (cls) at (7.2, 2.0) {Class};
      \draw[arr] (head) -- (cls);

      % --- Annotations ---
      \node[font=\scriptsize, text=gray, align=center] at (-5.2, -1.5) {Split into $N = \frac{HW}{P^2}$\\patches of size $P{\times}P$};
      \node[font=\scriptsize, text=gray, align=center] at (0.8, -1.9) {Patch embeddings\\$\in \mathbb{R}^{(N+1) \times D}$};
    \end{tikzpicture}
  \end{center}

  \vspace{0.2em}
  \begin{itemize}
    \item Split image into $N = HW/P^2$ fixed-size patches ($P{=}16$), flatten and linearly project to $D$ dimensions (implemented as \texttt{Conv2d} with kernel and stride $=P$).
    \item Prepend a learnable \textcolor{violet}{\textbf{[CLS]}} token whose final state is used for classification.
    \item \textbf{Positional encoding}: add \textbf{learnable 1D absolute} embeddings $E_\text{pos} \in \mathbb{R}^{(N+1) \times D}$ (one trainable vector per patch slot $+$ [CLS]), init.\ trunc.\ normal std $0.02$; added once at the input. Attention is otherwise permutation-invariant.
    \[
      \mathbf{z}_0 = [\mathbf{x}_\text{class};\; \mathbf{x}_p^1 E;\; \ldots;\; \mathbf{x}_p^N E] + E_\text{pos}
    \]
    \item Standard Transformer encoder; classify from [CLS] output.
    \item Requires large-scale pre-training (JFT-300M) to outperform CNNs.
  \end{itemize}
\end{frame}

% -----------------------------------------------------------------------
\begin{frame}[fragile]{ViT: Transformer Block}

  \begin{columns}
    \begin{column}{0.48\textwidth}
      \textbf{Transformer encoder block} (pre-norm):
      \begin{align*}
        \mathbf{z}'_\ell &= \operatorname{MSA}(\operatorname{LN}(\mathbf{z}_{\ell-1})) + \mathbf{z}_{\ell-1} \\
        \mathbf{z}_\ell &= \operatorname{MLP}(\operatorname{LN}(\mathbf{z}'_\ell)) + \mathbf{z}'_\ell
      \end{align*}
      \begin{itemize}
        \item $\operatorname{LN}$ before each sublayer (pre-norm), residual connections around both.
        \item $\operatorname{MSA}$: multi-head self-attention over all patch tokens.
        \item $\operatorname{MLP}$: two linear layers with a GELU nonlinearity.
        \item Classification head: $\mathbf{y} = \operatorname{LN}(\mathbf{z}_L^0)$ (the [CLS] token of the last layer).
      \end{itemize}
    \end{column}
    \begin{column}{0.48\textwidth}
      \begin{center}
        \begin{tikzpicture}[
          scale=0.85, every node/.style={scale=0.85},
          box/.style={draw, rounded corners, fill=#1, minimum width=2.2cm,
                      minimum height=0.45cm, align=center, font=\small},
          arr/.style={-Stealth, thick},
        ]
          % Input
          \node[font=\small] (in) at (0, -0.3) {$\mathbf{z}_{\ell-1}$};

          % LN1
          \node[box=gray!10] (ln1) at (0, 0.6) {Layer Norm};
          \draw[arr] (in) -- (ln1);

          % MSA
          \node[box=red!15] (msa) at (0, 1.6) {Multi-Head SA};
          \draw[arr] (ln1) -- (msa);

          \node[draw, circle, inner sep=1pt, font=\scriptsize, fill=white] (add1) at (0, 2.2) {$+$};
          \draw[arr] (msa) -- (add1);

          % LN2
          \node[box=gray!10] (ln2) at (0, 3.0) {Layer Norm};
          \draw[arr] (add1) -- (ln2);

          % MLP
          \node[box=blue!12] (mlp) at (0, 4.0) {MLP};
          \draw[arr] (ln2) -- (mlp);

          \node[draw, circle, inner sep=1pt, font=\scriptsize, fill=white] (add2) at (0, 4.6) {$+$};
          \draw[arr] (mlp) -- (add2);

          % Residual skip connections (both at the same x, pushed to the right of the boxes)
          \draw[arr, gray!50] (in.east) -- (1.8, -0.3) -- (1.8, 2.2) -- (add1.east);
          \draw[arr, gray!50] (add1.east) -- (1.8, 2.2) -- (1.8, 4.6) -- (add2.east);

          % Output
          \node[font=\small] (out) at (0, 5.3) {$\mathbf{z}_\ell$};
          \draw[arr] (add2) -- (out);
        \end{tikzpicture}
      \end{center}
    \end{column}
  \end{columns}

  \vspace{0.3em}
  \begin{center}
    \small
    \begin{tabular}{lccccc}
      \toprule
      \textbf{Variant} & \textbf{Layers} & \textbf{Hidden $D$} & \textbf{MLP dim} & \textbf{Heads} & \textbf{Params} \\
      \midrule
      ViT-B/16 & 12 & 768 & 3072 & 12 & 86M \\
      ViT-L/16 & 24 & 1024 & 4096 & 16 & 307M \\
      ViT-H/14 & 32 & 1280 & 5120 & 16 & 632M \\
      \bottomrule
    \end{tabular}
  \end{center}
\end{frame}

% -----------------------------------------------------------------------
\begin{frame}{ViT: Training Recipe --- Data Augmentation}

  ViT has \textbf{little image-specific inductive bias} (no convolutions, no locality), so it overfits small data. Augmentation and regularization substitute for missing bias and can match a $10\times$ larger pre-training set \cite{steiner2022augreg}.

  \vspace{0.4em}
  \begin{description}
    \item[Inception-style crop] random resized crop: sample a scale (e.g.\ $8\%$--$100\%$ of area) and aspect ratio, then resize to $224^2$.
    \item[Horizontal flip] independently mirror left/right with probability $0.5$.
    \item[RandAugment] apply $N$ randomly chosen operations (rotate, shear, color/contrast, solarize, posterize, \dots) at magnitude $M$; typical $N{=}2$, $M{\in}[9,15]$.
    \item[Mixup] convex combination of two images and their (soft) labels, $\tilde{x}=\lambda x_i + (1{-}\lambda) x_j$, $\lambda \sim \mathrm{Beta}(\alpha,\alpha)$.
    \item[CutMix] paste a rectangular patch of one image into another; mix labels by patch area.
  \end{description}

  \vspace{0.3em}
  \begin{itemize}
    \item Original ViT used \emph{minimal} augmentation but relied on JFT-300M; on smaller data (ImageNet-1k/21k) strong AugReg is essential.
    \item Augmentation strength should scale \emph{down} as the dataset grows: too much aug on huge data hurts.
  \end{itemize}
\end{frame}

% -----------------------------------------------------------------------
\begin{frame}{Augmentation: Inception-style Crop \& Horizontal Flip}
  \begin{center}
    \resizebox{0.92\textwidth}{!}{%
    \begin{tikzpicture}[
        arr/.style={-Stealth, thick},
        sub/.style={font=\footnotesize, text=gray},
      ]
      % --- crop ---
      \begin{scope}\clip (0,0) rectangle (2,2);\visionSketchScene\end{scope}
      \draw[thick] (0,0) rectangle (2,2);
      \draw[red, very thick, dashed] (0.6,0.3) rectangle (1.7,1.35);
      \node[sub] at (1,-0.4) {sample box};
      \draw[arr] (2.2,1) -- (2.8,1);
      \begin{scope}[shift={(3.0,0)}]
        \begin{scope}
          \clip (0,0) rectangle (2,2);
          \begin{scope}[xscale={2/1.1}, yscale={2/1.05}, shift={(-0.6,-0.3)}]
            \visionSketchScene
          \end{scope}
        \end{scope}
        \draw[red, very thick] (0,0) rectangle (2,2);
        \node[sub] at (1,-0.4) {resize $\to 224^2$};
      \end{scope}
      % --- flip ---
      \begin{scope}[shift={(7.0,0)}]
        \begin{scope}\clip (0,0) rectangle (2,2);\visionSketchScene\end{scope}
        \draw[thick] (0,0) rectangle (2,2);
        \node[sub] at (1,-0.4) {original};
        \draw[arr] (2.2,1) -- (2.8,1);
        \node[font=\footnotesize] at (2.5,1.35) {$p{=}0.5$};
        \begin{scope}[shift={(3.0,0)}]
          \begin{scope}
            \clip (0,0) rectangle (2,2);
            \begin{scope}[shift={(2,0)}]\begin{scope}[xscale=-1]\visionSketchScene\end{scope}\end{scope}
          \end{scope}
          \draw[thick] (0,0) rectangle (2,2);
          \node[sub] at (1,-0.4) {mirrored};
        \end{scope}
      \end{scope}
    \end{tikzpicture}}
  \end{center}
  \vspace{0.6em}
  \begin{itemize}
    \item \textbf{Random resized crop}: sample an area scale ($8\%$--$100\%$) and aspect ratio ($\tfrac34$--$\tfrac43$), crop that region, resize to $224^2$. Teaches scale/translation invariance.
    \item \textbf{Horizontal flip}: independent left/right mirror with $p{=}0.5$. Cheap; assumes horizontal symmetry, so skip it for orientation-sensitive data (text, many medical/satellite images).
  \end{itemize}
\end{frame}

% -----------------------------------------------------------------------
\begin{frame}{Augmentation: RandAugment}
  \begin{center}
    \resizebox{0.92\textwidth}{!}{%
    \begin{tikzpicture}[
        arr/.style={-Stealth, thick},
        sub/.style={font=\footnotesize, text=gray},
      ]
      \sceneAframed{}
      \node[sub] at (1,-0.4) {original};
      \draw[arr] (2.2,1) -- (2.8,1);
      \node[font=\footnotesize] at (2.5,1.35) {op 1};
      \begin{scope}[shift={(3.0,0)}]
        \begin{scope}
          \clip (0,0) rectangle (2,2);
          \fill[gray!25] (0,0) rectangle (2,2);
          \begin{scope}[rotate around={14:(1,1)}]\visionSketchScene\end{scope}
        \end{scope}
        \draw[thick] (0,0) rectangle (2,2);
        \node[sub] at (1,-0.4) {rotate};
      \end{scope}
      \draw[arr] (5.2,1) -- (5.8,1);
      \node[font=\footnotesize] at (5.5,1.35) {op 2};
      \begin{scope}[shift={(6.0,0)}]
        \begin{scope}
          \clip (0,0) rectangle (2,2);
          \fill[gray!25] (0,0) rectangle (2,2);
          \begin{scope}[cm={1,0,0.25,1,(-0.25,0)}]\visionSketchScene\end{scope}
          \fill[magenta!30, opacity=0.4] (0,0) rectangle (2,2);
        \end{scope}
        \draw[thick] (0,0) rectangle (2,2);
        \node[sub] at (1,-0.4) {shear + color};
      \end{scope}
    \end{tikzpicture}}
  \end{center}
  \vspace{0.6em}
  \begin{itemize}
    \item Two integer hyperparameters control everything:
      \begin{description}
        \item[$N$] \textbf{how many} operations are applied to each image: sample $N$ ops \emph{uniformly at random} from a fixed pool of $\sim$14 (rotate, shear, translate, contrast, brightness, color, solarize, posterize, equalize, sharpness, \dots) and apply them \emph{in sequence}.
        \item[$M$] \textbf{how strong} each operation is: a single magnitude \emph{shared} by all $N$ chosen ops.
      \end{description}
    \item Replaces learned AutoAugment policies with just these \textbf{two hyperparameters} $(N, M)$, tunable by a simple grid search; typical $N{=}2$, $M{\in}[9,15]$.
    \item Geometric ops leave empty borders, filled with a constant (gray) or reflection.
  \end{itemize}
\end{frame}

% -----------------------------------------------------------------------
\begin{frame}{Augmentation: Mixup}
  \begin{center}
    \resizebox{0.92\textwidth}{!}{%
    \begin{tikzpicture}[
        arr/.style={-Stealth, thick},
        sub/.style={font=\footnotesize, text=gray},
      ]
      \sceneAframed{}
      \node[sub] at (1,-0.4) {house};
      \node[font=\Large] at (2.5,1) {$+$};
      \begin{scope}[shift={(3.0,0)}]
        \begin{scope}\clip (0,0) rectangle (2,2);\visionSketchSceneB\end{scope}
        \draw[thick] (0,0) rectangle (2,2);
        \node[sub] at (1,-0.4) {mountain};
      \end{scope}
      \draw[arr] (5.2,1) -- (5.8,1);
      \begin{scope}[shift={(6.0,0)}]
        \begin{scope}
          \clip (0,0) rectangle (2,2);
          \visionSketchScene
          \begin{scope}[opacity=0.55]\visionSketchSceneB\end{scope}
        \end{scope}
        \draw[thick] (0,0) rectangle (2,2);
        \node[sub] at (1,-0.4) {$0.6\,$house $+\,0.4\,$mountain};
      \end{scope}
    \end{tikzpicture}}
  \end{center}
  \vspace{0.6em}
  \begin{itemize}
    \item \textbf{Pixel-wise} convex blend of two samples: $\tilde{x}=\lambda x_i+(1{-}\lambda)x_j$, with $\lambda\sim\mathrm{Beta}(\alpha,\alpha)$.
    \item Labels blend with the \emph{same} $\lambda$: $\tilde{y}=\lambda y_i+(1{-}\lambda)y_j$.
    \item \textbf{Loss}: since cross-entropy is linear in the target, the blended-label loss splits into a $\lambda$-weighted sum of two ordinary CE terms:
      \[
        \mathcal{L}_\text{mixup} = \lambda\,\mathrm{CE}\big(f(\tilde{x}),\, y_i\big) + (1{-}\lambda)\,\mathrm{CE}\big(f(\tilde{x}),\, y_j\big).
      \]
    \item Encourages linear behavior between classes; smooths decision boundaries and improves robustness. Produces a ghosted, semi-transparent overlay.
  \end{itemize}
\end{frame}

% -----------------------------------------------------------------------
\begin{frame}{Augmentation: CutMix}
  \begin{center}
    \resizebox{0.92\textwidth}{!}{%
    \begin{tikzpicture}[
        arr/.style={-Stealth, thick},
        sub/.style={font=\footnotesize, text=gray},
      ]
      \sceneAframed{}
      \node[sub] at (1,-0.4) {house};
      \node[font=\Large] at (2.5,1) {$+$};
      \begin{scope}[shift={(3.0,0)}]
        \begin{scope}\clip (0,0) rectangle (2,2);\visionSketchSceneB\end{scope}
        \draw[thick] (0,0) rectangle (2,2);
        \node[sub] at (1,-0.4) {mountain};
      \end{scope}
      \draw[arr] (5.2,1) -- (5.8,1);
      \begin{scope}[shift={(6.0,0)}]
        \begin{scope}\clip (0,0) rectangle (2,2);\visionSketchScene\end{scope}
        \begin{scope}
          \clip (1.0,0.15) rectangle (1.85,1.1);
          \visionSketchSceneB
        \end{scope}
        \draw[red, very thick] (1.0,0.15) rectangle (1.85,1.1);
        \draw[thick] (0,0) rectangle (2,2);
        \node[sub] at (1,-0.4) {$0.8\,$house $+\,0.2\,$mountain};
      \end{scope}
    \end{tikzpicture}}
  \end{center}
  \vspace{0.4em}
  \begin{itemize}
    \item Cut a rectangular patch from image $B$ and paste it onto $A$: with a binary mask $M$, \;$\tilde{x}=M\odot x_A+(1{-}M)\odot x_B$. Patches stay \textbf{sharp} (no blending).
    \item \textbf{Label by area}: let $\lambda=\dfrac{\text{pixels kept from }A}{H\cdot W}$. The composite is literally $\lambda$ class-$A$ pixels and $(1{-}\lambda)$ class-$B$ pixels, so the soft label mirrors the visible proportion: $\tilde{y}=\lambda y_A+(1{-}\lambda)y_B$.
    \item Box sized as $r_w{=}W\sqrt{1{-}\lambda},\; r_h{=}H\sqrt{1{-}\lambda}$ (area $=(1{-}\lambda)HW$), then $\lambda$ recomputed from the actual box area. Here the mountain patch covers $\sim20\%$ $\Rightarrow$ weight $0.2$.
  \end{itemize}
\end{frame}

% -----------------------------------------------------------------------
\begin{frame}{Augmentation: Repeated Augmentation (Batching)}

  \cite{hoffer2020augment}: build a batch from only $B/m$ \emph{distinct} images, each entering $m$ times with different augmentations (DeiT uses $m{=}3$). All copies share one label.

  \vspace{0.1em}
  \begin{center}
    \resizebox{0.61\textwidth}{!}{%
    \begin{tikzpicture}[
        arr/.style={-Stealth, thick},
        sub/.style={font=\footnotesize, text=gray},
      ]
      % source images
      \begin{scope}\clip (0,2.2) rectangle (1.2,3.4);
        \begin{scope}[shift={(0,2.2)},scale=0.6]\visionSketchScene\end{scope}\end{scope}
      \draw[thick] (0,2.2) rectangle (1.2,3.4);
      \node[sub] at (0.6,2.0) {image $i$ (house)};
      \begin{scope}\clip (0,0.0) rectangle (1.2,1.2);
        \begin{scope}[shift={(0,0.0)},scale=0.6]\visionSketchSceneB\end{scope}\end{scope}
      \draw[thick] (0,0.0) rectangle (1.2,1.2);
      \node[sub] at (0.6,-0.2) {image $j$ (mountain)};

      % batch container
      \draw[rounded corners, draw=gray!60, very thick] (3.0,-0.8) rectangle (7.4,4.2);
      \node[sub] at (5.2,4.45) {mini-batch (size $B$)};

      % augmented copies of A (top row): arrows curve up and over the thumbnails
      \rathumb{3.4}{2.5}{rotate around={12:(1,1)}}{\visionSketchScene}
      \rathumb{4.7}{2.5}{shift={(2,0)},xscale=-1}{\visionSketchScene}
      \rathumb{6.0}{2.5}{shift={(-0.5,-0.4)},scale=1.3}{\visionSketchScene}
      \foreach \x in {3.4,4.7,6.0} {
        \draw[arr, gray!70] (1.25,3.0) to[out=34,in=146] (\x+0.5,3.55);
        \node[sub] at (\x+0.5,2.22) {house};
      }
      % augmented copies of B (bottom row): arrows curve down and under the thumbnails
      \rathumb{3.4}{0.1}{rotate around={-10:(1,1)}}{\visionSketchSceneB}
      \rathumb{4.7}{0.1}{shift={(2,0)},xscale=-1}{\visionSketchSceneB}
      \rathumb{6.0}{0.1}{shift={(-0.45,-0.3)},scale=1.25}{\visionSketchSceneB}
      \foreach \x in {3.4,4.7,6.0} {
        \draw[arr, gray!70] (1.25,0.4) to[out=-34,in=-146] (\x+0.5,0.05);
        \node[sub] at (\x+0.5,1.3) {mountain};
      }
      \node[sub] at (5.2,1.75) {$m{=}3$ views each};
    \end{tikzpicture}}
  \end{center}

  \vspace{0.1em}
  \begin{itemize}
    \item The loss averages several augmented \emph{views of the same image} per step, lowering augmentation gradient variance (batch is no longer i.i.d.).
    \item DeiT finds RA \textbf{essential} on ImageNet-1k; little benefit once data is abundant.
  \end{itemize}
\end{frame}

% -----------------------------------------------------------------------
\begin{frame}{Regularization: Stochastic Depth}
  \begin{columns}[c]
    \begin{column}{0.40\textwidth}
      \begin{center}
        \begin{tikzpicture}[
            scale=0.9, every node/.style={scale=0.9},
            arr/.style={-Stealth, thick},
            sub/.style={font=\scriptsize, text=gray},
          ]
          \node[font=\scriptsize] (in) at (0,-0.4) {$\mathbf{z}_{\ell-1}$};
          % --- Block 1 (kept) ---
          \node[draw, rounded corners, fill=blue!12, minimum width=2.0cm, minimum height=0.6cm,
                font=\scriptsize] (b1) at (0,0.7) {Block};
          \node[draw, circle, inner sep=1pt, font=\scriptsize, fill=white] (a1) at (0,1.4) {$+$};
          \draw[arr] (in) -- (b1);
          \draw[arr] (b1) -- (a1);
          \draw[arr, gray!55] (0,-0.1) -- (1.35,-0.1) -- (1.35,1.4) -- (a1.east);
          \node[sub] at (2.5,0.7) {$p_\ell{=}0.95$};
          % --- Block 2 (dropped) ---
          \node[draw, rounded corners, fill=gray!12, dashed, minimum width=2.0cm, minimum height=0.6cm,
                font=\scriptsize, text=gray!70] (b2) at (0,2.1) {Block};
          \draw[red, thick] (b2.south west) -- (b2.north east);
          \draw[red, thick] (b2.north west) -- (b2.south east);
          \node[sub, text=red] at (-1.55,2.1) {dropped};
          \node[draw, circle, inner sep=1pt, font=\scriptsize, fill=white] (a2) at (0,2.8) {$+$};
          \draw[arr, gray!40, dashed] (a1) -- (b2);
          \draw[arr, gray!40, dashed] (b2) -- (a2);
          \draw[arr, green!55!black, very thick] (0,1.4) -- (1.35,1.4) -- (1.35,2.8) -- (a2.east);
          \node[sub] at (2.5,2.1) {$p_\ell{=}0.78$};
          % --- Block 3 (kept) ---
          \node[draw, rounded corners, fill=blue!12, minimum width=2.0cm, minimum height=0.6cm,
                font=\scriptsize] (b3) at (0,3.5) {Block};
          \node[draw, circle, inner sep=1pt, font=\scriptsize, fill=white] (a3) at (0,4.2) {$+$};
          \draw[arr] (a2) -- (b3);
          \draw[arr] (b3) -- (a3);
          \draw[arr, gray!55] (0,2.8) -- (1.35,2.8) -- (1.35,4.2) -- (a3.east);
          \node[sub] at (2.5,3.5) {$p_\ell{=}0.60$};
          \node[font=\scriptsize] (out) at (0,4.9) {$\mathbf{z}_L$};
          \draw[arr] (a3) -- (out);
        \end{tikzpicture}
      \end{center}
    \end{column}
    \begin{column}{0.60\textwidth}
      \begin{itemize}
        \item During training, drop each residual block independently with prob.\ $1-p_\ell$; a dropped block becomes the \textbf{identity} (only the skip passes, green).
        \item Survival prob.\ decays \emph{linearly} with depth, $p_\ell = 1 - \frac{\ell}{L}(1-p_L)$, so later blocks are dropped more often.
        \item Acts as an implicit \textbf{ensemble} over sub-networks of varying depth; lets very deep ViT-L/H train stably and regularizes strongly.
        \item At test time \textbf{all} blocks are kept; the residual branch is scaled by $p_\ell$ to match the training expectation.
      \end{itemize}
    \end{column}
  \end{columns}
\end{frame}

% -----------------------------------------------------------------------
\begin{frame}{Regularization: Label Smoothing}
  \begin{center}
    \begin{tikzpicture}[
        scale=0.95, every node/.style={scale=0.95},
        arr/.style={-Stealth, thick},
        sub/.style={font=\scriptsize, text=gray},
      ]
      % --- left: one-hot ---
      \begin{scope}
        \def\bh{3.0}
        \draw[->] (-0.3,0) -- (2.6,0);
        \draw[->] (-0.3,0) -- (-0.3,3.5);
        \node[sub, rotate=90] at (-0.65,1.6) {prob.};
        \foreach \i/\h in {0/0, 1/0, 2/1, 3/0} {
          \pgfmathsetmacro{\x}{0.05+\i*0.62}
          \pgfmathsetmacro{\bb}{\h*\bh}
          \ifnum\h=1
            \fill[red!55] (\x,0) rectangle (\x+0.42,\bb);
            \node[sub] at (\x+0.21,\bb+0.2) {1.0};
          \else
            \draw[gray!50] (\x,0) -- (\x+0.42,0);
            \node[sub] at (\x+0.21,0.18) {0};
          \fi
          \node[sub] at (\x+0.21,-0.28) {$c_\i$};
        }
        \node[sub] at (1.1,-0.7) {one-hot $y$};
      \end{scope}
      % arrow
      \draw[arr] (2.85,1.5) -- (3.55,1.5);
      % --- right: smoothed ---
      \begin{scope}[shift={(3.95,0)}]
        \def\bh{3.0}
        \draw[->] (-0.3,0) -- (2.6,0);
        \draw[->] (-0.3,0) -- (-0.3,3.5);
        \foreach \i/\h/\v in {0/0.033/0.03, 1/0.033/0.03, 2/0.9/0.9, 3/0.033/0.03} {
          \pgfmathsetmacro{\x}{0.05+\i*0.62}
          \pgfmathsetmacro{\bb}{\h*\bh}
          \ifnum\i=2
            \fill[red!55] (\x,0) rectangle (\x+0.42,\bb);
          \else
            \fill[blue!40] (\x,0) rectangle (\x+0.42,\bb);
          \fi
          \node[sub] at (\x+0.21,\bb+0.2) {\v};
          \node[sub] at (\x+0.21,-0.28) {$c_\i$};
        }
        \node[sub] at (1.1,-0.7) {smoothed target};
      \end{scope}
    \end{tikzpicture}
  \end{center}
  \vspace{0.6em}
  \begin{itemize}
    \item Replace the one-hot target by $\tilde{y} = (1{-}\alpha)\,y + \alpha/K$ (here $\alpha{=}0.1$, $K$ classes): the true class gets $1{-}\alpha+\alpha/K$, every other class $\alpha/K$.
    \item Stops the model from driving the correct logit to $+\infty$: penalizes over-confidence, improves \textbf{calibration} and generalization.
    \item Standard for ViT/DeiT classification ($\alpha{=}0.1$); composes naturally with the soft labels from Mixup/CutMix.
  \end{itemize}
\end{frame}

% -----------------------------------------------------------------------
\begin{frame}{ViT: Weight Initialization}
  \begin{columns}[c]
    \begin{column}{0.56\textwidth}
      \begin{itemize}
        \item \textbf{Linear layers \& patch-embed projection}: truncated normal, std $=0.02$, truncated at $\pm 2\sigma$; biases $=0$.
        \item \textbf{[CLS] token} and \textbf{positional embeddings} $E_\text{pos}$: truncated normal, std $=0.02$.
        \item \textbf{LayerNorm}: $\gamma=1$, $\beta=0$.
        \item \textbf{Classification head}: zero-initialized $\Rightarrow$ uniform initial predictions (pairs well with label smoothing); re-zeroed when a new head is attached at fine-tuning.
      \end{itemize}
      \vspace{0.3em}
      \footnotesize
      The patch embed is a \texttt{Conv2d} (kernel $=$ stride $=P$) but initialized like a \emph{linear} layer, not a vision conv.
    \end{column}
    \begin{column}{0.44\textwidth}
      \begin{center}
        \begin{tikzpicture}[scale=0.95]
          \draw[->] (-3.2,0) -- (3.3,0) node[right, font=\scriptsize] {$w$};
          \draw[->] (0,-0.05) -- (0,2.5);
          % tails (removed)
          \draw[gray!55, dashed, domain=-3:-2, smooth, samples=40] plot (\x,{2*exp(-\x*\x/2)});
          \draw[gray!55, dashed, domain=2:3, smooth, samples=40] plot (\x,{2*exp(-\x*\x/2)});
          % kept region
          \fill[blue!15] (-2,0) -- plot[domain=-2:2, smooth, samples=60] (\x,{2*exp(-\x*\x/2)}) -- (2,0) -- cycle;
          \draw[blue!60!black, thick, domain=-2:2, smooth, samples=60] plot (\x,{2*exp(-\x*\x/2)});
          % cutoffs
          \draw[red, dashed] (-2,0) -- (-2,{2*exp(-2)});
          \draw[red, dashed] (2,0) -- (2,{2*exp(-2)});
          \node[font=\scriptsize] at (-2,-0.28) {$-2\sigma$};
          \node[font=\scriptsize] at (2,-0.28) {$+2\sigma$};
          \node[red, font=\scriptsize] at (2.55,0.85) {cut};
          \node[red, font=\scriptsize] at (-2.55,0.85) {cut};
          \node[font=\scriptsize, text=blue!60!black] at (0,2.25) {std $=0.02$};
        \end{tikzpicture}
      \end{center}
      \footnotesize
      Small std keeps the residual stream bounded across many summed blocks; truncation drops rare large draws ($>2\sigma$) that destabilize early training.
    \end{column}
  \end{columns}
\end{frame}

% -----------------------------------------------------------------------
\begin{frame}{ViT: Training Recipe --- Regularization, Optimization \& Fine-tuning}

  \begin{columns}[T]
    \begin{column}{0.5\textwidth}
      \textbf{Regularization}
      \begin{itemize}
        \item \textbf{Weight decay}: large, $0.1$ (decoupled / AdamW); dominant regularizer for pre-training.
        \item \textbf{Dropout} $0.1$ on attention/MLP/embeddings (pre-training only).
        \item \textbf{Stochastic depth}: randomly drop residual blocks, rate grows with depth (key for ViT-L/H).
        \item \textbf{Label smoothing} $0.1$.
      \end{itemize}

      \vspace{0.4em}
      \textbf{Optimizer}
      \begin{itemize}
        \item \textbf{AdamW}, $\beta_1{=}0.9$, $\beta_2{=}0.999$.
        \item Large \textbf{batch size} $4096$.
        \item \textbf{Gradient clipping} at global norm $1.0$ (stabilizes early training).
      \end{itemize}
    \end{column}
    \begin{column}{0.5\textwidth}
      \textbf{Learning-rate schedule}
      \begin{itemize}
        \item Linear \textbf{warmup} ($\sim$10k steps), then cosine (or linear) decay.
        \item Peak LR $\sim 10^{-3}$ for pre-training.
      \end{itemize}

      \vspace{0.4em}
      \textbf{Pre-train $\rightarrow$ fine-tune}
      \begin{itemize}
        \item Pre-train at $224^2$; fine-tune at higher resolution ($384^2$) for accuracy.
        \item Higher resolution $\Rightarrow$ more patches: \textbf{2D-interpolate} positional embeddings to the new grid.
        \item Fine-tune with \textbf{SGD + momentum} $0.9$, smaller batch, short cosine schedule, no weight decay.
        \item Replace pre-training MLP head with a single zero-initialized linear layer.
      \end{itemize}
    \end{column}
  \end{columns}

  \vspace{0.4em}
  {\footnotesize Use \textbf{EMA} of weights and average over fine-tuning checkpoints for a small extra boost.}
\end{frame}

% =======================================================================
\subsection{DeiT: Data-efficient Image Transformers}

% -----------------------------------------------------------------------
\begin{frame}[fragile]{DeiT: Knowledge Distillation for ViT}

  \cite{touvron2021deit}: Train ViT \textbf{on ImageNet-1K only} (no JFT-300M) via a \textbf{distillation token} and strong data augmentation.

  \vspace{0.1em}
  \begin{center}
    \resizebox{0.80\textwidth}{!}{%
    \begin{tikzpicture}[
        box/.style={draw, rounded corners, minimum height=0.55cm, align=center, font=\small},
        tok/.style={draw, rounded corners=1pt, minimum width=0.75cm, minimum height=0.36cm, font=\tiny},
        loss/.style={draw, rounded corners, fill=yellow!25, minimum width=1.2cm, minimum height=0.5cm, font=\small},
        arr/.style={-Stealth, thick},
      ]
      % ---------- Encoder ----------
      \node[box, fill=red!12, minimum width=4.6cm, minimum height=1.5cm] (enc) at (0,2.0) {};
      \node[font=\small, text=red!60!black] at (0,2.25) {\textbf{Transformer Encoder}};
      \node[font=\scriptsize, text=gray] at (0,1.78) {(ViT student, $L$ layers)};

      % ---------- Input tokens ----------
      \node[tok, fill=violet!20] (cls)  at (-1.8,0.4) {[CLS]};
      \node[tok, fill=teal!20]   (dist) at (-0.9,0.4) {[DIST]};
      \node[tok, fill=green!15]  (p1)   at (0.0,0.4) {$\mathbf{x}_1$};
      \node[font=\small] at (0.75,0.4) {$\cdots$};
      \node[tok, fill=green!15]  (pn)   at (1.6,0.4) {$\mathbf{x}_N$};

      % merge bus into encoder
      \foreach \t in {cls,dist,p1,pn} { \draw[gray!60] (\t.north) -- (\t.north |- 0,0.92); }
      \draw[gray!60] (-1.8,0.92) -- (1.6,0.92);
      \draw[arr] (-0.1,0.92) -- (-0.1,1.25);
      \node[font=\scriptsize, text=gray, anchor=west] at (1.75,0.7) {patch + pos. embed};

      % ---------- Heads ----------
      \node[box, fill=violet!15] (clshead) at (-1.7,3.9) {Class\\head};
      \node[box, fill=teal!15]   (disthead) at (1.7,3.9) {Distill\\head};
      \draw[arr, violet]        (-1.7,2.75) -- (clshead.south);
      \draw[arr, teal!70!black] (1.7,2.75)  -- (disthead.south);
      \node[font=\scriptsize, text=violet, anchor=west] at (-1.6,3.25) {[CLS] out};
      \node[font=\scriptsize, text=teal!70!black, anchor=east] at (1.6,3.25) {[DIST] out};

      % ---------- Losses ----------
      \node[loss] (lce)  at (-1.7,5.4) {$\mathcal{L}_\text{CE}$};
      \node[loss] (ldist) at (1.7,5.4) {$\mathcal{L}_\text{distill}$};
      \draw[arr] (clshead.north)  -- (lce.south);
      \draw[arr] (disthead.north) -- (ldist.south);

      % ground truth
      \node[box, fill=gray!12] (gt) at (-4.4,5.4) {ground\\truth $y$};
      \draw[arr] (gt.east) -- (lce.west);

      % teacher
      \node[box, fill=orange!20, minimum width=2.0cm] (teacher) at (4.5,3.9) {Teacher CNN};
      \node[font=\scriptsize, text=gray] at (4.5,3.4) {(RegNetY)};
      \draw[arr, orange!80!black] (teacher.north) -- (4.5,5.4) -- (ldist.east);
      \node[font=\scriptsize, text=orange!70!black, anchor=south west] at (4.55,4.6) {$y_t=\arg\max Z_t$};
    \end{tikzpicture}}
  \end{center}

  \vspace{0.2em}
  \begin{itemize}
    \item Input: $[\text{CLS}];\; [\text{DIST}];\; \mathbf{x}_1;\; \ldots;\; \mathbf{x}_N$ \quad ($N{+}2$ tokens)
    \item \textbf{Positional encoding}: same \textbf{learnable 1D absolute} $E_\text{pos}$ as ViT, with one extra trainable entry for the [DIST] token; the [DIST] token itself is learned and position-agnostic.
    \item \textcolor{violet}{\textbf{[CLS]}} head: cross-entropy vs.\ ground truth labels
    \item \textcolor{teal}{\textbf{[DIST]}} head: cross-entropy vs.\ teacher's hard labels (\textbf{hard distillation} works best)
    \item At inference: average both heads' predictions
  \end{itemize}
\end{frame}

% -----------------------------------------------------------------------
\begin{frame}{DeiT: Hard vs Soft Distillation}

  \begin{columns}
    \begin{column}{0.48\textwidth}
      \textbf{Soft distillation} (KL divergence):
      \begin{align*}
        \mathcal{L}_\text{soft} = {}& (1{-}\lambda)\,\mathcal{L}_\text{CE}(\psi(Z_s), y) \\
        &+ \lambda\,\tau^2\, \text{KL}(\psi(Z_s/\tau) \,\|\, \psi(Z_t/\tau))
      \end{align*}
      \begin{itemize}
        \item $Z_s, Z_t$: student/teacher logits
        \item $\tau$: temperature, $\psi$: softmax
      \end{itemize}

      \vspace{0.5em}
      \textbf{Hard distillation} (best results):
      \[
        \mathcal{L}_\text{hard} = \tfrac{1}{2}\,\mathcal{L}_\text{CE}(\psi(Z_s), y) + \tfrac{1}{2}\,\mathcal{L}_\text{CE}(\psi(Z_s), y_t)
      \]
      where $y_t = \arg\max Z_t$ (teacher's hard prediction)

      \vspace{0.4em}
      \textbf{Why hard distillation?}
      {\footnotesize
      \begin{itemize}\itemsep1pt
        \item Teacher \textbf{relabels each augmented view}: after crop/Mixup/CutMix the GT label may not match the pixels, but $y_t$ does.
        \item Targets are \emph{already} soft (label smoothing, Mixup) $\Rightarrow$ the teacher's soft distribution adds little.
        \item Parameter-free (no $\tau$, no $\lambda$) and more robust.
      \end{itemize}}
    \end{column}
    \begin{column}{0.48\textwidth}
      \textbf{Training strategy} (data-efficient):
      {\footnotesize
      \begin{itemize}\itemsep1pt
        \item ImageNet-1K only (1.2M images); AdamW + cosine schedule, $<3$ days on 8 GPUs
        \item Strong augmentation + regularization (see next slide)
      \end{itemize}}

      \vspace{0.3em}
      \textbf{ImageNet-1K top-1} (trained on IN-1K only):
      \begin{center}
        \small
        \begin{tabular}{lccc}
          \toprule
          \textbf{Variant} & \textbf{Params} & \textbf{Base} & \textbf{$+$Distill} \\
          \midrule
          DeiT-Ti & 5M  & 72.2 & 74.5 \\
          DeiT-S  & 22M & 79.8 & 81.2 \\
          DeiT-B  & 86M & 81.8 & 83.4 \\
          DeiT-B\,$\uparrow$384 & 86M & 83.1 & \textbf{85.2} \\
          \bottomrule
        \end{tabular}
      \end{center}
      {\footnotesize Baseline \textbf{ViT-B/16} on IN-1K only: \textbf{77.9} (vs.\ 84.2 with JFT-300M pre-training). DeiT-B already beats it; distillation adds $+1.6$. \par}

      \vspace{0.3em}
      \textbf{The teacher: RegNetY} \cite{radosavovic2020regnet}
      {\footnotesize
      \begin{itemize}\itemsep1pt
        \item Efficient ConvNet from \emph{designing a parameterized design space} (per-stage widths/depths from a quantized linear rule); \textbf{Y} adds Squeeze-and-Excitation. DeiT uses RegNetY-16GF (84M, $\approx$84\% top-1).
        \item CNN teacher \textbf{beats} a Transformer teacher: its convolutional inductive bias transfers priors the ViT student lacks.
      \end{itemize}}
    \end{column}
  \end{columns}
\end{frame}

% -----------------------------------------------------------------------
\begin{frame}{DeiT: Data Augmentation Recipe}

  Heavy augmentation substitutes for large-scale pre-training. All techniques were detailed in the ViT training section; the \textbf{DeiT-specific settings}:

  \vspace{0.4em}
  \begin{columns}[c]
    \begin{column}{0.62\textwidth}
      \begin{description}
        \item[RandAugment] \cite{cubuk2020randaugment} $N{=}2$ ops at magnitude $M{=}9$ (pool of 14).
        \item[Random Erasing] \cite{zhong2020randomerasing} blank a random rectangle ($p{=}0.25$); occlusion robustness.
        \item[Mixup] \cite{zhang2018mixup} blend image+label pairs, $\lambda\sim\mathrm{Beta}(0.8,0.8)$.
        \item[CutMix] \cite{yun2019cutmix} paste a box, mix labels by area.
        \item[Repeated aug.] \cite{hoffer2020augment} $m{=}3$ augmented views of each image per batch.
      \end{description}
      \vspace{0.3em}
      {\footnotesize Plus random-resized-crop $+$ horizontal flip, label smoothing $0.1$, and stochastic depth.}
    \end{column}
    \begin{column}{0.38\textwidth}
      \begin{center}
        {\footnotesize \textbf{Random Erasing} (new here)}\\[0.5em]
        \begin{tikzpicture}[scale=0.5, every node/.style={scale=0.5}]
          \begin{scope}\clip (0,0) rectangle (2,2);\visionSketchScene\end{scope}
          \draw[thick] (0,0) rectangle (2,2);
          \draw[-Stealth, thick] (2.2,1) -- (3.0,1);
          \node[font=\small, above] at (2.6,1) {erase};
          \begin{scope}[shift={(3.2,0)}]
            \begin{scope}\clip (0,0) rectangle (2,2);\visionSketchScene\end{scope}
            \fill[gray!55] (0.75,0.7) rectangle (1.65,1.35);
            \draw[thick] (0,0) rectangle (2,2);
          \end{scope}
        \end{tikzpicture}
      \end{center}
    \end{column}
  \end{columns}
\end{frame}

% -----------------------------------------------------------------------
\begin{frame}{DeiT: ImageNet-1K Results}

  Same ViT-B architecture: the gains come entirely from the \textbf{training recipe} (heavy augmentation $+$ distillation), with \textbf{no external data}. Top-1 accuracy on ImageNet-1K val.

  \vspace{0.3em}
  \begin{columns}[T]
    \begin{column}{0.56\textwidth}
      \begin{center}
        \footnotesize
        \renewcommand{\arraystretch}{1.15}
        \begin{tabular}{l r c r}
          \toprule
          \textbf{Model} & \textbf{Params} & \textbf{Extra data} & \textbf{Top-1} \\
          \midrule
          \multicolumn{4}{l}{\textit{ImageNet-1K only (from scratch)}}\\
          ViT-B/16              & 86M & none & 77.9 \\
          DeiT-S                & 22M & none & 79.8 \\
          DeiT-B                & 86M & none & 81.8 \\
          DeiT-B \distill       & 86M & none & \textbf{83.4} \\
          DeiT-B \distill\,$\uparrow$384 & 86M & none & \textbf{85.2} \\
          \midrule
          \multicolumn{4}{l}{\textit{Reference (large-scale pre-training)}}\\
          ViT-B/16              & 86M & IN-21k    & 84.0 \\
          ViT-L/16              & 307M & JFT-300M & 87.8 \\
          \midrule
          \multicolumn{4}{l}{\textit{Reference CNNs}}\\
          RegNetY-16GF (teacher) & 84M & none & 82.9 \\
          EfficientNet-B7        & 66M & none & 84.3 \\
          \bottomrule
        \end{tabular}
      \end{center}
      {\scriptsize \distill\ $=$ distillation token; $\uparrow$384 $=$ fine-tuned at $384^2$.}
    \end{column}
    \begin{column}{0.44\textwidth}
      \textbf{Takeaways}
      \begin{itemize}
        \item DeiT-B \textbf{matches JFT/21k-scale ViT} using ImageNet-1K alone: the recipe substitutes for $\sim$100$\times$ more data.
        \item \textbf{Distillation token}: $81.8 \rightarrow 83.4$ ($+1.6$); the student even surpasses its CNN teacher (82.9).
        \item \textbf{Fine-tune at $384^2$}: $+1.8$ more.
        \item \textbf{Augmentation is essential}: ablating the stack (RandAugment, Mixup/CutMix, repeated aug, stochastic depth) collapses accuracy by several points; without it ViT-B overfits ImageNet-1K.
        \item Trains in $\sim$3 days on one 8-GPU node.
      \end{itemize}
    \end{column}
  \end{columns}
\end{frame}

% =======================================================================
\subsection{MAE: Masked Autoencoders}

% -----------------------------------------------------------------------
\begin{frame}[fragile]{MAE: Masked Autoencoders Are Scalable Vision Learners}

  \cite{he2022mae}: Self-supervised pre-training by masking \textbf{75\%} of image patches and reconstructing them.

  \vspace{0.3em}
  \begin{center}
    \begin{tikzpicture}[
      scale=0.72, every node/.style={scale=0.72},
      box/.style={draw, rounded corners, fill=#1, minimum width=1.6cm, minimum height=0.5cm, align=center, font=\small},
      arr/.style={-Stealth, thick},
    ]
      % --- Image patches (sketch scene with patch grid) ---
      \begin{scope}[shift={(-6.5, -1.0)}]
        \node[font=\scriptsize\bfseries] at (1.0, 2.3) {Image patches};
        \begin{scope}
          \clip (0,0) rectangle (2,2);
          \visionSketchScene
        \end{scope}
        \draw[thick] (0,0) rectangle (2,2);
        \foreach \i in {1,2,3} {
          \pgfmathsetmacro{\xpos}{0.5*\i}
          \draw[gray!70] (\xpos, 0) -- (\xpos, 2);
          \draw[gray!70] (0, \xpos) -- (2, \xpos);
        }
      \end{scope}

      % Arrow
      \draw[arr] (-4.3, 0) -- (-3.2, 0);
      \node[font=\scriptsize, text=gray] at (-3.75, 0.3) {mask 75\%};

      % --- After masking (sketch scene with 12/16 patches masked) ---
      \begin{scope}[shift={(-3.0, -1.0)}]
        \node[font=\scriptsize\bfseries] at (1.0, 2.3) {After masking};
        \begin{scope}
          \clip (0,0) rectangle (2,2);
          \visionSketchScene
        \end{scope}
        % Mask 12 of 16 patches; leave 4 visible: (r=0,c=1), (r=1,c=3), (r=2,c=0), (r=3,c=2)
        \foreach \r in {0,...,3} {
          \foreach \c in {0,...,3} {
            \pgfmathtruncatemacro{\visible}{
              (\r==0 && \c==1) ? 1 :
              (\r==1 && \c==3) ? 1 :
              (\r==2 && \c==0) ? 1 :
              (\r==3 && \c==2) ? 1 : 0}
            \pgfmathsetmacro{\xx}{0.5*\c}
            \pgfmathsetmacro{\yy}{0.5*(3-\r)}
            \ifnum\visible=0
              \fill[gray!55] (\xx, \yy) rectangle (\xx+0.5, \yy+0.5);
              \node[font=\tiny, text=gray!90] at (\xx+0.25, \yy+0.25) {$\times$};
            \fi
          }
        }
        \draw[thick] (0,0) rectangle (2,2);
        \foreach \i in {1,2,3} {
          \pgfmathsetmacro{\xpos}{0.5*\i}
          \draw[gray!70] (\xpos, 0) -- (\xpos, 2);
          \draw[gray!70] (0, \xpos) -- (2, \xpos);
        }
      \end{scope}

      % Arrow to encoder (only visible patches)
      \draw[arr, green!50!black] (-0.5, 0) -- (0.5, 0);
      \node[font=\scriptsize, text=green!50!black, align=center] at (0.0, 0.5) {visible\\only};

      % --- Encoder ---
      \node[box=red!12, minimum width=1.8cm, minimum height=2.2cm] (enc) at (2.0, 0) {};
      \node[font=\small, text=red!60!black] at (2.0, 0.5) {\textbf{ViT}};
      \node[font=\small, text=red!60!black] at (2.0, 0.1) {\textbf{Encoder}};
      \node[font=\tiny, text=gray] at (2.0, -0.4) {(large)};
      \node[font=\tiny, text=gray] at (2.0, -0.7) {25\% tokens};

      % Arrow: insert mask tokens
      \draw[arr] (3.0, 0) -- (4.2, 0);
      \node[font=\scriptsize, text=violet, align=center] at (3.6, 0.55) {insert\\mask tokens};

      % --- Decoder ---
      \node[box=blue!12, minimum width=1.8cm, minimum height=1.6cm] (dec) at (5.5, 0) {};
      \node[font=\small, text=blue!60!black] at (5.5, 0.3) {\textbf{Decoder}};
      \node[font=\tiny, text=gray] at (5.5, -0.2) {(small: 8 layers)};
      \node[font=\tiny, text=gray] at (5.5, -0.5) {100\% tokens};

      % Arrow to reconstruction
      \draw[arr] (6.5, 0) -- (7.5, 0);

      % --- Reconstruction (full sketch with reconstructed patches tinted) ---
      \begin{scope}[shift={(7.7, -1.0)}]
        \node[font=\scriptsize\bfseries] at (1.0, 2.3) {Reconstruction};
        \begin{scope}
          \clip (0,0) rectangle (2,2);
          \visionSketchScene
        \end{scope}
        % Tint the 12 reconstructed patches (everything except the 4 visible)
        \foreach \r in {0,...,3} {
          \foreach \c in {0,...,3} {
            \pgfmathtruncatemacro{\visible}{
              (\r==0 && \c==1) ? 1 :
              (\r==1 && \c==3) ? 1 :
              (\r==2 && \c==0) ? 1 :
              (\r==3 && \c==2) ? 1 : 0}
            \pgfmathsetmacro{\xx}{0.5*\c}
            \pgfmathsetmacro{\yy}{0.5*(3-\r)}
            \ifnum\visible=0
              \fill[orange!50, opacity=0.45] (\xx, \yy) rectangle (\xx+0.5, \yy+0.5);
            \fi
          }
        }
        \draw[thick] (0,0) rectangle (2,2);
        \foreach \i in {1,2,3} {
          \pgfmathsetmacro{\xpos}{0.5*\i}
          \draw[gray!70] (\xpos, 0) -- (\xpos, 2);
          \draw[gray!70] (0, \xpos) -- (2, \xpos);
        }
        \node[font=\scriptsize, text=red!60!black] at (1.0, -0.5) {MSE loss on \textcolor{orange!70!black}{masked} only};
      \end{scope}
    \end{tikzpicture}
  \end{center}

  \vspace{0.2em}
  \begin{itemize}
    \item \textbf{Encoder}: standard ViT, processes only \textcolor{green!50!black}{visible patches} (25\%) --- $3\times$+ speedup
    \item \textbf{Decoder}: lightweight (8 layers, width 512), operates on all tokens incl.\ shared learned \textcolor{violet}{mask tokens}
    \item \textbf{Target}: per-patch normalized pixel values; loss $= \frac{1}{|\mathcal{M}|}\sum_{i \in \mathcal{M}} \|\hat{x}_i - x_i\|^2$ on masked patches only
  \end{itemize}
\end{frame}

% -----------------------------------------------------------------------
\begin{frame}{MAE: Decoder \& Reconstruction Loss}

  The decoder is responsible for predicting raw pixels for the missing 75\% of patches.

  \vspace{0.3em}
  \begin{columns}[T]
    \begin{column}{0.55\textwidth}
      \textbf{Decoder input} (full sequence, $N{+}1$ tokens):
      \begin{itemize}\itemsep0pt
        \item Encoded \textcolor{green!50!black}{visible patch} embeddings (from encoder)
        \item One \textcolor{violet}{mask token} $[M]$ inserted at each masked position
        \item $[M]$ is a \emph{single shared learned vector} reused at every masked position
        \item \textbf{Positional encoding}: \textbf{fixed 2D sin-cos} (non-learned; row and column sinusoids concatenated), added to \textbf{all} tokens. Since every $[M]$ is identical, $E_\text{pos}$ is the decoder's \emph{only} cue to where a masked patch belongs.
      \end{itemize}

      \vspace{0.3em}
      \textbf{Architecture}: lightweight Transformer
      \begin{itemize}\itemsep0pt
        \item 8 blocks, hidden dim 512 ($\sim 9\%$ of encoder FLOPs)
        \item Standard pre-norm MSA + MLP blocks
        \item Asymmetric design: only the \emph{encoder} is kept for downstream tasks; decoder is \textbf{discarded} after pre-training
      \end{itemize}

      \vspace{0.3em}
      \textbf{Output head}: a single linear layer maps each output token to a vector $\hat{x}_i \in \mathbb{R}^{P^2 C}$, the predicted pixel values of patch $i$ (reshaped to $P{\times}P{\times}C$).
    \end{column}
    \begin{column}{0.42\textwidth}
      \begin{center}
        \begin{tikzpicture}[
          scale=0.7, every node/.style={scale=0.7},
          tk/.style={draw, minimum width=0.32cm, minimum height=0.4cm, inner sep=0pt, font=\tiny, rounded corners=1pt},
          arr/.style={-Stealth, thick},
        ]
          % Input tokens row: mix of visible (0) and mask (1)
          \node[font=\scriptsize, anchor=base] at (2, 5.6) {Decoder input ($N{+}1$ tokens)};
          \foreach \i/\type in {0/1, 1/0, 2/1, 3/1, 4/0, 5/1, 6/1, 7/0, 8/1, 9/1} {
            \pgfmathsetmacro{\xx}{0.4 + \i*0.35}
            \ifnum\type=0
              \node[tk, fill=green!25] at (\xx, 5) {};
            \else
              \node[tk, fill=violet!25, font=\tiny] at (\xx, 5) {$\scriptscriptstyle[M]$};
            \fi
          }
          \node[font=\scriptsize, text=orange!70!black] at (2, 4.55) {$+\,E_\text{pos}$ (on all tokens)};

          \draw[arr] (2, 4.25) -- (2, 3.85);

          \node[draw, rounded corners, fill=blue!12, minimum width=3.0cm, minimum height=1.0cm, align=center, font=\small] (dec) at (2, 3.3) {Decoder\\\scriptsize 8 blocks, dim 512};

          \draw[arr] (2, 2.75) -- (2, 2.35);

          \node[draw, rounded corners, fill=yellow!20, minimum width=3.0cm, font=\scriptsize] at (2, 2.05) {Linear head (per token)};

          \draw[arr] (2, 1.75) -- (2, 1.35);

          % Predicted patches row
          \foreach \i in {0,...,9} {
            \pgfmathsetmacro{\xx}{0.4 + \i*0.35}
            \node[tk, fill=orange!25, font=\tiny] at (\xx, 1.0) {};
          }
          \node[font=\scriptsize] at (2, 0.55) {$\hat{x}_i \in \mathbb{R}^{P^2 C}$};

          \draw[arr, red!60!black] (2, 0.25) -- (2, -0.2);
          \node[font=\scriptsize, text=red!60!black, align=center] at (2, -0.5) {MSE loss\\(masked only)};
        \end{tikzpicture}
      \end{center}
    \end{column}
  \end{columns}

  \vspace{0.3em}
  \textbf{Training loss}: per-patch MSE, computed \textbf{only on masked positions} (visible positions are ignored, like BERT):
  \[
    \mathcal{L} = \frac{1}{|\mathcal{M}|} \sum_{i \in \mathcal{M}} \bigl\| \hat{x}_i - \mathrm{norm}(x_i) \bigr\|_2^2
  \]
  where $\mathrm{norm}(x_i) = (x_i - \mu_i)/\sigma_i$ uses the patch's \emph{own} mean and std. This per-patch normalization improves representation quality compared to raw-pixel targets.
\end{frame}

% -----------------------------------------------------------------------
\begin{frame}{MAE: Design Choices \& Results}

  \begin{columns}
    \begin{column}{0.48\textwidth}
      \textbf{Why 75\% masking?}
      \begin{itemize}
        \item Images have high spatial redundancy. Low masking allows trivial interpolation
        \item 75\% forces the model to learn \emph{semantic} understanding, not local texture
        \item Much higher than BERT's 15\% for text
      \end{itemize}

      \vspace{0.3em}
      \textbf{Asymmetric encoder-decoder}:
      \begin{itemize}
        \item Encoder: large ViT (e.g.\ 24 layers, 1024 dim)
        \item Decoder: small (8 layers, 512 dim)
        \item Decoder \textbf{discarded} after pre-training
        \item Only encoder used for downstream tasks
      \end{itemize}

      \vspace{0.3em}
      \textbf{Downstream usage}:
      \begin{description}[style=nextline, leftmargin=0.5cm]
        \item[Fine-tuning] Full encoder + classification head; best results
        \item[Linear probing] Frozen encoder + linear layer; weaker than contrastive methods but gap narrows at scale
      \end{description}
    \end{column}
    \begin{column}{0.48\textwidth}
      \textbf{ImageNet-1K results} (fine-tuning):
      \begin{center}
        \small
        \begin{tabular}{lcc}
          \toprule
          \textbf{Model} & \textbf{Params} & \textbf{Top-1} \\
          \midrule
          ViT-B/16 & 86M & 83.6\% \\
          ViT-L/16 & 307M & 85.9\% \\
          ViT-H/14 & 632M & 86.9\% \\
          \bottomrule
        \end{tabular}
      \end{center}
      Using ImageNet-1K only. No external data.

      \vspace{0.3em}
      \textbf{Key}: MAE scales excellently to ViT-H; contrastive methods plateau at larger sizes.
    \end{column}
  \end{columns}
\end{frame}

% =======================================================================
\subsection{Swin Transformer}

% -----------------------------------------------------------------------
\begin{frame}{Swin Transformer: Motivation \& Key Ideas}

  \cite{liu2021swin}: make the Transformer an \textbf{efficient backbone} for dense tasks (detection, segmentation), not only image-level classification. Plain ViT \emph{can} be adapted to dense prediction (SETR, DPT, later ViTDet), but it is single-scale and quadratic in resolution; Swin builds the multi-scale hierarchy and linear cost in by design.

  \vspace{0.4em}
  \begin{columns}[T]
    \begin{column}{0.46\textwidth}
      \textbf{Plain ViT as a dense backbone}
      \begin{itemize}\itemsep3pt
        \item Global attention is $\mathcal{O}(n^2)$ in the \#tokens $\to$ costly at the high resolutions dense tasks need
        \item Single fixed-resolution feature map ($16{\times}$ downsample) $\to$ multi-scale features must be bolted on afterwards
        \item Weak locality prior $\to$ data hungry
      \end{itemize}
    \end{column}
    \begin{column}{0.52\textwidth}
      \textbf{Swin's ideas}
      \begin{itemize}\itemsep3pt
        \item \textbf{Window attention (W-MSA)}: attend only within fixed $M{\times}M$ windows $\to$ \textbf{linear} cost $+$ locality
        \item \textbf{Shifted windows (SW-MSA)}: shift the partition every other block $\to$ cross-window flow, still linear
        \item \textbf{Hierarchical stages}: patch merging halves resolution \& doubles channels $\to$ CNN-like feature pyramid
      \end{itemize}
    \end{column}
  \end{columns}

  \vspace{0.4em}
  \textbf{In one line}: local + shifted window attention for linear cost and cross-window mixing, stacked hierarchically for multi-scale features.
\end{frame}

% -----------------------------------------------------------------------
\begin{frame}{Swin: Window \& Shifted Window Attention}

  \begin{columns}
    \begin{column}{0.48\textwidth}
      \textbf{Window-MSA (W-MSA)}:
      \begin{itemize}
        \item Partition feature map into non-overlapping $M{\times}M$ windows ($M{=}7$)
        \item Self-attention computed \emph{within} each window
        \item Complexity: $\mathcal{O}(M^2 \cdot hw \cdot C)$, \textbf{linear} in image size
      \end{itemize}

      \vspace{0.3em}
      \textbf{Problem}: No cross-window information flow.

      \vspace{0.3em}
      \textbf{Shifted Window-MSA (SW-MSA)}:
      \begin{itemize}
        \item Shift window grid by $(\lfloor M/2\rfloor, \lfloor M/2\rfloor)$
        \item Windows now straddle previous boundaries
        \item Consecutive blocks alternate W-MSA / SW-MSA
      \end{itemize}

      \vspace{0.3em}
      \textbf{Efficient cyclic shift}:
      \begin{enumerate}
        \item Cyclically shift feature map
        \item Regular window partition + attention masks
        \item Reverse shift
      \end{enumerate}
    \end{column}
    \begin{column}{0.48\textwidth}
      \begin{center}
        \begin{tikzpicture}[scale=0.62, every node/.style={scale=0.62}]
          % --- Layer l: 9 regular windows ---
          \def\cs{0.5}
          \node[font=\small\bfseries] at (1.5,3.95) {Layer $\ell$: W-MSA};
          \foreach \wx in {0,1,2}{\foreach \wy in {0,1,2}{
            \pgfmathtruncatemacro{\t}{mod(\wx+\wy,2)}
            \ifnum\t=0 \fill[blue!10] (\wx,\wy) rectangle (\wx+1,\wy+1);
            \else \fill[red!8] (\wx,\wy) rectangle (\wx+1,\wy+1);\fi}}
          \draw[step=\cs, black!22, thin] (0,0) grid (3,3);
          \foreach \b in {1,2}{
            \draw[very thick, red!70!black] (\b,0) -- (\b,3);
            \draw[very thick, red!70!black] (0,\b) -- (3,\b);}
          \draw[thick] (0,0) rectangle (3,3);
          \foreach \wx/\nx in {0/1,1/2,2/3}{\foreach \wy/\ny in {2/1,1/4,0/7}{
            \pgfmathtruncatemacro{\lab}{\nx+\ny-1}
            \node[font=\scriptsize] at (\wx+0.5,\wy+0.5) {W\lab};}}
          \node[font=\scriptsize, text=red!70!black, align=center] at (1.5,-0.4)
            {9 windows, attention \emph{within} each};
        \end{tikzpicture}

        \vspace{0.25cm}

        \begin{tikzpicture}[scale=0.62, every node/.style={scale=0.62}]
          % --- Layer l+1: shifted windows (9-window grid shifted by M/2) ---
          \def\cs{0.5}
          \node[font=\small\bfseries] at (1.5,3.95) {Layer $\ell{+}1$: SW-MSA};
          \foreach \xa/\xb in {0/0.5,0.5/1.5,1.5/2.5,2.5/3}{
            \foreach \ya/\yb in {0/0.5,0.5/1.5,1.5/2.5,2.5/3}{
              \pgfmathsetmacro{\cx}{(\xa+\xb)/2}\pgfmathsetmacro{\cy}{(\ya+\yb)/2}
              \pgfmathtruncatemacro{\t}{mod(round(\cx*2)+round(\cy*2),2)}
              \ifnum\t=0 \fill[green!12] (\xa,\ya) rectangle (\xb,\yb);
              \else \fill[orange!12] (\xa,\ya) rectangle (\xb,\yb);\fi}}
          \draw[step=\cs, black!22, thin] (0,0) grid (3,3);
          \foreach \b in {1,2}{
            \draw[dashed, red!45] (\b,0) -- (\b,3);
            \draw[dashed, red!45] (0,\b) -- (3,\b);}
          \foreach \b in {0.5,1.5,2.5}{
            \draw[very thick, blue!70!black] (\b,0) -- (\b,3);
            \draw[very thick, blue!70!black] (0,\b) -- (3,\b);}
          \draw[thick] (0,0) rectangle (3,3);
          \draw[-Stealth, violet, thick] (-0.4,3) -- (-0.4,2.5) node[midway,left,font=\tiny]{$\tfrac{M}{2}$};
          \draw[-Stealth, violet, thick] (3.0,3.35) -- (3.5,3.35) node[midway,above,font=\tiny]{$\tfrac{M}{2}$};
          \node[font=\scriptsize, text=blue!60!black, align=center] at (1.5,-0.4)
            {windows straddle old borders\\$\Rightarrow$ cross-window mixing};
        \end{tikzpicture}
      \end{center}
    \end{column}
  \end{columns}

  \vspace{0.2em}
  \textbf{Relative position bias}: $\operatorname{Attn}(Q,K,V) = \operatorname{softmax}\!\bigl(\frac{QK^\top}{\sqrt{d}} + B\bigr)\,V$,\quad $B$ from learnable $(2M{-}1){\times}(2M{-}1)$ table.
\end{frame}

% -----------------------------------------------------------------------
\begin{frame}[fragile]{Swin: How Window Attention Works (W-MSA)}

  W-MSA is just \textbf{ordinary self-attention run independently inside each window}. A window of $M{\times}M$ tokens is flattened to a length-$M^2$ sequence; tokens attend to all $M^2$ peers in that window and to nothing outside.

  \vspace{0.1em}
  \begin{center}
    \begin{tikzpicture}[
        scale=0.82, every node/.style={scale=0.82},
        font=\small,
        arr/.style={-Stealth, thick},
        att/.style={-Stealth, red!70!black, line width=0.5pt, opacity=0.85},
      ]
      % ---------- Part A: feature map partitioned into windows ----------
      \def\cs{0.5}
      \fill[red!8]   (0,0)     rectangle (1.5,1.5);
      \fill[gray!6]  (1.5,0)   rectangle (3,1.5);
      \fill[gray!6]  (0,1.5)   rectangle (1.5,3);
      \fill[gray!6]  (1.5,1.5) rectangle (3,3);
      \draw[step=\cs, black!20, thin] (0,0) grid (3,3);
      \draw[very thick, red!70!black] (1.5,0) -- (1.5,3);
      \draw[very thick, red!70!black] (0,1.5) -- (3,1.5);
      \draw[thick] (0,0) rectangle (3,3);
      \foreach \c in {0,...,5}{\foreach \r in {0,...,5}{
        \fill[black!45] (\c*\cs+\cs/2,\r*\cs+\cs/2) circle (1.1pt);}}
      \coordinate (q) at (1*\cs+\cs/2,1*\cs+\cs/2);
      \foreach \c in {0,1,2}{\foreach \r in {0,1,2}{
        \draw[att] (q) -- (\c*\cs+\cs/2,\r*\cs+\cs/2);}}
      \fill[red!90!black] (q) circle (2.2pt);
      \node[font=\scriptsize, text=red!70!black, anchor=north] at (0.75,-0.12) {query};
      \draw[gray, dashed, -Stealth] (q) -- (2.25,2.25);
      \draw[red, line width=1pt] (1.62,1.62) -- (1.92,1.92);
      \draw[red, line width=1pt] (1.62,1.92) -- (1.92,1.62);
      \node[font=\scriptsize, text=gray, align=center, anchor=west] at (2.4,2.4) {no attention\\across windows};
      \node[font=\small\bfseries] at (1.5,3.55) {Feature map, $M{\times}M$ windows ($M{=}3$)};
      \node[font=\scriptsize, text=red!70!black, align=center] at (1.5,-0.7)
        {each token attends to \textbf{all $M^2$ tokens in its own window}};
    \end{tikzpicture}
  \end{center}

  \vspace{0.1em}
  {\small Same $Q,K,V$ projections and softmax as a Transformer, only the attention is \textbf{block-diagonal} over windows: an $hw{\times}hw$ attention map collapses to $\tfrac{hw}{M^2}$ independent $M^2{\times}M^2$ maps $\Rightarrow$ cost grows \textbf{linearly}, not quadratically, in $hw$.}
\end{frame}

% -----------------------------------------------------------------------
\begin{frame}[fragile]{Swin Transformer: The Swin Block}

  \cite{liu2021swin}: replace global attention by \textbf{local window attention}. The basic unit is a \textbf{pair} of pre-norm transformer blocks that always come together: the first uses \textbf{W-MSA}, the second \textbf{SW-MSA}.

  \vspace{0.2em}
  \begin{columns}[T]
    \begin{column}{0.40\textwidth}
      \begin{center}
        \begin{tikzpicture}[
            scale=0.52, every node/.style={scale=0.52},
            op/.style={draw, rounded corners=2pt, minimum width=2.0cm, minimum height=0.5cm, align=center, font=\small},
            msa/.style={op, fill=#1},
            ln/.style={op, fill=gray!12},
            mlp/.style={op, fill=blue!12},
            add/.style={draw, circle, inner sep=1.2pt, font=\footnotesize, fill=white},
            arr/.style={-Stealth, thick},
            res/.style={-Stealth, thick, gray!55},
          ]
          \def\x{0}
          \node[font=\small] (in) at (\x,0) {$\mathbf{z}^{\ell-1}$};
          \node[ln]  (ln1)  at (\x,0.9) {LN};
          \node[msa=red!18] (wmsa) at (\x,1.8) {W-MSA};
          \node[add] (a1) at (\x,2.7) {$+$};
          \node[ln]  (ln2)  at (\x,3.6) {LN};
          \node[mlp] (mlp1) at (\x,4.5) {MLP};
          \node[add] (a2) at (\x,5.4) {$+$};
          \node[ln]  (ln3)  at (\x,6.4) {LN};
          \node[msa=orange!22] (swmsa) at (\x,7.3) {SW-MSA};
          \node[add] (a3) at (\x,8.2) {$+$};
          \node[ln]  (ln4)  at (\x,9.1) {LN};
          \node[mlp] (mlp2) at (\x,10.0) {MLP};
          \node[add] (a4) at (\x,10.9) {$+$};
          \node[font=\small] (out) at (\x,11.8) {$\mathbf{z}^{\ell+1}$};

          \draw[arr] (in) -- (ln1);  \draw[arr] (ln1) -- (wmsa);
          \draw[arr] (wmsa) -- (a1); \draw[arr] (a1) -- (ln2);
          \draw[arr] (ln2) -- (mlp1);\draw[arr] (mlp1) -- (a2);
          \draw[arr] (a2) -- (ln3);  \draw[arr] (ln3) -- (swmsa);
          \draw[arr] (swmsa) -- (a3);\draw[arr] (a3) -- (ln4);
          \draw[arr] (ln4) -- (mlp2);\draw[arr] (mlp2) -- (a4);
          \draw[arr] (a4) -- (out);

          \def\rx{1.5}
          \draw[res] (in.east) -- (\rx,0)   -- (\rx,2.7)  -- (a1.east);
          \draw[res] (a1.east) -- (\rx,2.7) -- (\rx,5.4)  -- (a2.east);
          \draw[res] (a2.east) -- (\rx,5.4) -- (\rx,8.2)  -- (a3.east);
          \draw[res] (a3.east) -- (\rx,8.2) -- (\rx,10.9) -- (a4.east);

          \def\bx{-1.65}
          \draw[decorate, decoration={brace, amplitude=5pt}, gray]
                (\bx,1.15) -- (\bx,5.55)
                node[midway, left=5pt, font=\scriptsize, align=center, text=red!55!black] {block\\$\ell$\\(W-MSA)};
          \draw[decorate, decoration={brace, amplitude=5pt}, gray]
                (\bx,6.65) -- (\bx,11.05)
                node[midway, left=5pt, font=\scriptsize, align=center, text=orange!60!black] {block\\$\ell{+}1$\\(SW-MSA)};
        \end{tikzpicture}
      \end{center}
    \end{column}
    \begin{column}{0.58\textwidth}
      \textbf{Two-block unit} (pre-norm, standard residual + MLP):
      \begin{align*}
        \hat{\mathbf{z}}^\ell     &= \text{W-MSA}(\text{LN}(\mathbf{z}^{\ell-1})) + \mathbf{z}^{\ell-1} \\
        \mathbf{z}^\ell           &= \text{MLP}(\text{LN}(\hat{\mathbf{z}}^\ell)) + \hat{\mathbf{z}}^\ell \\
        \hat{\mathbf{z}}^{\ell+1} &= \text{SW-MSA}(\text{LN}(\mathbf{z}^\ell)) + \mathbf{z}^\ell \\
        \mathbf{z}^{\ell+1}       &= \text{MLP}(\text{LN}(\hat{\mathbf{z}}^{\ell+1})) + \hat{\mathbf{z}}^{\ell+1}
      \end{align*}

      \begin{description}[style=nextline, leftmargin=0.4cm]
        \item[W-MSA] attention inside fixed $M{\times}M$ windows: cheap, but windows never talk.
        \item[SW-MSA] same, but the window grid is \emph{shifted} so the next block's windows straddle the previous boundaries $\Rightarrow$ cross-window information flow.
      \end{description}
      Blocks are \textbf{always added in pairs}. $(2,2,6,2)$ = blocks per stage (Swin-T); even, since W-MSA/SW-MSA come together. Mechanism of the shift: next slides.
    \end{column}
  \end{columns}
\end{frame}

% -----------------------------------------------------------------------
\begin{frame}[fragile]{Swin Transformer: Hierarchical Feature Maps}

  Stacking Swin blocks between \textbf{Patch Merging} layers yields a \textbf{CNN-like feature pyramid}: resolution shrinks, channels grow.

  \vspace{0.2em}
  \begin{center}
    \begin{tikzpicture}[
        scale=0.74, every node/.style={scale=0.74},
        font=\small,
        arr/.style={-Stealth, thick},
        lbl/.style={font=\scriptsize, text=gray},
        op/.style={font=\scriptsize, text=blue!55!black, align=center},
      ]
      \def\base{1.2}
      \def\srow{3.05}
      % feature volume centred on \base; sets \rgt to its right edge
      \newcommand{\fvol}[5]{%
        \pgfmathsetmacro{\side}{#2*#3}%
        \pgfmathsetmacro{\yb}{\base-\side/2}%
        \foreach \d in {#4,...,1}{%
          \pgfmathsetmacro{\off}{(\d-1)*0.06}%
          \fill[#5, draw=black!55] (#1+\off,\yb+\off) rectangle (#1+\off+\side,\yb+\off+\side);%
        }%
        \foreach \i in {0,...,#2}{
          \draw[black!30, thin] (#1+\i*#3,\yb) -- (#1+\i*#3,\yb+\side);
          \draw[black!30, thin] (#1,\yb+\i*#3) -- (#1+\side,\yb+\i*#3);
        }
        \draw[black!70, thick] (#1,\yb) rectangle (#1+\side,\yb+\side);
        \pgfmathsetmacro{\rgt}{#1+\side+(#4-1)*0.06}%
        \global\let\rgt\rgt
      }

      \node[align=center] (img) at (-0.5,\base) {Image\\$224^2{\times}3$};

      \fvol{1.5}{8}{0.30}{2}{red!15}      \edef\rA{\rgt}
      \node[lbl, align=center] at (2.7,-0.35) {$\tfrac{H}{4}{\times}\tfrac{W}{4}$,\; $56^2$\\$C$};
      \node[font=\scriptsize\bfseries, text=red!55!black] at (2.7,\srow) {Stage 1};

      \fvol{5.2}{4}{0.30}{3}{orange!18}   \edef\rB{\rgt}
      \node[lbl, align=center] at (5.95,-0.35) {$\tfrac{H}{8}{\times}\tfrac{W}{8}$,\; $28^2$\\$2C$};
      \node[font=\scriptsize\bfseries, text=orange!60!black] at (5.9,\srow) {Stage 2};

      \fvol{7.8}{2}{0.40}{4}{yellow!35}   \edef\rC{\rgt}
      \node[lbl, align=center] at (8.4,-0.35) {$\tfrac{H}{16}{\times}\tfrac{W}{16}$,\; $14^2$\\$4C$};
      \node[font=\scriptsize\bfseries, text=olive] at (8.4,\srow) {Stage 3};

      \fvol{10.0}{1}{0.55}{5}{green!22}   \edef\rD{\rgt}
      \node[lbl, align=center] at (10.5,-0.35) {$\tfrac{H}{32}{\times}\tfrac{W}{32}$,\; $7^2$\\$8C$};
      \node[font=\scriptsize\bfseries, text=green!45!black] at (10.5,\srow) {Stage 4};

      \draw[arr] (0.5,\base) -- (1.45,\base);
      \node[op] at (1.0,\base+0.5) {Patch\\Partition};
      \draw[arr] (\rA+0.05,\base) -- (5.15,\base);  \node[op] at (4.55,\base+0.55) {Patch\\Merge};
      \draw[arr] (\rB+0.05,\base) -- (7.75,\base);  \node[op] at (7.15,\base+0.55) {Patch\\Merge};
      \draw[arr] (\rC+0.05,\base) -- (9.95,\base);  \node[op] at (9.4,\base+0.55) {Patch\\Merge};
      \draw[arr] (\rD+0.05,\base) -- (11.7,\base);
      \node[align=center] at (12.5,\base) {Global Avg\\Pool (GAP)\\$\to$ class};

      \draw[decorate, decoration={brace, amplitude=4pt}, gray]
            (1.5,\srow+0.55) -- (10.6,\srow+0.55)
            node[midway, above=4pt, font=\scriptsize, text=blue!55!black, align=center]
            {each Patch Merge: resolution $\downarrow 2{\times}$, channels $\uparrow 2{\times}$\quad(vs.\ ViT's flat $14^2$ grid)};
    \end{tikzpicture}
  \end{center}

  \vspace{0.1em}
  \begin{description}[style=nextline, leftmargin=2cm]
    \item[Patch Partition] Split image into $4{\times}4$ non-overlapping patches ($56{\times}56$ tokens, $C{=}96$)
    \item[Patch Merging] Concatenate $2{\times}2$ neighbor patches $\to$ linear projection to $2C$; halves spatial resolution, doubles channels
    \item[Stages] $2,2,6,2$ Swin blocks; hierarchical maps make Swin a drop-in \textbf{backbone} for detection and segmentation
    \item[GAP] Global Average Pooling: average the final $7^2$ feature map over all spatial positions $\to$ one $8C$ vector for the classifier (no \texttt{[CLS]} token, unlike ViT)
  \end{description}
\end{frame}

% -----------------------------------------------------------------------
\begin{frame}{Swin: Training Recipe}

  Trained \textbf{from scratch on ImageNet-1K} (no extra data); recipe mostly follows DeiT~\cite{touvron2021deit}, so gains are architectural, not from a better recipe.

  \vspace{0.3em}
  \begin{columns}[T]
    \begin{column}{0.5\textwidth}
      \textbf{Optimization} (224$^2$)
      \begin{description}[style=nextline, leftmargin=0pt, itemsep=1pt]
        \item[Optimizer] AdamW, weight decay $0.05$
        \item[Epochs] $300$, $+20$ linear warm-up
        \item[Schedule] cosine decay
        \item[Batch / LR] $1024$ / peak $10^{-3}$
        \item[Grad clip] max norm $1$
      \end{description}

      \vspace{0.3em}
      \textbf{Regularization} (scaled with model size)
      \begin{itemize}\itemsep1pt
        \item Stochastic depth: $0.2$ / $0.3$ / $0.5$ for Swin-T / S / B
        \item Label smoothing $0.1$
      \end{itemize}
    \end{column}
    \begin{column}{0.48\textwidth}
      \textbf{Augmentation} (DeiT-style)
      \begin{itemize}\itemsep1pt
        \item RandAugment, Mixup, CutMix, random erasing
        \item \emph{No} repeated augmentation, \emph{no} EMA (did not help)
      \end{itemize}

      \vspace{0.3em}
      \textbf{Larger models / resolution}
      \begin{description}[style=nextline, leftmargin=0pt, itemsep=2pt]
        \item[IN-22K pre-train] AdamW, $90$ ep ($+5$ warm-up), batch $4096$, LR $10^{-3}$, wd $0.01$
        \item[Fine-tune @ 384$^2$] $30$ ep, constant LR $10^{-5}$, wd $10^{-8}$, stoch.\ depth $0.1$
      \end{description}
    \end{column}
  \end{columns}

  \vspace{0.4em}
  {\footnotesize Swin-L (197M) overfits on IN-1K from scratch $\Rightarrow$ needs IN-22K pre-training.}
\end{frame}

% -----------------------------------------------------------------------
\begin{frame}{Swin: ImageNet-1K Results \& Analysis}

  \begin{columns}[T]
    \begin{column}{0.52\textwidth}
      \textbf{ImageNet-1K classification} (from scratch, $224^2$ unless noted):
      \begin{center}
        \small
        \begin{tabular}{lccc}
          \toprule
          \textbf{Model} & \textbf{Params} & \textbf{GFLOPs} & \textbf{Top-1} \\
          \midrule
          Swin-T            & 29M & 4.5  & 81.3 \\
          Swin-S            & 50M & 8.7  & 83.0 \\
          Swin-B            & 88M & 15.4 & 83.5 \\
          Swin-B\,$\uparrow$384 & 88M & 47.0 & \textbf{84.5} \\
          \midrule
          \textit{DeiT-B} (ref) & 86M & 17.5 & 81.8 \\
          \textit{DeiT-B}\,$\uparrow$384 & 86M & 55.4 & 83.1 \\
          \bottomrule
        \end{tabular}
      \end{center}
      {\footnotesize Configs: $C{=}96/96/128/192$, blocks/stage $(2,2,6,2)$ for T, $(2,2,18,2)$ for S/B/L. Swin-L (197M) needs ImageNet-22K pre-training ($87.3$ @384); from-scratch IN-1K overfits.}
    \end{column}
    \begin{column}{0.46\textwidth}
      \textbf{Analysis}:
      \begin{itemize}\itemsep2pt
        \item At \textbf{matched compute}, Swin $>$ DeiT: Swin-B \textbf{83.5} vs DeiT-B 81.8 ($-2.1$ FLOPs), \emph{no} distillation, IN-1K only.
        \item Gains come from \textbf{window locality + hierarchy}, not extra data.
        \item Linear-in-resolution cost $\Rightarrow$ fine-tuning at $384^2$ is cheap and adds $+1$.
        \item \textbf{Main payoff is dense prediction}, where the built-in pyramid plugs straight into existing detectors/segmenters:
        \begin{itemize}\footnotesize
          \item COCO det.: Mask R-CNN + Swin-T $\approx 46$ box AP
          \item ADE20K seg.: UperNet + Swin-T $\approx 45$ mIoU
        \end{itemize}
        Hierarchical multi-scale features make Swin a drop-in CNN-style \textbf{backbone} (plain ViT needs extra adaptation, e.g.\ ViTDet, to match this).
      \end{itemize}
    \end{column}
  \end{columns}
\end{frame}

% =======================================================================
\subsection{MLP-Mixer}

% -----------------------------------------------------------------------
\begin{frame}[fragile]{MLP-Mixer: No Attention, Only MLPs}

  \cite{tolstikhin2021mlpmixer}: Replace attention entirely with two types of MLPs: \textbf{token-mixing} (across patches) and \textbf{channel-mixing} (within patches).

  \vspace{0.1em}
  \begin{center}
    \begin{tikzpicture}[
        scale=0.7, every node/.style={scale=0.7},
        box/.style={draw, rounded corners, fill=#1, minimum width=2.4cm, minimum height=0.5cm, align=center, font=\small},
        arr/.style={-Stealth, thick},
        res/.style={-Stealth, thick, gray!55},
        dot/.style={circle, fill=gray!55, inner sep=1.3pt},
      ]
      % --- Mixer layer main path ---
      \node[font=\small] (in) at (0,-0.3) {$X \in \mathbb{R}^{S \times C}$};
      \node[box=gray!12] (ln1)  at (0,0.6) {Layer Norm};
      \node[box=red!14]  (tmix) at (0,1.5) {Token-Mixing MLP};
      \node[draw, circle, inner sep=1pt, font=\scriptsize, fill=white] (add1) at (0,2.4) {$+$};
      \node[box=gray!12] (ln2)  at (0,3.3) {Layer Norm};
      \node[box=blue!14] (cmix) at (0,4.2) {Channel-Mixing MLP};
      \node[draw, circle, inner sep=1pt, font=\scriptsize, fill=white] (add2) at (0,5.1) {$+$};
      \node[font=\small] (out) at (0,6.0) {$Y \in \mathbb{R}^{S \times C}$};

      \draw[arr] (in) -- (ln1);   \draw[arr] (ln1) -- (tmix);
      \draw[arr] (tmix) -- (add1); \draw[arr] (add1) -- (ln2);
      \draw[arr] (ln2) -- (cmix);  \draw[arr] (cmix) -- (add2);
      \draw[arr] (add2) -- (out);

      % residual streams: equal distance from main path
      \def\rx{-1.9}
      \node[dot] (b1) at (0,0.1) {};
      \draw[res] (b1) -- (\rx,0.1) -- (\rx,2.4) -- (add1.west);
      \node[dot] (b2) at (0,2.75) {};
      \draw[res] (b2) -- (\rx,2.75) -- (\rx,5.1) -- (add2.west);
      \node[font=\scriptsize, text=gray, rotate=90] at (\rx-0.25,1.25) {skip};
      \node[font=\scriptsize, text=gray, rotate=90] at (\rx-0.25,3.9) {skip};

      % ===== striped matrices: row = token (patch), column = channel =====
      % Token-Mixing: vertical stripes = channels; mix DOWN the S axis
      \begin{scope}[shift={(3.4,0.55)}]
        \node[font=\small\bfseries, text=red!60!black] at (1.26,4.6) {Token-Mixing};
        \def\cs{0.42}
        \foreach \c in {0,...,5}{\pgfmathsetmacro{\sh}{10+12*mod(\c,2)}\fill[red!\sh] (\c*\cs,0) rectangle (\c*\cs+\cs,6*\cs);}
        \draw[step=\cs, black!25, thin] (0,0) grid (6*\cs,6*\cs);
        \draw[thick] (0,0) rectangle (6*\cs,6*\cs);
        \foreach \c in {0,...,5}{\draw[-Stealth, red!70!black, thick] (\c*\cs+\cs/2,5.7*\cs) -- (\c*\cs+\cs/2,0.3*\cs);}
        \foreach \r [evaluate=\r as \rn using int(6-\r)] in {0,...,5}{
          \node[font=\tiny, anchor=east] at (-0.08,\r*\cs+\cs/2) {token \rn};}
        \foreach \c [evaluate=\c as \cn using int(\c+1)] in {0,...,5}{
          \node[font=\tiny, rotate=90, anchor=west] at (\c*\cs+\cs/2,6*\cs+0.06) {channel \cn};}
        \node[font=\scriptsize, text=red!60!black, align=center] at (1.26,-0.75)
          {one MLP per \textbf{column} (channel),\\shared across channels;\\mixes across \textbf{tokens} ($S$)};
      \end{scope}

      % Channel-Mixing: horizontal stripes = tokens; mix ACROSS the C axis
      \begin{scope}[shift={(7.8,0.55)}]
        \node[font=\small\bfseries, text=blue!60!black] at (1.26,4.6) {Channel-Mixing};
        \def\cs{0.42}
        \foreach \r in {0,...,5}{\pgfmathsetmacro{\sh}{10+12*mod(\r,2)}\fill[blue!\sh] (0,\r*\cs) rectangle (6*\cs,\r*\cs+\cs);}
        \draw[step=\cs, black!25, thin] (0,0) grid (6*\cs,6*\cs);
        \draw[thick] (0,0) rectangle (6*\cs,6*\cs);
        \foreach \r in {0,...,5}{\draw[-Stealth, blue!70!black, thick] (0.3*\cs,\r*\cs+\cs/2) -- (5.7*\cs,\r*\cs+\cs/2);}
        \foreach \r [evaluate=\r as \rn using int(6-\r)] in {0,...,5}{
          \node[font=\tiny, anchor=east] at (-0.08,\r*\cs+\cs/2) {token \rn};}
        \foreach \c [evaluate=\c as \cn using int(\c+1)] in {0,...,5}{
          \node[font=\tiny, rotate=90, anchor=west] at (\c*\cs+\cs/2,6*\cs+0.06) {channel \cn};}
        \node[font=\scriptsize, text=blue!60!black, align=center] at (1.26,-0.75)
          {one MLP per \textbf{row} (token),\\shared across tokens;\\mixes across \textbf{channels} ($C$)};
      \end{scope}
    \end{tikzpicture}
  \end{center}

  \vspace{0.1em}
  {\small
  Each mixing block is a \textbf{2-layer MLP} (exactly a Transformer FFN): $\;\text{Linear} \to \text{GELU} \to \text{Linear}\;$ with an expansion hidden width $-$ \emph{not} a single matrix. Token-mix weights $W_1,W_2$ act on the $S$ axis; channel-mix weights $W_3,W_4$ act on the $C$ axis:}
  \[
    U_{*,i} = X_{*,i} + W_2\, \sigma(W_1\, \text{LN}(X)_{*,i}), \qquad
    Y_{j,*} = U_{j,*} + W_4\, \sigma(W_3\, \text{LN}(U)_{j,*})
  \]
\end{frame}

% -----------------------------------------------------------------------
\begin{frame}{MLP-Mixer: Architecture Comparison}

  \begin{columns}
    \begin{column}{0.48\textwidth}
      \textbf{Full pipeline}:
      \begin{enumerate}
        \item Split image into $S = HW/P^2$ patches
        \item Per-patch linear projection $\to X \in \mathbb{R}^{S \times C}$
        \item $N$ Mixer layers (token-mix + channel-mix)
        \item Global average pooling over $S$ $\to \mathbb{R}^C$
        \item Linear classification head
      \end{enumerate}

      \vspace{0.3em}
      \textbf{Why it works}:
      \begin{itemize}
        \item Token-mixing MLPs learn \emph{fixed} spatial mixing patterns (no data-dependent attention)
        \item With enough data, fixed mixing is competitive with attention
        \item Much simpler architecture
      \end{itemize}

      \vspace{0.3em}
      \textbf{Training} (no locality prior $\Rightarrow$ heavy regularization on IN-1K): see next slide for the full recipe.
    \end{column}
    \begin{column}{0.48\textwidth}
      \textbf{Comparison with Transformer}:
      \begin{center}
        \small
        \begin{tabular}{lcc}
          \toprule
          & \textbf{Transformer} & \textbf{Mixer} \\
          \midrule
          Spatial mix & Attention & Token MLP \\
          Channel mix & FFN & Channel MLP \\
          Data-dep.\ & Yes & No \\
          Complexity & $\mathcal{O}(S^2)$ & $\mathcal{O}(S)$ \\
          \bottomrule
        \end{tabular}
      \end{center}

      \vspace{0.4em}
      \textbf{ImageNet-1K top-1}:
      \begin{center}
        \small
        \begin{tabular}{lccc}
          \toprule
          \textbf{Variant} & \textbf{Params} & \textbf{IN-1K} & \textbf{$+$21K pre} \\
          \midrule
          Mixer-B/16 & 59M  & 76.4 & 80.6 \\
          Mixer-L/16 & 207M & 71.8 & 82.2 \\
          \midrule
          \textit{ViT-B/16} & 86M & 77.9 & 84.0 \\
          \bottomrule
        \end{tabular}
      \end{center}

      {\footnotesize From scratch on IN-1K, Mixer \textbf{trails ViT} (76.4 vs 77.9) and the bigger \textbf{L/16 overfits} ($71.8 < 76.4$). With ImageNet-21K pre-training the ordering flips ($L{>}B$): fixed token-mixing is \textbf{data-hungry} but scales well with data.}
    \end{column}
  \end{columns}
\end{frame}

% -----------------------------------------------------------------------
\begin{frame}{MLP-Mixer: Training Recipe}

  No convolution/attention $\Rightarrow$ \textbf{no locality or translation-equivariance prior}. Mixer overfits badly on IN-1K and leans on \textbf{strong regularization}; it only shines with large-scale pre-training.

  \vspace{0.3em}
  \begin{columns}[T]
    \begin{column}{0.5\textwidth}
      \textbf{Optimization} (same for pre-train \& scratch)
      \begin{description}[style=nextline, leftmargin=0pt, itemsep=1pt]
        \item[Optimizer] Adam, $\beta_1{=}0.9$, $\beta_2{=}0.999$
        \item[Peak LR] $10^{-3}$, $10$k-step linear warm-up
        \item[Schedule] linear decay
        \item[Batch size] $4096$
        \item[Weight decay] $0.1$ (IN-1K / IN-21K)
        \item[Grad clip] global norm $1$
        \item[Epochs] $300$ (IN-1K \& IN-21K); JFT-300M only $7$--$14$
      \end{description}
    \end{column}
    \begin{column}{0.48\textwidth}
      \textbf{Regularization / augmentation} (IN-1K from scratch)
      \begin{description}[style=nextline, leftmargin=0pt, itemsep=1pt]
        \item[RandAugment] magnitude $15$
        \item[Mixup] $\alpha = 0.5$
        \item[Stochastic depth] $0.1$
        \item[Dropout] $0.0$
      \end{description}

      \vspace{0.3em}
      \textbf{Fine-tuning} (after pre-training)
      \begin{description}[style=nextline, leftmargin=0pt, itemsep=1pt]
        \item[Optimizer] SGD + momentum, batch $512$
        \item[Schedule] cosine + linear warm-up, grad clip $1$
        \item[Resolution] $224$ or $448$ (higher helps)
      \end{description}
    \end{column}
  \end{columns}

  \vspace{0.4em}
  {\footnotesize Large-scale pre-training (IN-21K, JFT-300M) needs \emph{less} augmentation; the data itself regularizes. This is what flips Mixer from below ViT to competitive. \cite{tolstikhin2021mlpmixer}}
\end{frame}

% =======================================================================
\subsection{DETR}

% -----------------------------------------------------------------------
\begin{frame}{The Object Detection Problem}

  \textbf{Task}: localize \emph{and} classify every object instance in an image. Output a \textbf{variable-size set} $\{(c_i, b_i)\}_{i=1}^{M}$ of (class, bounding box) pairs.

  \vspace{0.3em}
  \begin{columns}[T]
    \begin{column}{0.42\textwidth}
      \begin{center}
        \begin{tikzpicture}[scale=1.5, every node/.style={scale=1.0}]
          \begin{scope}
            \clip (0,0) rectangle (2,2);
            \visionSketchScene
          \end{scope}
          \draw[thick] (0,0) rectangle (2,2);
          % Detection boxes with labels
          \draw[very thick, green!50!black] (0.03, 0.49) rectangle (0.68, 1.58);
          \node[fill=green!50!black, text=white, font=\tiny, inner sep=1pt, anchor=south west]
            at (0.03, 1.58) {tree};
          \draw[very thick, red!70!black] (0.76, 0.49) rectangle (1.54, 1.52);
          \node[fill=red!70!black, text=white, font=\tiny, inner sep=1pt, anchor=south west]
            at (0.76, 1.52) {house};
          \draw[very thick, violet] (1.68, 0.43) rectangle (1.9, 1.02);
          \node[fill=violet, text=white, font=\tiny, inner sep=1pt, anchor=south east]
            at (1.9, 1.02) {person};
        \end{tikzpicture}
      \end{center}
      \vspace{0.2em}
      {\footnotesize Box $b_i=(c_x, c_y, w, h)$, normalized to $[0,1]$.}
    \end{column}
    \begin{column}{0.56\textwidth}
      \textbf{Harder than classification}
      \begin{itemize}
        \item Classification outputs \emph{one} label; detection outputs an \emph{unknown number} of objects.
        \item Output is a \textbf{set}: no canonical order, no fixed size.
        \item Must not report the same object twice (\textbf{duplicate suppression}).
      \end{itemize}

      \vspace{0.3em}
      \textbf{Classic pipelines} \;{\footnotesize\cite{ren2015fasterrcnn,redmon2016yolo}}
      \begin{itemize}
        \item Tile the image with dense \textbf{anchors} / region proposals.
        \item Per anchor: classify + regress a box offset.
        \item \textbf{Non-Maximum Suppression} (NMS) prunes overlapping boxes.
      \end{itemize}
      \vspace{0.2em}
      {\footnotesize Relies on hand-crafted components: anchor priors, assignment heuristics, NMS. Not end-to-end.}
    \end{column}
  \end{columns}

  \vspace{0.4em}
  \begin{center}
    \footnotesize\itshape DETR's idea: cast detection as \textbf{direct set prediction} and learn it end-to-end, dropping anchors and NMS.
  \end{center}
\end{frame}

% -----------------------------------------------------------------------
\begin{frame}[fragile]{DETR: End-to-End Object Detection with Transformers}

  \cite{carion2020detr}: Reformulate object detection as a \textbf{direct set prediction} problem: no anchors, no NMS, no hand-crafted components.

  \vspace{0.2em}
  \begin{center}
    \begin{tikzpicture}[
      scale=0.66, every node/.style={scale=0.66},
      box/.style={draw, rounded corners, fill=#1, minimum width=1.8cm, minimum height=0.5cm, align=center, font=\small},
      qbox/.style={draw=#1, fill=#1, fill opacity=0.20, text opacity=1, thick,
                   minimum width=0.46cm, minimum height=0.34cm, inner sep=1pt, font=\tiny},
      arr/.style={-Stealth, thick},
    ]
      % Input image (actual scene)
      \node[inner sep=0] (img) at (-5.6, 0) {
        \begin{tikzpicture}[scale=0.7]
          \begin{scope}\clip (0,0) rectangle (2,2);\visionSketchScene\end{scope}
          \draw[thick] (0,0) rectangle (2,2);
        \end{tikzpicture}
      };
      \node[font=\tiny, below=0.04cm of img] {Image $H{\times}W{\times}3$};

      % CNN Backbone (trapezoid funnel)
      \draw[thick, fill=orange!15] (-3.7, 0.85) -- (-2.1, 0.5) -- (-2.1, -0.5) -- (-3.7, -0.85) -- cycle;
      \node[font=\small] at (-2.9, 0) {CNN};
      \node[font=\tiny, text=gray] at (-2.9, -1.15) {ResNet-50};
      \draw[arr] (img.east) -- (-3.7, 0);

      % 1x1 conv
      \node[box=orange!8, minimum width=0.9cm, font=\tiny] (conv) at (-0.9, 0) {$1{\times}1$\\Conv};
      \draw[arr] (-2.1, 0) -- (conv);

      % Positional encoding
      \node[font=\tiny, text=orange!70!black] at (-0.9, 0.8) {$+ \text{pos}_{2D}$};

      % Transformer Encoder
      \node[box=red!12, minimum width=1.8cm, minimum height=1.4cm] (enc) at (1.4, 0) {};
      \node[font=\small, text=red!60!black] at (1.4, 0.25) {\textbf{Encoder}};
      \node[font=\tiny, text=gray] at (1.4, -0.2) {$6$ layers};
      \draw[arr] (conv) -- (enc);

      % Transformer Decoder
      \node[box=blue!12, minimum width=1.8cm, minimum height=1.4cm] (dec) at (4.2, 0) {};
      \node[font=\small, text=blue!60!black] at (4.2, 0.25) {\textbf{Decoder}};
      \node[font=\tiny, text=gray] at (4.2, -0.2) {$6$ layers};
      \draw[arr] (enc) -- node[above, font=\tiny] {memory} (dec);

      % --- Input object queries (5), centered below decoder; colored per slot
      %     so each matches its output embedding o_i (same color, same size) ---
      \foreach \i/\c in {0/blue!70, 1/orange!90!black, 2/teal, 3/magenta!80, 4/brown!70!black} {
        \node[qbox=\c] at (3.16 + \i*0.52, -1.83) {};
      }
      \node[font=\tiny, text=gray!40!black] at (4.2, -2.28) {object queries (in)};
      \draw[arr, gray!55!black] (4.2, -1.5) -- (dec.south);

      % --- Output queries: one decoder output embedding per slot, colored to
      %     match the corresponding input query ---
      \node[font=\tiny, text=gray!40!black] at (5.6, 1.62) {output queries};
      \foreach \i/\yy/\c in {1/1.2/blue!70, 2/0.6/orange!90!black, 3/0.0/teal, 4/-0.6/magenta!80, 5/-1.2/brown!70!black} {
        \node[qbox=\c] (o\i) at (5.6, \yy) {$\mathbf{o}_{\i}$};
        \draw[arr, \c, thin] (dec.east) -- (o\i.west);
      }

      % --- Shared heads, applied to every output query ---
      \node[draw, rounded corners, fill=gray!10, minimum width=0.95cm, minimum height=2.85cm,
            align=center, font=\tiny] (heads) at (6.95, 0) {Class\\FFN\\[3pt]Box\\MLP};
      \node[font=\tiny, text=gray] at (6.95, -1.66) {shared};
      \foreach \i in {1,...,5} {\draw[arr, thin] (o\i.east) -- (heads.west |- o\i);}

      % --- Result image (assembled predicted boxes) ---
      \node[inner sep=0] (result) at (9.3, 0) {
        \begin{tikzpicture}[scale=0.62]
          \begin{scope}\clip (0,0) rectangle (2,2);\visionSketchScene\end{scope}
          \draw[thick] (0,0) rectangle (2,2);
          \draw[very thick, green!50!black] (0.03, 0.49) rectangle (0.68, 1.58);
          \draw[very thick, red!70!black]   (0.76, 0.49) rectangle (1.54, 1.52);
          \draw[very thick, violet]         (1.68, 0.43) rectangle (1.9, 1.02);
        \end{tikzpicture}
      };
      \node[font=\tiny, below=0.02cm of result] {predictions};
      \node[font=\tiny] at (7.95, 1.62) {$(c_i,b_i)$ / $\varnothing$};

      % --- Per-query head output: a (class, box), or no-object ---
      % o1 -> tree (green)
      \draw[arr, thin, green!55!black] (heads.east |- o1) -- (7.62,1.2);
      \draw[green!55!black, very thick, fill=green!15] (7.64,1.06) rectangle (7.98,1.34);
      \draw[arr, thin, green!55!black] (7.98,1.2) -- ([yshift=8pt]result.west);
      % o2 -> no object
      \draw[arr, thin, gray] (heads.east |- o2) -- (7.62,0.6);
      \node[font=\tiny, text=gray, anchor=west] at (7.66,0.6) {$\varnothing$};
      % o3 -> house (red)
      \draw[arr, thin, red!70!black] (heads.east |- o3) -- (7.62,0.0);
      \draw[red!70!black, very thick, fill=red!15] (7.64,-0.14) rectangle (7.98,0.14);
      \draw[arr, thin, red!70!black] (7.98,0.0) -- (result.west);
      % o4 -> person (violet)
      \draw[arr, thin, violet] (heads.east |- o4) -- (7.62,-0.6);
      \draw[violet, very thick, fill=violet!18] (7.64,-0.74) rectangle (7.98,-0.46);
      \draw[arr, thin, violet] (7.98,-0.6) -- ([yshift=-8pt]result.west);
      % o5 -> no object
      \draw[arr, thin, gray] (heads.east |- o5) -- (7.62,-1.2);
      \node[font=\tiny, text=gray, anchor=west] at (7.66,-1.2) {$\varnothing$};
    \end{tikzpicture}
  \end{center}

  \vspace{0.3em}
  \begin{columns}
    \begin{column}{0.48\textwidth}
      \textbf{Key components}:
      \begin{itemize}
        \item CNN extracts features $\to$ flatten to sequence
        \item Encoder: self-attention over spatial features
        \item Decoder: $N$ \textbf{object queries} cross-attend to encoder memory
        \item Each query predicts one (class, box) or ``no object''
      \end{itemize}
    \end{column}
    \begin{column}{0.48\textwidth}
      \textbf{Why it matters}:
      \begin{itemize}
        \item No anchors, NMS, or hand-crafted post-processing
        \item Set prediction via \textbf{Hungarian matching}
        \item Object queries self-attend $\to$ reason about inter-object relations
        \item Simple, elegant, end-to-end trainable
      \end{itemize}
    \end{column}
  \end{columns}
\end{frame}

% -----------------------------------------------------------------------
\begin{frame}[fragile]{DETR: Decoder and Object Queries}

  The decoder turns $N$ \textbf{object queries} into $N$ output embeddings \textbf{in parallel} (non-autoregressive). Each of the $6$ identical layers applies self-attention, cross-attention, then an FFN.

  \vspace{0.2em}
  \begin{columns}[T]
    \begin{column}{0.46\textwidth}
      \begin{center}
        \begin{tikzpicture}[
          scale=0.82, every node/.style={scale=0.82},
          box/.style={draw, rounded corners, fill=#1, minimum width=2.6cm,
                      minimum height=0.5cm, align=center, font=\scriptsize},
          tok/.style={draw=#1, fill=#1, fill opacity=0.22, text opacity=1, thick,
                      minimum width=0.34cm, minimum height=0.26cm, inner sep=0pt},
          arr/.style={-Stealth, thick},
        ]
          % Inputs: object queries as a row of tokens
          \foreach \i in {0,...,3} {\node[tok=violet] at (-0.6 + \i*0.42, -0.55) {};}
          \coordinate (q) at (0.03, -0.55);
          \node[font=\tiny, text=violet] at (0.03, -1.0) {object queries (tokens)};

          % Self-attention
          \node[box=violet!18] (sa) at (0, 0.5) {Self-Attention};
          \draw[arr] (q) -- (sa);
          \node[draw, circle, inner sep=1pt, font=\tiny, fill=white] (a1) at (0, 1.15) {$+$};
          \draw[arr] (sa) -- (a1);

          % Cross-attention
          \node[box=blue!15] (ca) at (0, 1.9) {Cross-Attention};
          \draw[arr] (a1) -- (ca);
          \node[draw, circle, inner sep=1pt, font=\tiny, fill=white] (a2) at (0, 2.55) {$+$};
          \draw[arr] (ca) -- (a2);

          % Encoder memory: spatial feature-map grid flattened into K,V tokens
          \begin{scope}[shift={(-3.45,1.9)}]
            \draw[thick, fill=red!8] (-0.27,-0.27) rectangle (0.27,0.27);
            \foreach \g in {-0.09,0.09} {
              \draw[gray!60] (\g,-0.27) -- (\g,0.27);
              \draw[gray!60] (-0.27,\g) -- (0.27,\g);
            }
          \end{scope}
          \node[font=\tiny, text=gray, align=center] at (-3.45, 1.3) {encoder\\feature map};
          \draw[arr, red!60!black] (-3.1,1.9) -- (-2.75,1.9);
          \foreach \j in {0,...,3} {\node[tok=red!70!black] at (-2.45, 2.34 - \j*0.30) {};}
          \draw[arr, red!60!black] (-2.28,1.9) -- node[above, font=\tiny] {$K,V$} (ca.west);

          % FFN
          \node[box=green!12] (ffn) at (0, 3.3) {FFN};
          \draw[arr] (a2) -- (ffn);
          \node[draw, circle, inner sep=1pt, font=\tiny, fill=white] (a3) at (0, 3.95) {$+$};
          \draw[arr] (ffn) -- (a3);

          % Output
          \node[font=\scriptsize] (out) at (0, 4.6) {output embeddings};
          \draw[arr] (a3) -- (out);

          % query Q label into self/cross attn
          \node[font=\tiny, text=violet, anchor=west] at (1.5, 0.5) {$Q,K,V$};
          \node[font=\tiny, text=violet, anchor=west] at (1.5, 1.9) {$Q$};

          % residual hints
          \draw[arr, gray!50] (0,-0.1) -- (1.85,-0.1) -- (1.85,1.15) -- (a1.east);
          \draw[arr, gray!50] (a1.east) ++(0,0) -- (1.85,1.15) -- (1.85,2.55) -- (a2.east);
          \draw[arr, gray!50] (a2.east) ++(0,0) -- (1.85,2.55) -- (1.85,3.95) -- (a3.east);
        \end{tikzpicture}
      \end{center}
      \vspace{0.1em}
      {\footnotesize $\times 6$ layers; LayerNorm omitted. Heads applied to each output embedding.}
    \end{column}
    \begin{column}{0.52\textwidth}
      \textbf{Object queries} -- the key idea
      \begin{itemize}
        \item $N{=}100$ \textbf{learned}, input-independent embeddings: a fixed set of ``slots''.
        \item Act as \textbf{learned positional encodings}, added to $Q$ (and $K$) at \emph{every} layer; each specializes to a region/size band.
      \end{itemize}

      \vspace{0.2em}
      \textbf{Cross-attention} (query $\to$ image)
      \begin{itemize}
        \item $Q$ from queries; $K,V$ from \textbf{encoder memory}, with fixed \textbf{2D sinusoidal} pos.\ enc.\ on $K$ so queries attend by \emph{location}.
        \item Each query pools its object's evidence into one vector.
      \end{itemize}

      \vspace{0.2em}
      \textbf{Self-attention} (query $\to$ query)
      \begin{itemize}
        \item Queries communicate $\to$ avoid two claiming the \emph{same} object (implicit NMS).
      \end{itemize}
    \end{column}
  \end{columns}

  \vspace{0.2em}
  {\footnotesize Shared heads per output embedding: class FFN $\to$ softmax over classes${+}\varnothing$; box MLP $\to$ sigmoid $(c_x,c_y,w,h)$. Auxiliary loss after \emph{every} decoder layer.}
\end{frame}

% -----------------------------------------------------------------------
\begin{frame}{DETR: Set Loss via Hungarian Matching: Idea}

  \begin{columns}[T]
    \begin{column}{0.54\textwidth}
      \textbf{The problem}
      \begin{itemize}
        \item Every prediction needs a training \emph{target}.
        \item But the $N$ predictions are an \textbf{unordered set}: no fixed object$\to$slot correspondence.
        \item Any hand-chosen assignment is \textbf{arbitrary}.
        \item Uncoordinated, several predictions \textbf{chase the same object} $\Rightarrow$ duplicates.
      \end{itemize}

      \vspace{0.4em}
      \textbf{The idea: let the assignment be optimal}
      \begin{itemize}
        \item Pad the $M$ ground-truth objects to $N$ with $\varnothing$ (``no object'').
        \item \textbf{Cost} every (prediction, object) pair: how well that prediction fits the object (class $+$ box).
        \item Pick the \textbf{one-to-one matching of lowest total cost} (a bipartite assignment).
        \item Train the usual loss on \textbf{matched pairs}; unmatched predictions $\to \varnothing$.
      \end{itemize}

      \vspace{0.4em}
      \textbf{Why it works}
      \begin{itemize}
        \item One-to-one $\Rightarrow$ \textbf{one detector per object}: no duplicates, \textbf{no NMS}.
      \end{itemize}
    \end{column}
    \begin{column}{0.44\textwidth}
      \begin{center}
        \begin{tikzpicture}[scale=0.6, font=\scriptsize,
            arr/.style={-Stealth, semithick},
            qbox/.style={fill opacity=0.18, text opacity=1, thick,
                         rounded corners=1pt, minimum width=0.7cm, minimum height=0.36cm, font=\scriptsize}]
          % --- object queries (predictions) on the left ---
          \node[font=\scriptsize, text=violet, align=center] at (0.3, 3.95) {object queries\\(predictions)};
          \foreach \i/\yy/\c in {1/3.2/green!55!black, 2/2.4/red!70!black, 3/1.6/violet, 4/0.8/gray} {
            \node[qbox, draw=\c, fill=\c] (q\i) at (0.3,\yy) {$\hat y_\i$};
          }
          % --- image with ground-truth objects ---
          \begin{scope}[shift={(2.7,0.55)}, scale=0.95]
            \begin{scope}\clip (0,0) rectangle (2,2);\visionSketchScene\end{scope}
            \draw[thick] (0,0) rectangle (2,2);
            \draw[very thick, green!55!black] (0.04,0.49) rectangle (0.66,1.56);
            \draw[very thick, red!70!black]   (0.76,0.49) rectangle (1.54,1.52);
            \draw[very thick, violet]         (1.66,0.42) rectangle (1.92,1.02);
            \coordinate (tree)  at (0.66,1.30);
            \coordinate (house) at (1.54,1.25);
            \coordinate (pers)  at (1.92,0.85);
          \end{scope}
          \node[font=\scriptsize, text=gray!50!black] at (3.65,2.75) {image $+$ ground truth};
          % --- no-object slot ---
          \node[draw=gray, fill=gray!14, thick, rounded corners=1pt,
                minimum width=0.6cm, minimum height=0.36cm, font=\scriptsize] (none) at (3.0,-0.55) {$\varnothing$};
          % --- one-to-one matching edges, colored by target ---
          \draw[arr, green!55!black] (q1.east) to[out=0,in=180] (tree);
          \draw[arr, red!70!black]   (q2.east) to[out=0,in=180] (house);
          \draw[arr, violet]         (q3.east) to[out=0,in=185] (pers);
          \draw[arr, gray]           (q4.east) to[out=0,in=180] (none);
        \end{tikzpicture}
      \end{center}
      \vspace{0.2em}
      {\footnotesize Each \textbf{object query} yields one prediction $\hat y_i$. Hungarian matching assigns each prediction to \textbf{one} ground-truth object (or $\varnothing$): a one-to-one prediction$\leftrightarrow$object pairing of minimum total cost.}
    \end{column}
  \end{columns}
\end{frame}

% -----------------------------------------------------------------------
\begin{frame}{DETR: Set Loss via Hungarian Matching: Details}

  \textbf{Ingredients.} Ground truth $y_i = (c_i, b_i)$ (class, box); prediction $\hat y_j = (\hat p_j, \hat b_j)$.
  \begin{itemize}
    \item $\hat p_j(c)$: prediction $j$'s \textbf{softmax probability} of class $c$ (over the $K$ classes $+\,\varnothing$). So $\hat p_{\sigma(i)}(c_i)$ is the mass it puts on object $i$'s \emph{true} class.
    \item $\mathcal{L}_\text{box}$: the \textbf{box cost}, scale-invariant GIoU $+$ $L_1$ on the $(c_x,c_y,w,h)$ boxes:
      \[
        \mathcal{L}_\text{box}(b_i, \hat{b}) = \lambda_\text{giou}\,\mathcal{L}_\text{giou}(b_i, \hat{b}) + \lambda_{L_1}\,\| b_i - \hat{b} \|_1
      \]
  \end{itemize}

  \vspace{0.15em}
  \textbf{Step 1: Optimal assignment.} Over all permutations $\sigma \in \mathfrak{S}_N$ (predictions $\leftrightarrow$ padded targets), minimize the total \emph{matching cost}:
  \[
    \hat{\sigma} = \arg\min_{\sigma \in \mathfrak{S}_N} \sum_{i=1}^{N} \mathcal{L}_\text{match}(y_i, \hat{y}_{\sigma(i)}),
    \qquad
    \mathcal{L}_\text{match} = \mathbbm{1}_{c_i \neq \varnothing}\bigl[-\hat{p}_{\sigma(i)}(c_i) + \mathcal{L}_\text{box}(b_i, \hat{b}_{\sigma(i)})\bigr]
  \]
  \begin{itemize}
    \item Uses the class \emph{probability} (not $\log$), so it sits on the same scale as the box cost.
    \item Solved by the \textbf{Hungarian algorithm} in $\mathcal{O}(N^3)$; \textbf{non-differentiable} (run under \texttt{no\_grad}).
  \end{itemize}

  \vspace{0.15em}
  \textbf{Step 2: Hungarian loss.} With $\hat{\sigma}$ fixed, the trained loss on the matched pairs:
  \[
    \mathcal{L} = \sum_{i=1}^{N} \bigl[-\log \hat{p}_{\hat{\sigma}(i)}(c_i) + \mathbbm{1}_{c_i \neq \varnothing}\, \mathcal{L}_\text{box}(b_i, \hat{b}_{\hat{\sigma}(i)})\bigr]
  \]
  \begin{itemize}
    \item Class term: now \textbf{cross-entropy} ($-\log\hat p$, for good gradients); $\varnothing$ down-weighted $10\times$ (class imbalance).
  \end{itemize}
\end{frame}

% -----------------------------------------------------------------------
\begin{frame}{DETR: Generalized IoU (GIoU)}

  \cite{rezatofighi2019giou}: a \textbf{scale-invariant} box overlap measure that, unlike plain IoU, still gives a useful signal when boxes \emph{do not overlap}.

  \vspace{0.3em}
  \begin{center}
    \begin{tikzpicture}[
        scale=0.62, every node/.style={scale=0.62},
        font=\small,
        gt/.style={green!55!black, very thick},
        pr/.style={blue!70!black, very thick},
        enc/.style={densely dashed, very thick, black!75},
        capt/.style={font=\scriptsize, text=gray},
      ]
      % ===== Panel 1: overlapping =====
      \begin{scope}
        \fill[gray!22] (0,0) rectangle (3,2.2);                 % enclosing C
        \fill[green!30] (0,0) rectangle (2,1.5);                % A (GT)
        \fill[blue!30]  (1,0.6) rectangle (3,2.2);              % B (pred)
        \fill[orange!55] (1,0.6) rectangle (2,1.5);             % A cap B
        \draw[gt] (0,0) rectangle (2,1.5);
        \draw[pr] (1,0.6) rectangle (3,2.2);
        \draw[enc] (0,0) rectangle (3,2.2);
        \node[font=\scriptsize, text=green!45!black] at (0.5,0.28) {$A$ (GT)};
        \node[font=\scriptsize, text=blue!70!black] at (2.5,1.95) {$B$ (pred)};
        \node[font=\scriptsize] at (1.5,1.05) {$A\!\cap\!B$};
        \node[capt] at (1.5,2.55) {smallest enclosing box $C$};
        % penalty region callout
        \draw[-Stealth, gray] (0.45,3.1) to[out=-90,in=120] (0.5,1.85);
        \node[capt, align=center] at (0.45,3.35) {$C\setminus(A\cup B)$\\(wasted area)};
        \node[font=\scriptsize\bfseries] at (1.5,-0.55) {overlap: $\mathrm{IoU}>0$};
      \end{scope}

      % ===== Panel 2: disjoint =====
      \begin{scope}[shift={(6.2,0)}]
        \fill[gray!22] (0,0) rectangle (3.4,2.4);               % enclosing C
        \fill[green!30] (0,0) rectangle (1.3,1.1);              % A
        \fill[blue!30]  (2.3,1.3) rectangle (3.4,2.4);          % B
        \draw[gt] (0,0) rectangle (1.3,1.1);
        \draw[pr] (2.3,1.3) rectangle (3.4,2.4);
        \draw[enc] (0,0) rectangle (3.4,2.4);
        \node[font=\scriptsize, text=green!45!black] at (0.65,0.55) {$A$};
        \node[font=\scriptsize, text=blue!70!black] at (2.85,1.85) {$B$};
        \node[capt] at (1.7,2.75) {$C$ still encloses both};
        \node[font=\scriptsize\bfseries, align=center] at (1.7,-0.55)
          {disjoint: $\mathrm{IoU}=0$,\; but $\mathrm{GIoU}<0$};
      \end{scope}
    \end{tikzpicture}
  \end{center}

  \vspace{0.2em}
  \begin{columns}[T]
    \begin{column}{0.46\textwidth}
      \begin{align*}
        \mathrm{IoU}  &= \frac{|A\cap B|}{|A\cup B|} \\[2pt]
        \mathrm{GIoU} &= \mathrm{IoU} - \frac{|C\setminus(A\cup B)|}{|C|} \\[2pt]
        \mathcal{L}_\text{giou} &= 1 - \mathrm{GIoU} \;\in\; [0,2)
      \end{align*}
    \end{column}
    \begin{column}{0.54\textwidth}
      \begin{itemize}\itemsep2pt
        \item $C$: smallest axis-aligned box enclosing $A$ and $B$.
        \item $\mathrm{GIoU}\in(-1,1]$; $=1$ iff $A{=}B$, and $\mathrm{GIoU}\le\mathrm{IoU}$ always.
        \item \textbf{Disjoint boxes}: plain IoU $=0$ everywhere $\Rightarrow$ no gradient; the enclosing term keeps shrinking as boxes approach $\Rightarrow$ a usable pull.
        \item \textbf{Scale-invariant} (ratios of areas) $\Rightarrow$ consistent across object sizes, unlike $L_1$ which grows with box size.
      \end{itemize}
    \end{column}
  \end{columns}
\end{frame}

% -----------------------------------------------------------------------
\begin{frame}{DETR: Empirical Results}

  \cite{carion2020detr} on \textbf{COCO 2017} detection (val), ResNet backbones.

  \vspace{0.3em}
  \begin{columns}[T]
    \begin{column}{0.55\textwidth}
      \begin{center}
        \small
        \begin{tabular}{lccccc}
          \toprule
          \textbf{Model} & \textbf{AP} & \textbf{AP$_S$} & \textbf{AP$_M$} & \textbf{AP$_L$} & \textbf{Ep.} \\
          \midrule
          Faster R-CNN-FPN+ & 42.0 & \textbf{26.6} & 45.4 & 53.4 & 109 \\
          DETR-R50          & 42.0 & 20.5 & 45.8 & 61.1 & 500 \\
          DETR-DC5-R101     & \textbf{44.9} & 23.7 & 49.5 & \textbf{62.3} & 500 \\
          \bottomrule
        \end{tabular}
      \end{center}
      \vspace{0.2em}
      {\footnotesize Table 1 of \cite{carion2020detr}. AP$_S$/AP$_M$/AP$_L$: small / medium / large objects. ``$+$'': Faster R-CNN with GIoU, crop aug.\ and the long $9\times$ ($\sim$109-epoch) schedule, for a fair match to DETR's 500.}
    \end{column}
    \begin{column}{0.43\textwidth}
      \begin{itemize}
        \item \textbf{Matches} a strong Faster R-CNN in overall AP, with no anchors and no NMS.
        \item \textbf{Best on large objects} (AP$_L$ $+8$): global self-attention sees the whole image at once.
        \item \textbf{Worse on small objects} (AP$_S$ $-6$): single low-res feature map, no multi-scale.
        \item \textbf{Slow to train}: $500$ epochs $\Rightarrow$ motivates Deformable DETR ($50$ ep).
      \end{itemize}
    \end{column}
  \end{columns}

  \vspace{0.2em}
  {\footnotesize Adding a mask head extends the same set-prediction model to \textbf{panoptic segmentation} competitively. Takeaway: a simple, end-to-end detector that is on par with mature pipelines, trading small-object accuracy and training time for architectural simplicity.}
\end{frame}

% -----------------------------------------------------------------------
\begin{frame}{DETR: Training Details and Practicalities}

  \begin{columns}[T]
    \begin{column}{0.5\textwidth}
      \textbf{Data \& optimization}
      \begin{itemize}
        \item \textbf{COCO 2017} ($\sim$118k images). Backbone: ImageNet-pretrained ResNet-50/101 with \textbf{frozen BatchNorm}.
        \item \textbf{AdamW}: lr $10^{-4}$ (transformer), $10^{-5}$ (backbone), weight decay $10^{-4}$; dropout $0.1$; $N{=}100$ queries.
        \item \textbf{300 ep} (lr ${\div}10$ at 200) or \textbf{500 ep} (at 400) for best AP.
        \item Augmentation: \textbf{scale aug} (shortest side $480$--$800$, longest $\le 1333$) $+$ random crop.
      \end{itemize}

      \vspace{0.2em}
      \textbf{Loss weights}
      \begin{itemize}
        \item Box: $\lambda_{L_1}{=}5$, $\lambda_\text{giou}{=}2$.
        \item \textbf{Down-weight $\varnothing$ by $10\times$}: most of the 100 queries match no object, else the ``no-object'' term dominates.
      \end{itemize}
    \end{column}
    \begin{column}{0.48\textwidth}
      \textbf{Auxiliary (deep-supervision) losses}
      \begin{itemize}
        \item Prediction heads $+$ Hungarian loss after \emph{every} one of the 6 decoder layers, not just the last.
        \item Heads \textbf{share parameters} across layers; a shared LayerNorm normalizes their inputs. Speeds convergence, helps output the right object \emph{count}.
      \end{itemize}

      \vspace{0.2em}
      \textbf{Inference / post-processing}
      \begin{itemize}
        \item \textbf{No NMS, no anchors}: one forward pass $\to$ 100 (class, box); keep $\arg\max$ class $\neq\varnothing$ per query, ranked by score.
      \end{itemize}

      \vspace{0.2em}
      \textbf{What else matters}
      \begin{itemize}
        \item The very \textbf{long schedule} and $\varnothing$ down-weighting are essential; GIoU adds scale-invariance.
      \end{itemize}
    \end{column}
  \end{columns}
\end{frame}

% =======================================================================
\subsection{Deformable DETR}

% -----------------------------------------------------------------------
\begin{frame}{Deformable DETR: Sparse Attention for Detection}

  \cite{zhu2021deformabledetr}: Fix DETR's slow convergence (500 epochs) and poor small-object performance via \textbf{deformable attention}. Each query attends to only $K$ learned sampling points.

  \vspace{0.3em}
  \begin{center}
    \begin{tikzpicture}[
      scale=0.63, every node/.style={scale=0.63},
      arr/.style={-Stealth, thick},
    ]
      % --- Standard Attention (left) ---
      \begin{scope}[shift={(-4.5, 0)}]
        \node[font=\small\bfseries] at (1.5, 3.5) {DETR: Global Attention};
        \draw[thick, fill=gray!8] (0, 0) rectangle (3, 3);
        % Grid
        \foreach \i in {0.5, 1.0, ..., 2.5} {
          \draw[gray!20] (\i, 0) -- (\i, 3);
          \draw[gray!20] (0, \i) -- (3, \i);
        }
        % Reference point
        \fill[red] (1.5, 1.5) circle (3pt);
        % Attention to all positions
        \foreach \x in {0.25, 0.75, ..., 2.75} {
          \foreach \y in {0.25, 0.75, ..., 2.75} {
            \draw[red!30, thin, -Stealth] (1.5, 1.5) -- (\x, \y);
          }
        }
        \fill[red] (1.5, 1.5) circle (3pt);
        \node[font=\scriptsize, text=red!60!black] at (1.5, -0.4) {$\mathcal{O}(H^2 W^2)$};
      \end{scope}

      % --- Deformable Attention (right) ---
      \begin{scope}[shift={(2.0, 0)}]
        \node[font=\small\bfseries] at (1.5, 3.5) {Deformable: $K$ Points};
        \draw[thick, fill=gray!8] (0, 0) rectangle (3, 3);
        \foreach \i in {0.5, 1.0, ..., 2.5} {
          \draw[gray!20] (\i, 0) -- (\i, 3);
          \draw[gray!20] (0, \i) -- (3, \i);
        }
        % Reference point
        \fill[red] (1.5, 1.5) circle (3pt);
        \node[font=\tiny, text=red!70!black, below right=1pt] at (1.5, 1.5) {$\mathbf{p}_q$};
        % K=4 sampling points with offset arrows
        \draw[blue!70!black, thick, -Stealth] (1.5, 1.5) -- (0.8, 2.3) node[circle, fill=blue!60, inner sep=2pt] {};
        \draw[blue!70!black, thick, -Stealth] (1.5, 1.5) -- (2.4, 2.1) node[circle, fill=blue!60, inner sep=2pt] {};
        \draw[blue!70!black, thick, -Stealth] (1.5, 1.5) -- (2.1, 0.6) node[circle, fill=blue!60, inner sep=2pt] {};
        \draw[blue!70!black, thick, -Stealth] (1.5, 1.5) -- (0.5, 0.9) node[circle, fill=blue!60, inner sep=2pt] {};
        % Labels
        \node[font=\tiny, text=blue!70!black] at (0.5, 2.5) {$\Delta\mathbf{p}_1$};
        \node[font=\tiny, text=blue!70!black] at (2.7, 2.3) {$\Delta\mathbf{p}_2$};
        \node[font=\tiny, text=blue!70!black] at (2.5, 0.4) {$\Delta\mathbf{p}_3$};
        \node[font=\tiny, text=blue!70!black] at (0.2, 0.7) {$\Delta\mathbf{p}_4$};
        \node[font=\scriptsize, text=blue!60!black] at (1.5, -0.4) {$\mathcal{O}(HW \cdot K)$, $K{=}4$};
      \end{scope}
    \end{tikzpicture}
  \end{center}

  \vspace{0.2em}
  \textbf{Deformable attention}:
  \[
    \operatorname{DeformAttn}(\mathbf{z}_q, \mathbf{p}_q, \mathbf{x}) = \sum_{m=1}^{M} W_m \biggl[\sum_{k=1}^{K} A_{mqk} \cdot W'_m \, \mathbf{x}(\mathbf{p}_q + \Delta\mathbf{p}_{mqk})\biggr]
  \]
  \begin{itemize}
    \item $\mathbf{z}_q$: \textbf{query feature} (content embedding of query $q$); \;$\mathbf{p}_q$: its 2D \textbf{reference point} (where the query ``looks'').
    \item $\mathbf{x}$: input feature map; \;$\mathbf{x}(\mathbf{p}_q{+}\Delta\mathbf{p})$ read by \textbf{bilinear interpolation} (continuous positions).
    \item $\Delta\mathbf{p}_{mqk}$ (sampling \textbf{offsets}) and $A_{mqk}$ (\textbf{attention weights}, $\textstyle\sum_k A_{mqk}{=}1$) are both \emph{predicted from} $\mathbf{z}_q$ by linear layers followed by softmax.
    \item Indices: $m{=}1{\dots}M$ heads, $k{=}1{\dots}K$ sampled points ($K{=}4$); \;$W_m, W'_m$: output / value projection matrices.
  \end{itemize}
\end{frame}

% -----------------------------------------------------------------------
\begin{frame}{Deformable DETR: Multi-Scale Features}

  \begin{columns}
    \begin{column}{0.55\textwidth}
      \textbf{Multi-scale deformable attention}:
      \begin{align*}
        \operatorname{MSDeformAttn} = \sum_{m=1}^{M} W_m \biggl[ & \sum_{l=1}^{L}\sum_{k=1}^{K} A_{mlqk} \cdot W'_m \, \\
        & \mathbf{x}^l(\phi_l(\hat{\mathbf{p}}_q) + \Delta\mathbf{p}_{mlqk})\biggr]
      \end{align*}
      \begin{itemize}
        \item $\{\mathbf{x}^l\}_{l=1}^{L}$: the $L{=}4$ \textbf{feature-map levels} (scales); $l$ indexes the level.
        \item $\hat{\mathbf{p}}_q\in[0,1]^2$: \textbf{normalized} reference point; $\phi_l$ \textbf{rescales} it to level $l$'s grid.
        \item $K$ points sampled \emph{per level} ($L{\cdot}K$ total); weights normalized over levels \emph{and} points: $\sum_l\sum_k A_{mlqk}{=}1$.
      \end{itemize}

      \vspace{0.3em}
      \textbf{Additional features}:
      \begin{itemize}
        \item \textbf{Iterative box refinement}: each decoder layer refines the box from the previous layer
        \item \textbf{Two-stage variant}: encoder produces proposals as initial reference points
      \end{itemize}
    \end{column}
    \begin{column}{0.42\textwidth}
      \begin{center}
        \begin{tikzpicture}[scale=0.92, font=\scriptsize, >=Stealth]
          % feature pyramid: same image, coarser (smaller) maps higher up; ref pt at rel (0.86,0.40)
          % maps centered on x=1.1; wide vertical gaps for centered downsample arrows
          % --- 1/8 (finest, largest), bottom ---
          \begin{scope}\clip (0.25,0.0) rectangle (1.95,1.7);
            \begin{scope}[shift={(0.25,0.0)}, scale=0.85, opacity=0.5]\visionSketchScene\end{scope}\end{scope}
          \foreach \g in {1,...,7}{\draw[gray!40] (0.25+\g*0.2125,0)--(0.25+\g*0.2125,1.7);
            \draw[gray!40] (0.25,\g*0.2125)--(1.95,\g*0.2125);}
          \draw[thick] (0.25,0) rectangle (1.95,1.7);
          \node[anchor=east] at (0.18,0.85) {$\tfrac{1}{8}$};
          \draw[blue!70!black,->] (1.712,0.68)--(1.40,1.12); \fill[blue!70!black] (1.40,1.12) circle (1.3pt);
          \draw[blue!70!black,->] (1.712,0.68)--(1.86,0.34); \fill[blue!70!black] (1.86,0.34) circle (1.3pt);
          \fill[red] (1.712,0.68) circle (1.8pt);
          % --- 1/16 (medium) ---
          \begin{scope}\clip (0.5,2.3) rectangle (1.7,3.5);
            \begin{scope}[shift={(0.5,2.3)}, scale=0.6, opacity=0.5]\visionSketchScene\end{scope}\end{scope}
          \foreach \g in {1,2,3}{\draw[gray!50] (0.5+\g*0.3,2.3)--(0.5+\g*0.3,3.5);
            \draw[gray!50] (0.5,2.3+\g*0.3)--(1.7,2.3+\g*0.3);}
          \draw[thick] (0.5,2.3) rectangle (1.7,3.5);
          \node[anchor=east] at (0.43,2.9) {$\tfrac{1}{16}$};
          \draw[blue!70!black,->] (1.532,2.78)--(1.18,3.22); \fill[blue!70!black] (1.18,3.22) circle (1.3pt);
          \draw[blue!70!black,->] (1.532,2.78)--(1.60,2.48); \fill[blue!70!black] (1.60,2.48) circle (1.3pt);
          \fill[red] (1.532,2.78) circle (1.8pt);
          % --- 1/32 (coarsest, smallest), top ---
          \begin{scope}\clip (0.7,4.1) rectangle (1.5,4.9);
            \begin{scope}[shift={(0.7,4.1)}, scale=0.4, opacity=0.5]\visionSketchScene\end{scope}\end{scope}
          \draw[gray!55] (1.1,4.1)--(1.1,4.9); \draw[gray!55] (0.7,4.5)--(1.5,4.5);
          \draw[thick] (0.7,4.1) rectangle (1.5,4.9);
          \node[anchor=east] at (0.63,4.5) {$\tfrac{1}{32}$};
          \draw[blue!70!black,->] (1.388,4.42)--(1.00,4.66); \fill[blue!70!black] (1.00,4.66) circle (1.3pt);
          \draw[blue!70!black,->] (1.388,4.42)--(1.28,4.20); \fill[blue!70!black] (1.28,4.20) circle (1.3pt);
          \fill[red] (1.388,4.42) circle (1.8pt);
          % --- downsampling arrows: centered (x=1.1), vertical, in the gaps ---
          \draw[orange!85!black, ->, line width=1.1pt] (1.1,1.78)--(1.1,2.26);
          \draw[orange!85!black, ->, line width=1.1pt] (1.1,3.58)--(1.1,4.06);
          \node[orange!85!black, anchor=west, font=\scriptsize] at (1.2,2.02) {$\times\tfrac12$};
          \node[orange!85!black, anchor=west, font=\scriptsize] at (1.2,3.82) {$\times\tfrac12$};
          % legend
          \node[anchor=west, red, font=\scriptsize] at (2.05, 4.42) {same ref.\ $\hat{\mathbf p}_q$};
          \node[anchor=west, blue!70!black, font=\scriptsize, align=left] at (2.05, 2.9) {$K$ sampled\\pts / level};
          \node[anchor=west, orange!85!black, font=\scriptsize, align=left] at (2.05, 0.85) {downsample\\image $\times 2$};
        \end{tikzpicture}
      \end{center}
      \vspace{-0.2em}
      {\footnotesize Coarser levels are the image \textbf{downsampled} (orange). The reference point is the same image location at every scale; the $K$ offsets (hence sampled points) are learned \emph{per level}.}
    \end{column}
  \end{columns}
\end{frame}

% -----------------------------------------------------------------------
\begin{frame}{Deformable DETR: Results \& Training Details}

  \begin{columns}
    \begin{column}{0.46\textwidth}
      \textbf{Training recipe} (COCO, ResNet-50):
      \begin{itemize}
        \item \textbf{50 epochs}, lr $\times 0.1$ at epoch 40 ($\approx 10\times$ fewer than DETR's 500)
        \item Adam, lr $2{\times}10^{-4}$ (backbone $2{\times}10^{-5}$), weight decay $10^{-4}$, batch $32$
        \item \textbf{Focal loss} for classification (DETR used cross-entropy)
        \item $4$ levels from backbone stages C3--C5 $+$ one extra stride-64 map
      \end{itemize}

      \vspace{0.3em}
      \textbf{Hyperparameters}:
      \begin{description}
        \item[$M{=}8$] attention heads
        \item[$L{=}4$] feature levels
        \item[$K{=}4$] sampling points / level
        \item[$300$] object queries
      \end{description}

      \vspace{0.3em}
      \textbf{Why it converges fast}: sparse, \emph{localized} sampling around a reference point gives a strong spatial prior, so attention need not learn to localize from scratch.
    \end{column}
    \begin{column}{0.52\textwidth}
      \begin{center}
        \small
        \begin{tabular}{lccccc}
          \toprule
          \textbf{Model} & \textbf{Ep.} & \textbf{AP} & \textbf{AP$_S$} & \textbf{AP$_L$} \\
          \midrule
          DETR-DC5        & 500 & 43.3 & 22.5 & 61.1 \\
          DETR-DC5        &  50 & 35.3 & 15.2 & 53.6 \\
          \midrule
          Deformable      &  50 & 43.8 & 26.4 & 59.0 \\
          \;+ box refine  &  50 & 45.4 & 26.8 & 61.7 \\
          \;+ two-stage   &  50 & \textbf{46.2} & \textbf{28.8} & \textbf{61.7} \\
          \bottomrule
        \end{tabular}
      \end{center}
      \vspace{-0.2em}
      {\footnotesize COCO val2017, all ResNet-50. AP$_S$/AP$_L$: small / large objects.}

      \vspace{0.6em}
      \textbf{Takeaways}:
      \begin{itemize}
        \item \textbf{$10\times$ faster training}, yet higher AP than 500-epoch DETR
        \item Largest gain on \textbf{small objects} (AP$_S$ $22.5\!\to\!28.8$): multi-scale features fix DETR's main weakness
        \item Linear-in-tokens attention ($\mathcal{O}(N_q\,L\,K)$) makes high-resolution multi-scale input affordable
        \item Iterative refinement \& two-stage proposals each add $\sim\!1$--$2$ AP
      \end{itemize}
    \end{column}
  \end{columns}

  \vspace{0.3em}
  {\scriptsize Source: \textcite{zhu2021deformabledetr}; baseline \textcite{carion2020detr}.}
\end{frame}

% =======================================================================
\subsection{TrackFormer}

% -----------------------------------------------------------------------
\begin{frame}{The Multi-Object Tracking Problem}

  \textbf{Task}: in a video, detect every object instance \emph{and} assign each a \textbf{consistent identity} (track) across frames. Output a set of trajectories $\{(c_i, b_i^t)\}_t$.

  \vspace{0.2em}
  \begin{center}
    \pgfdeclarelayer{fglinks}\pgfsetlayers{main,fglinks}%
    \resizebox{0.82\textwidth}{!}{%
      \begin{tikzpicture}[arr/.style={-Stealth, thick},
        person/.pic={
          \begin{scope}[line cap=round, line join=round]
            \fill (0,0.85) circle (0.13);                    % head
            \draw[line width=0.9pt] (0,0.72) -- (0,0.32);    % torso
            \draw[line width=0.9pt] (0,0.62) -- (-0.18,0.40); % arms
            \draw[line width=0.9pt] (0,0.62) -- ( 0.18,0.40);
            \draw[line width=0.9pt] (0,0.32) -- (-0.15,0.0); % legs
            \draw[line width=0.9pt] (0,0.32) -- ( 0.15,0.0);
          \end{scope}}]
        % Three mini frames showing two objects moving + one appearing
        \foreach \fr/\dx/\ax/\bx in {t/0/0.35/1.3, {t{+}1}/3.0/0.6/1.05, {t{+}2}/6.0/0.9/0.8} {
          \begin{scope}[shift={(\dx,0)}]
            % scene background: sky, sun, clouds, skyline, ground, road
            \begin{scope}
              \clip (0,0) rectangle (2,1.5);
              \fill[blue!12] (0,0) rectangle (2,1.5);              % sky
              % sun
              \fill[yellow!55!orange!80] (1.7,1.28) circle (0.12);
              % clouds
              \fill[white] (0.35,1.22) circle (0.09) (0.48,1.25) circle (0.11) (0.62,1.22) circle (0.09);
              \fill[white] (1.15,1.32) circle (0.07) (1.26,1.34) circle (0.09) (1.37,1.32) circle (0.07);
              % distant buildings (skyline)
              \fill[gray!22] (0.18,0.45) rectangle (0.55,1.02);    % building 1
              \fill[gray!16] (0.62,0.45) rectangle (0.95,0.82);    % building 2
              \fill[gray!30] (1.00,0.45) rectangle (1.30,1.10);    % building 3
              \fill[gray!28] (1.45,0.45) rectangle (1.82,1.18);    % building 4
              % windows on buildings
              \foreach \wy in {0.58,0.74,0.90} {
                \fill[yellow!40!gray!40] (0.26,\wy) rectangle (0.32,\wy+0.05);
                \fill[yellow!40!gray!40] (0.41,\wy) rectangle (0.47,\wy+0.05);
              }
              \foreach \wy in {0.58,0.78,0.98} {
                \fill[yellow!40!gray!40] (1.53,\wy) rectangle (1.59,\wy+0.05);
                \fill[yellow!40!gray!40] (1.68,\wy) rectangle (1.74,\wy+0.05);
              }
              % trees
              \fill[brown!60!black] (0.06,0.30) rectangle (0.10,0.50);
              \fill[green!45!black] (0.08,0.55) circle (0.11);
              \fill[brown!60!black] (1.90,0.30) rectangle (1.94,0.50);
              \fill[green!45!black] (1.92,0.55) circle (0.10);
              % ground
              \fill[green!10!gray!18] (0,0) rectangle (2,0.45);    % ground
              \draw[gray!45, line width=0.4pt] (0,0.45) -- (2,0.45); % horizon
            \end{scope}
            \draw[thick] (0,0) rectangle (2,1.5);
            \node[font=\scriptsize, above] at (1.0, 1.5) {Frame $\fr$};
            % object A (red), moves right: person inside its box
            \pic[red!70!black, scale=0.76] at (\ax+0.225,0.2) {person};
            \draw[very thick, red!70!black] (\ax,0.2) rectangle (\ax+0.45,0.95);
            \node[fill=red!70!black, text=white, font=\tiny, inner sep=0.5pt] at (\ax+0.22,1.05) {1};
            % object B (green), moves left: person inside its box
            \pic[green!55!black, scale=0.76] at (\bx+0.225,0.45) {person};
            \draw[very thick, green!55!black] (\bx,0.45) rectangle (\bx+0.45,1.2);
            \node[fill=green!55!black, text=white, font=\tiny, inner sep=0.5pt] at (\bx+0.22,1.3) {2};
          \end{scope}
        }
        % object C appears at t+2 (violet): person inside its box
        \pic[violet, scale=0.56] at (6.0+1.67,0.15) {person};
        \draw[very thick, violet] (6.0+1.5,0.15) rectangle (6.0+1.85,0.7);
        \node[fill=violet, text=white, font=\tiny, inner sep=0.5pt] at (6.0+1.67,0.8) {3};
        % Identity-preserving links, drawn on a foreground layer so they overlay the
        % frames (never hidden behind them) and each with a white casing for contrast.
        % Each link starts ON the tracked object's box edge in the source frame and
        % ends ON the same object's box edge in the next frame. Object 1 (red) runs
        % along the lower body height (below object 2's box); object 2 (green) runs
        % along the upper body height (above object 1's box), so neither link ever
        % crosses the other object.
        \begin{pgfonlayer}{fglinks}
          \foreach \sx/\sy/\ex/\ey/\col in {%
              0.8/0.32/3.6/0.32/red!70!black, 4.05/0.32/6.9/0.32/red!70!black,
              1.75/1.05/4.05/1.05/green!55!black, 4.5/1.05/6.8/1.05/green!55!black} {
            \draw[white, line width=2.6pt, line cap=butt] (\sx,\sy) -- (\ex,\ey);
            \draw[arr, \col, dashed, very thick] (\sx,\sy) -- (\ex,\ey);
          }
        \end{pgfonlayer}
      \end{tikzpicture}}
  \end{center}
  \vspace{-0.1em}
  \begin{center}{\footnotesize Same color/ID across frames; object \textbf{3} is a new appearance.}\end{center}

  \vspace{0.2em}
  \begin{columns}[T]
    \begin{column}{0.48\textwidth}
      \textbf{Beyond per-frame detection}
      \begin{itemize}
        \item Detection alone gives boxes per frame with \emph{no identity link}.
        \item Must handle objects that \textbf{appear}, \textbf{disappear}, and \textbf{reappear} (occlusion).
        \item Avoid \textbf{identity switches} when tracks cross or overlap.
      \end{itemize}
    \end{column}
    \begin{column}{0.48\textwidth}
      \textbf{Classic ``tracking-by-detection''} \;{\footnotesize\cite{bewley2016sort,wojke2017deepsort}}
      \begin{itemize}
        \item Run a detector independently per frame.
        \item \textbf{Associate} boxes across frames: motion model (Kalman) $+$ appearance features.
        \item Solve a bipartite assignment (Hungarian) on an IoU / similarity cost.
      \end{itemize}
      \vspace{0.1em}
      {\footnotesize Separate stages, hand-crafted matching. Not end-to-end.}
    \end{column}
  \end{columns}

  \vspace{0.3em}
  \begin{center}
    \footnotesize\itshape TrackFormer's idea: fold association \emph{into} the detector by letting each object's decoder embedding persist as a query into the next frame.
  \end{center}
\end{frame}

% -----------------------------------------------------------------------
\begin{frame}[fragile]{TrackFormer: Multi-Object Tracking with Transformers}

  \cite{meinhardt2022trackformer}: Extend DETR to video by introducing \textbf{track queries} that persist across frames: tracking by attention.

  \vspace{0.2em}
  \begin{center}
    \begin{tikzpicture}[
      scale=0.61, every node/.style={scale=0.61},
      box/.style={draw, rounded corners, fill=#1, minimum width=1.6cm, minimum height=0.5cm, align=center, font=\small},
      qbox/.style={draw=#1, fill=#1, fill opacity=0.20, text opacity=1, thick,
                   minimum width=0.42cm, minimum height=0.32cm, inner sep=1pt, font=\tiny},
      % carried-over track query: identity colour + violet hatch + violet frame
      tqbox/.style={draw=violet, line width=0.7pt, fill=#1, fill opacity=0.22, text opacity=1,
                    minimum width=0.34cm, minimum height=0.32cm, inner sep=1pt, font=\tiny,
                    postaction={pattern=north east lines, pattern color=violet!65}},
      % output embedding produced by a track query: same violet hatch, output-box width
      otqbox/.style={draw=violet, line width=0.7pt, fill=#1, fill opacity=0.20, text opacity=1,
                     minimum width=0.42cm, minimum height=0.32cm, inner sep=1pt, font=\tiny,
                     postaction={pattern=north east lines, pattern color=violet!65}},
      arr/.style={-Stealth, thick},
    ]
      % One frame = a DETR pipeline. Decoder anchored at local (1.4,0).
      \def\rowpipeline#1{%
        \node[inner sep=0] (img#1) at (-6.7,0) {};
        \draw[thick, fill=orange!15] (-4.9,0.7) -- (-3.6,0.42) -- (-3.6,-0.42) -- (-4.9,-0.7) -- cycle;
        \node[font=\tiny] at (-4.25,0) {CNN};
        \node[box=orange!8, minimum width=0.8cm, font=\tiny] (conv#1) at (-2.7,0) {$1{\times}1$};
        \node[font=\tiny, text=orange!70!black] at (-2.7,0.62) {$+\text{pos}$};
        \node[box=red!12, minimum width=1.5cm, minimum height=1.25cm] (enc#1) at (-0.9,0) {};
        \node[font=\tiny, text=red!60!black] at (-0.9,0.18) {\textbf{Encoder}};
        \node[box=blue!12, minimum width=1.5cm, minimum height=1.25cm] (dec#1) at (1.4,0) {};
        \node[font=\tiny, text=blue!60!black] at (1.4,0.18) {\textbf{Decoder}};
        \node[draw, rounded corners, fill=gray!12, minimum width=0.85cm, minimum height=2.2cm,
              align=center, font=\tiny] (heads#1) at (4.2,0) {Class\\FFN\\[2pt]Box\\MLP};
        \node[inner sep=0] (res#1) at (6.2,0) {};
        \draw[arr] (img#1.east) -- (-4.9,0);
        \draw[arr] (-3.6,0) -- (conv#1.west);
        \draw[arr] (conv#1) -- (enc#1);
        \draw[arr] (enc#1) -- node[above,font=\tiny]{mem} (dec#1);
        \draw[arr] (heads#1.east) -- (res#1.west);
      }
      % decoder -> output embeddings -> heads as a DETR-style fan (nothing pierces the stack)
      \def\outfan#1#2{%
        \foreach \i in {0,...,#2} {
          \draw[arr, thin, gray!55!black] (dec#1.east) -- (o#1\i.west);
          \draw[arr, thin, gray!55!black] (o#1\i.east) -- (heads#1.west |- o#1\i);
        }}

      % ===================== Frame t =====================
      \begin{scope}[shift={(0,1.6)}]
        \rowpipeline{t}
        \node[font=\small\bfseries] at (-6.7,1.55) {Frame $t$};
        \node[inner sep=0] at (imgt) {\begin{tikzpicture}[scale=0.62]
            \begin{scope}\clip(0,0)rectangle(2,2);\motBG
              \motFig{0.55}{blue!70}\motFig{1.25}{orange!90!black}\end{scope}
            \draw[thick](0,0)rectangle(2,2);\end{tikzpicture}};
        % decoder query set = the fixed pool of object queries (no tracks yet at frame $t$)
        \foreach \i/\c in {0/blue!70,1/orange!90!black,2/teal,3/magenta!80,4/brown!70!black}
          \node[qbox=\c, minimum width=0.34cm] (oqt\i) at (\i*0.4,-1.5) {};
        \draw[gray!55!black] (-0.18,-1.31) -- (1.78,-1.31);
        \draw[arr,gray!55!black] (0.8,-1.31) -- (dect.south);
        \node[font=\tiny,text=green!45!black] at (0.8,-2.0) {object queries (new)};
        \node[font=\tiny,text=gray,align=center] at (2.45,-1.62) {\itshape no track\\\itshape queries yet};
        % output embeddings: two queries detect A,B; the others emit no-object ($\varnothing$)
        \foreach \i/\c/\l in {0/blue!70/{},1/orange!90!black/{},2/gray!55/$\varnothing$,3/gray!55/$\varnothing$,4/gray!55/$\varnothing$}
          \node[qbox=\c,text=black] (ot\i) at (2.75,0.9-\i*0.45) {\l};
        \node[font=\tiny,text=gray!55!black,align=center] at (2.75,1.32) {output\\embeddings};
        \outfan{t}{4}
        \node[inner sep=0] at (rest) {\begin{tikzpicture}[scale=0.62]
            \begin{scope}\clip(0,0)rectangle(2,2);\motBG
              \motFig{0.55}{blue!70}\motFig{1.25}{orange!90!black}\end{scope}
            \motFigBox{0.55}{blue!70}\motFigBox{1.25}{orange!90!black}
            \draw[thick](0,0)rectangle(2,2);\end{tikzpicture}};
        \node[font=\tiny] at (6.2,-1.3) {detect $A,B$};
      \end{scope}

      % ===================== Frame t+1 =====================
      \begin{scope}[shift={(0,-1.6)}]
        \rowpipeline{u}
        \node[font=\small\bfseries] at (-6.7,1.55) {Frame $t{+}1$};
        \node[inner sep=0] at (imgu) {\begin{tikzpicture}[scale=0.62]
            \begin{scope}\clip(0,0)rectangle(2,2);\motBG
              \motFig{0.3}{magenta!80}\motFig{0.78}{blue!70}\motFig{1.5}{orange!90!black}\end{scope}
            \draw[thick](0,0)rectangle(2,2);\end{tikzpicture}};
        % decoder query set = object queries (new) CONCATENATED with track queries copied from t
        \foreach \i/\c in {0/blue!70,1/orange!90!black,2/teal,3/magenta!80,4/brown!70!black}
          \node[qbox=\c, minimum width=0.34cm] (oqu\i) at (\i*0.4,-1.5) {};
        \node[font=\footnotesize,text=violet] at (1.875,-1.5) {$\oplus$};
        \node[tqbox=blue!70] (tqA) at (2.15,-1.5) {};
        \node[tqbox=orange!90!black] (tqB) at (2.55,-1.5) {};
        \draw[gray!55!black] (-0.18,-1.31) -- (2.73,-1.31);
        \draw[arr,gray!55!black] (1.275,-1.31) -- (decu.south);
        \node[font=\tiny,text=green!45!black] at (0.8,-2.0) {object q (new)};
        \node[font=\tiny,text=violet] at (2.35,-2.0) {track q (from $t$)};
        % 7 outputs, one per input query, IN THE SAME ORDER as the concatenated query set:
        % the 5 object queries first, then the 2 track queries. Of the object queries only
        % the magenta one detects the new object C; the others emit no-object ($\varnothing$).
        % The two track-query outputs carry the violet hatch (continued tracks, same as inputs).
        \foreach \i/\c/\l in {0/gray!55/$\varnothing$,1/gray!55/$\varnothing$,2/gray!55/$\varnothing$,3/magenta!80/{},4/gray!55/$\varnothing$}
          \node[qbox=\c,text=black] (ou\i) at (2.75,1.05-\i*0.35) {\l};
        \node[otqbox=blue!70]         (ou5) at (2.75,1.05-5*0.35) {};
        \node[otqbox=orange!90!black] (ou6) at (2.75,1.05-6*0.35) {};
        \outfan{u}{6}
        \node[inner sep=0] at (resu) {\begin{tikzpicture}[scale=0.62]
            \begin{scope}\clip(0,0)rectangle(2,2);\motBG
              \motFig{0.3}{magenta!80}\motFig{0.78}{blue!70}\motFig{1.5}{orange!90!black}\end{scope}
            \motFigBoxTrack{0.78}{blue!70}\motFigBoxTrack{1.5}{orange!90!black}\motFigBox{0.3}{magenta!80}
            \draw[thick](0,0)rectangle(2,2);\end{tikzpicture}};
        \node[font=\tiny] at (6.2,-1.3) {track $A,B$ + new $C$};
      \end{scope}

      % ===================== temporal recurrence =====================
      % Each detected frame-t output re-enters as the SAME-identity track query of
      % frame t+1's decoder query set; the no-object ($\varnothing$) outputs are dropped.
      % Identity-coloured feedback: output A -> track query A, output B -> track query B.
      \draw[arr,blue!70,very thick,dashed]
         (ot0.east) -- (5.35,2.5) -- (5.35,-4.2) -- (2.15,-4.2) -- (tqA.south);
      \draw[arr,orange!90!black,very thick,dashed]
         (ot1.east) -- (5.0,2.05) -- (5.0,-3.95) -- (2.55,-3.95) -- (tqB.south);
    \end{tikzpicture}
  \end{center}

  \begin{itemize}
    \item \textbf{Object queries}: detect \emph{new} objects (as in DETR).
    \item \textbf{Track queries}: previous-frame output embeddings $\to$ track \emph{existing} objects.
    \item Decoder self-attention over the joint set avoids duplicates / identity switches.
    \item Track terminated when classified $\varnothing$; revivable from its stored embedding.
  \end{itemize}
\end{frame}

% -----------------------------------------------------------------------
\begin{frame}[fragile]{TrackFormer: Track Queries in Detail}

  The decoder query set at frame $t$ \textbf{concatenates} a fixed pool of object queries with the track queries carried over from $t{-}1$.

  \vspace{0.2em}
  \begin{columns}[T]
    \begin{column}{0.46\textwidth}
      \begin{center}\textbf{\footnotesize Track life cycle across frames}\end{center}
      \vspace{-0.1em}
      \begin{center}
        \resizebox{\linewidth}{!}{%
        \begin{tikzpicture}[
            x=1.15cm, y=1.0cm,
            every node/.style={font=\scriptsize},
            trk/.style={line width=2.4pt, #1},
            ghost/.style={line width=1.2pt, #1, dashed, dash pattern=on 2pt off 2pt},
            born/.style={circle, fill=#1, draw=white, line width=0.4pt, inner sep=1.6pt},
            xdead/.style={red!80!black, line width=1.2pt},
          ]
          % ---- frame grid + header ----
          \foreach \f/\lab in {0/$t$,1/$t{+}1$,2/$t{+}2$,3/$t{+}3$,4/$t{+}4$,5/$t{+}5$} {
            \draw[gray!22, line width=0.4pt] (\f,-0.55) -- (\f,3.05);
            \node[gray!55!black, font=\tiny] at (\f,3.3) {\lab};
          }
          \node[gray!55!black, font=\tiny, anchor=east] at (-0.55,3.3) {frame:};
          % ============ Track A (blue): one-frame break, then revived ============
          \node[blue!70!black, anchor=east] at (-0.35,2.4) {\bfseries A};
          \node[born=blue!70] at (0,2.4) {};
          \draw[trk=blue!70] (0,2.4) -- (2,2.4);
          \draw[ghost=blue!70] (2,2.4) -- (4,2.4);
          \draw[trk=blue!70] (4,2.4) -- (5,2.4);
          \foreach \f in {1,2,4,5} \fill[blue!70] (\f,2.4) circle (1.4pt);
          \node[blue!70!black, font=\tiny, fill=white, inner sep=0.5pt] at (3,2.4) {$\varnothing$};
          % ============ Track B (orange): occluded then terminated ============
          \node[orange!85!black, anchor=east] at (-0.35,1.2) {\bfseries B};
          \node[born=orange!90!black] at (0,1.2) {};
          \draw[trk=orange!90!black] (0,1.2) -- (2,1.2);
          \foreach \f in {1,2} \fill[orange!90!black] (\f,1.2) circle (1.4pt);
          \draw[ghost=orange!90!black] (2,1.2) -- (4,1.2);
          \foreach \f in {3,4} \node[orange!80!black, font=\tiny, fill=white, inner sep=0.5pt] at (\f,1.2) {$\varnothing$};
          \draw[xdead] (5-0.11,1.2-0.11) -- (5+0.11,1.2+0.11);
          \draw[xdead] (5-0.11,1.2+0.11) -- (5+0.11,1.2-0.11);
          \draw[decorate, decoration={brace, amplitude=3pt}, gray!60!black]
                (2.05,1.5) -- (3.95,1.5);
          \node[gray!55!black, font=\tiny, align=center] at (3,1.9) {$\varnothing$ for $\le K$ frames};
          % ============ Track C (magenta): appears mid-sequence ============
          \node[magenta!80!black, anchor=east] at (-0.35,0) {\bfseries C};
          \node[born=magenta!80] at (2,0) {};
          \draw[trk=magenta!80] (2,0) -- (5,0);
          \foreach \f in {3,4,5} \fill[magenta!80] (\f,0) circle (1.4pt);
          % ---- birth annotations (object query) ----
          \draw[-Stealth, green!45!black, line width=0.6pt] (0,2.72) -- (0,2.5);
          \node[green!45!black, font=\tiny, anchor=south] at (0.05,2.72) {spawned by object query};
          \draw[-Stealth, green!45!black, line width=0.6pt] (2,0.32) -- (2,0.1);
          \node[green!45!black, font=\tiny, anchor=south] at (2.05,0.28) {new object $\to$ new track};
          % ==================== legend ====================
          \begin{scope}[shift={(-0.35,-1.35)}]
            \node[born=gray!60] at (0,0) {};
            \node[anchor=west, font=\tiny] at (0.15,0) {birth (object query)};
            \draw[trk=gray!55!black] (2.55,0) -- (3.05,0);
            \node[anchor=west, font=\tiny] at (3.1,0) {tracked (track query)};
            \draw[ghost=gray!55!black] (0,-0.5) -- (0.5,-0.5);
            \node[anchor=west, font=\tiny] at (0.6,-0.5) {$\varnothing$ but kept ($\le K$)};
            \draw[xdead] (2.65,-0.11-0.5) -- (2.87,0.11-0.5);
            \draw[xdead] (2.65,0.11-0.5) -- (2.87,-0.11-0.5);
            \node[anchor=west, font=\tiny] at (3.1,-0.5) {terminated};
          \end{scope}
        \end{tikzpicture}}
      \end{center}
      \vspace{-0.1em}
      {\footnotesize Each identity is born from an object query, then carried frame-to-frame by \emph{its own} track query, bridging short $\varnothing$ gaps until terminated.}
    \end{column}
    \begin{column}{0.52\textwidth}
      \textbf{Two query types, one decoder}
      \begin{itemize}
        \item \textbf{Object queries}: $N$ learned, static (DETR) $\to$ spawn \emph{new} tracks.
        \item \textbf{Track queries}: \emph{not} learned; the output embedding of an alive track, fed back next frame (autoregressive).
      \end{itemize}

      \vspace{0.2em}
      \textbf{Why it works}
      \begin{itemize}
        \item Track query carries identity \emph{and} appearance; cross-attention relocates it $\Rightarrow$ implicit association, no test-time Hungarian.
        \item Self-attention over the \emph{joint} set lets track queries suppress object queries on the same object (no duplicates).
      \end{itemize}

      \vspace{0.2em}
      \textbf{Life cycle}: object query over threshold $\to$ \emph{birth}; track scored $\varnothing$ for $K$ frames $\to$ \emph{death}.

    \end{column}
  \end{columns}
\end{frame}

% -----------------------------------------------------------------------
\begin{frame}[fragile]{TrackFormer: Training and Pre-training}

  \begin{columns}[T]
    \begin{column}{0.5\textwidth}
      \textbf{Trained on frame pairs} $(t{-}1, t)$
      \begin{itemize}
        \item Run the model on $t{-}1$; its output embeddings become the \textbf{track queries} fed into $t$, next to the static object queries.
        \item \textbf{Set loss} (DETR-style: class CE $+\,L_1 +$ GIoU). \emph{New} objects assigned by \textbf{Hungarian matching}; track queries matched by carried-over \emph{identity}.
      \end{itemize}

      \vspace{0.3em}
      \textbf{Auxiliary losses}
      \begin{itemize}
        \item The full set loss is applied after \emph{every} decoder layer ($6\times$) $\Rightarrow$ faster, more stable convergence.
      \end{itemize}

      \vspace{0.3em}
      \textbf{Track-query augmentation}
      \begin{itemize}
        \item Randomly \textbf{drop} track queries: simulate occlusion / track death.
        \item Inject \textbf{false-positive} track queries from the background: teach the model to terminate tracks and not propagate errors.
      \end{itemize}
    \end{column}
    \begin{column}{0.48\textwidth}
      \textbf{Pseudo-trajectories from single images}
      \begin{center}
        \begin{tikzpicture}[scale=0.6, font=\tiny, arr/.style={-Stealth, thick}]
          \begin{scope}\clip (0,0) rectangle (2,2);\visionSketchScene\end{scope}
          \draw[thick] (0,0) rectangle (2,2);
          % box on the person (stick figure on the right), not the house
          \draw[very thick, violet] (1.62,0.48) rectangle (1.94,1.03);
          \node[below=0pt] at (1,0) {frame $t{-}1$};
          \draw[arr] (2.2,1) -- (3.05,1);
          \node[above=0pt] at (2.62,1.0) {augment};
          \begin{scope}[shift={(3.3,0)}]
            \begin{scope}\clip (0,0) rectangle (2,2);
              % box drawn inside the augmentation transform, so it tracks the person
              \begin{scope}[shift={(-0.28,0.12)}, rotate around={-5:(1,1)}, scale=1.07]%
                \visionSketchScene
                \draw[very thick, violet] (1.62,0.48) rectangle (1.94,1.03);
              \end{scope}
            \end{scope}
            \draw[thick] (0,0) rectangle (2,2);
            \node[below=0pt] at (1,0) {synthetic frame $t$};
          \end{scope}
        \end{tikzpicture}
      \end{center}
      \begin{itemize}
        \item Apply random \textbf{spatial augmentation} (shift, scale, rotate) to a still image to synthesize an adjacent frame; the \emph{same} boxes give the identities $\Rightarrow$ a 2-frame \textbf{pseudo-trajectory}.
        \item Turns large \emph{single-image detection} sets into tracking data.
      \end{itemize}

      \vspace{0.2em}
      \textbf{Pre-training data}
      \begin{itemize}
        \item Detector pre-trained on \textbf{COCO} (DETR detection).
        \item Trained on \textbf{CrowdHuman}~\cite{shao2018crowdhuman} (stills $\to$ pseudo-trajectories) jointly with \textbf{MOT17/20}~\cite{milan2016mot16} video.
      \end{itemize}
    \end{column}
  \end{columns}
\end{frame}

% =======================================================================
\subsection{MaskFormer}

% -----------------------------------------------------------------------
\begin{frame}[fragile]{Image Segmentation: Semantic, Instance, Panoptic}

  \textbf{Segmentation} = assign a label to every pixel. Classes split into two kinds:
  \textbf{stuff} (amorphous, uncountable: sky, grass, road) and \textbf{things}
  (countable objects: person, car).

  \vspace{0.4em}
  \begin{center}
    \begin{tikzpicture}[scale=1.45, every node/.style={scale=1.45}]
      % ---- original ----
      \begin{scope}\clip (0,0) rectangle (2,2);\segSceneRaw\end{scope}
      \draw[thick] (0,0) rectangle (2,2);
      \node[font=\scriptsize, below=1pt] at (1,0) {original};
      % ---- semantic ----
      \begin{scope}[xshift=2.4cm]
        \begin{scope}\clip (0,0) rectangle (2,2);
          \segStuff
          \segPerson{0.4}{red!75}{red!75}
          \segPerson{0.85}{red!75}{red!75}
          \segPerson{1.25}{red!75}{red!75}
        \end{scope}
        \draw[thick] (0,0) rectangle (2,2);
        \node[font=\scriptsize, below=1pt] at (1,0) {semantic};
      \end{scope}
      % ---- instance ----
      \begin{scope}[xshift=4.8cm]
        \begin{scope}\clip (0,0) rectangle (2,2);
          \fill[gray!18] (0,0) rectangle (2,2);   % stuff = ignored
          \segPerson{0.4}{red!80}{red!80}
          \segPerson{0.85}{blue!75}{blue!75}
          \segPerson{1.25}{orange!90!black}{orange!90!black}
        \end{scope}
        \draw[thick] (0,0) rectangle (2,2);
        \node[font=\scriptsize, below=1pt] at (1,0) {instance};
      \end{scope}
      % ---- panoptic ----
      \begin{scope}[xshift=7.2cm]
        \begin{scope}\clip (0,0) rectangle (2,2);
          \segStuff
          \segPerson{0.4}{red!80}{red!80}
          \segPerson{0.85}{blue!75}{blue!75}
          \segPerson{1.25}{orange!90!black}{orange!90!black}
        \end{scope}
        \draw[thick] (0,0) rectangle (2,2);
        \node[font=\scriptsize, below=1pt] at (1,0) {panoptic};
      \end{scope}
    \end{tikzpicture}
  \end{center}

  \vspace{0.3em}
  \begin{description}
    \item[Semantic] one label per pixel; \emph{instances of the same class merge}
      (all three people become one ``person'' region). Stuff and things treated alike.
    \item[Instance] detect and segment each \emph{thing} separately (three distinct people);
      stuff is ignored, and pixels may be unlabeled or overlap.
    \item[Panoptic] unifies both: every pixel gets a (class, instance id). Stuff = one region
      per class, things = one region per instance. Non-overlapping and complete.
  \end{description}

  \vspace{0.2em}
  {\footnotesize MaskFormer tackles all three with a single mask-classification formulation.}
\end{frame}

% -----------------------------------------------------------------------
\begin{frame}[fragile]{MaskFormer: From Per-Pixel to Mask Classification}

  \cite{cheng2021maskformer}: Reformulate segmentation as predicting a set of \textbf{(class, binary mask)} pairs instead of assigning a class to every pixel independently.

  \vspace{0.2em}
  \begin{center}
    \begin{tikzpicture}[
      scale=0.66, every node/.style={scale=0.66},
      box/.style={draw, rounded corners, fill=#1, minimum width=1.8cm, minimum height=0.5cm, align=center, font=\small},
      qbox/.style={draw=#1, fill=#1, fill opacity=0.20, text opacity=1, thick,
                   minimum width=0.9cm, minimum height=0.34cm, inner sep=1pt, font=\tiny},
      arr/.style={-Stealth, thick},
    ]
      % --- Input image (actual scene) ---
      \node[inner sep=0] (img) at (-7.4, 0) {
        \begin{tikzpicture}[scale=0.62]
          \begin{scope}\clip (0,0) rectangle (2,2);\visionSketchScene\end{scope}
          \draw[thick] (0,0) rectangle (2,2);
        \end{tikzpicture}
      };
      \node[font=\tiny, below=0.03cm of img] {Image $H{\times}W{\times}3$};

      % --- Backbone (trapezoid funnel) ---
      \draw[thick, fill=orange!15] (-5.5, 0.85) -- (-3.9, 0.5) -- (-3.9, -0.5) -- (-5.5, -0.85) -- cycle;
      \node[font=\small] at (-4.7, 0) {Backbone};
      \node[font=\tiny, text=gray] at (-4.7, -1.12) {ResNet / Swin};
      \draw[arr] (img.east) -- (-5.5, 0);

      % --- Top branch: Pixel Decoder -> per-pixel embeddings ---
      \node[box=cyan!12, minimum width=1.9cm, minimum height=0.9cm] (pixdec) at (-1.9, 1.7) {};
      \node[font=\small, text=cyan!60!black] at (-1.9, 1.86) {\textbf{Pixel}};
      \node[font=\small, text=cyan!60!black] at (-1.9, 1.5) {\textbf{Decoder}};
      \node[box=cyan!8, minimum width=1.5cm, font=\scriptsize] (epix) at (1.0, 1.7) {$\mathcal{E}_\text{pixel}$\\$H{\times}W{\times}D$};
      % features leave the backbone, run right a little, then split to both branches
      \coordinate (bsplit) at (-3.25, 0);
      \draw[thick] (-3.9, 0) -- (bsplit);
      \fill[black] (bsplit) circle (1.4pt);
      \draw[arr] (bsplit) |- (pixdec.west);
      \draw[arr] (pixdec) -- (epix);

      % --- Bottom branch: Transformer Decoder -> segment embeddings ---
      \node[box=red!12, minimum width=2.5cm, minimum height=0.9cm] (trdec) at (-1.9, -1.7) {};
      \node[font=\small, text=red!60!black] at (-1.9, -1.54) {\textbf{Transformer}};
      \node[font=\small, text=red!60!black] at (-1.9, -1.9) {\textbf{Decoder}};
      \draw[arr] (bsplit) |- node[pos=0.7, above, font=\tiny] {} (trdec.west);

      % N object queries (n boxes), colored per slot, feeding the decoder
      \foreach \i/\c in {0/blue!70, 1/orange!90!black, 2/teal, 3/magenta!80} {
        \node[qbox=\c, minimum width=0.4cm] at (-2.65 + \i*0.5, -3.25) {};
      }
      \node[font=\tiny, text=gray!40!black] at (-1.9, -3.7) {$N$ queries (in)};
      \draw[arr, gray!55!black] (-1.9, -3.0) -- (trdec.south);

      % Segment embeddings: N colored boxes (match query colors)
      \foreach \i/\yy/\c in {0/-0.95/blue!70, 1/-1.45/orange!90!black, 2/-1.95/teal, 3/-2.45/magenta!80} {
        \node[qbox=\c] (s\i) at (1.0, \yy) {};
        \draw[arr, \c, thin] (trdec.east) -- (s\i.west);
      }
      \node[font=\scriptsize, text=red!50!black] at (1.0, -0.6) {$\mathcal{E}_\text{segment}$ $N{\times}D$};

      % --- Class MLP (from segment embeddings) ---
      \node[box=green!12, minimum width=1.3cm, minimum height=1.85cm, font=\scriptsize] (clsp) at (3.9, -1.7) {Class MLP\\$N{\times}(K{+}1)$};
      \foreach \i in {0,...,3} {\draw[arr, thin] (s\i.east) -- (clsp.west |- s\i);}

      % --- Dot product: segment x pixel -> masks ---
      \node[draw, circle, inner sep=1.5pt, font=\small] (dot) at (3.3, 0.5) {$\cdot$};
      \draw[arr] (epix.east) -| (dot);
      \draw[arr] (s0.east) .. controls (2.4,-0.95) and (2.6,0.0) .. (dot.south);
      \node[box=yellow!15, minimum width=1.5cm, font=\scriptsize] (masks) at (5.6, 0.5) {Binary Masks\\$N{\times}H{\times}W$};
      \draw[arr] (dot) -- (masks);

      % --- Result image (assembled segmentation) ---
      \node[inner sep=0] (result) at (8.2, 0) {
        \begin{tikzpicture}[scale=0.62]
          \begin{scope}\clip (0,0) rectangle (2,2);
            \visionSketchScene
            % overlay predicted masks, colored per query
            \fill[blue!70, opacity=0.30] (0,0.5) rectangle (2,2);                       % sky
            \fill[teal, opacity=0.45] (0.05,0.5) -- (0.65,0.5) -- (0.65,0.85) -- (0.35,1.55) -- (0.05,0.85) -- cycle; % tree
            \fill[orange!90!black, opacity=0.40] (0.78,0.5) rectangle (1.52,1.05);       % house body
            \fill[orange!90!black, opacity=0.40] (0.78,1.05) -- (1.52,1.05) -- (1.15,1.5) -- cycle; % roof
          \end{scope}
          \draw[thick] (0,0) rectangle (2,2);
        \end{tikzpicture}
      };
      \node[font=\tiny, below=0.03cm of result] {segmentation};
      \draw[arr] (masks.east) -- (result.west);
      \draw[arr] (clsp.east) -| node[pos=0.3, above, font=\tiny] {+ class} ([yshift=-11pt]result.south);

      % Annotation
      \node[font=\scriptsize, text=gray, align=center] at (4.45, 2.0) {$M_i(h,w) = \mathcal{E}_\text{seg}^{i\,\top} \cdot \mathcal{E}_\text{pixel}(h,w)$};
    \end{tikzpicture}
  \end{center}

  \vspace{0.2em}
  \begin{itemize}
    \item \textbf{Pixel decoder}: produces dense per-pixel embeddings $\mathcal{E}_\text{pixel} \in \mathbb{R}^{H \times W \times D}$
    \item \textbf{Transformer decoder}: $N$ queries $\to$ $N$ segment embeddings $\mathcal{E}_\text{segment} \in \mathbb{R}^{N \times D}$
    \item \textbf{Mask}: dot product $M_i = \mathcal{E}_\text{segment}^i \cdot \mathcal{E}_\text{pixel}^\top$ \quad (binary mask per query)
    \item Loss: Hungarian matching + CE (class) + BCE + Dice (mask). Unified for semantic/instance/panoptic.
  \end{itemize}
\end{frame}

% -----------------------------------------------------------------------
\begin{frame}{MaskFormer: Per-Pixel vs Mask Classification}

  Two ways to formalize segmentation, differing in \emph{what} is predicted and \emph{how} labels are assigned.

  \vspace{0.4em}
  \begin{columns}[t]
    \begin{column}{0.48\textwidth}
      \textbf{Per-pixel classification} \quad {\footnotesize (FCN, DeepLab)}
      \begin{itemize}
        \item Predict one distribution per pixel:
          \[ p_{h,w} \in \Delta^{K}, \quad \forall (h,w) \]
        \item Loss: per-pixel cross-entropy
          \[ \mathcal{L} = \sum_{h,w} \mathrm{CE}\bigl(p_{h,w},\, c^{\text{gt}}_{h,w}\bigr) \]
        \item \textbf{\#outputs} fixed to image size $H{\times}W$
        \item Cannot separate two instances of the same class $\Rightarrow$ no instance/panoptic
      \end{itemize}
    \end{column}
    \begin{column}{0.48\textwidth}
      \textbf{Mask classification} \quad {\footnotesize (DETR, MaskFormer)}
      \begin{itemize}
        \item Predict a \emph{set} of $N$ pairs:
          \[ z = \{(p_i, m_i)\}_{i=1}^{N} \]
          $p_i \in \Delta^{K+1}$ (incl.\ $\varnothing$), \; $m_i \in [0,1]^{H\times W}$
        \item Each mask groups pixels of \emph{one} segment, regardless of class
        \item $N$ decoupled from $K$ and from image size
        \item One formulation for \textbf{semantic, instance, panoptic}
      \end{itemize}
    \end{column}
  \end{columns}

\end{frame}

% -----------------------------------------------------------------------
\begin{frame}[fragile]{MaskFormer: Architecture Modules in Detail}

  Three modules turn an image into $N$ (class, mask) predictions.

  \vspace{0.2em}
  \begin{columns}[t]
    \begin{column}{0.54\textwidth}
      \begin{description}
        \item[\textcolor{orange!70!black}{1. Backbone}] ResNet / Swin $\to$ low-res features $\mathcal{F} \in \mathbb{R}^{C_\mathcal{F} \times \frac{H}{32} \times \frac{W}{32}}$.
        \item[\textcolor{cyan!60!black}{2. Pixel decoder}] FPN-style upsampler $\to$ full-res \textbf{per-pixel embeddings} $\mathcal{E}_\text{pixel} \in \mathbb{R}^{C_\mathcal{E} \times H \times W}$.
        \item[\textcolor{red!60!black}{3. Transf.\ decoder}] DETR-style: $N$ \textbf{query} embeddings cross-attend to $\mathcal{F}$ ($6$ layers) $\to$ \textbf{per-segment embeddings} $\mathcal{Q} \in \mathbb{R}^{C_\mathcal{Q} \times N}$.
      \end{description}

      \vspace{0.3em}
      {\footnotesize Two shared heads, per segment embedding $\mathcal{Q}_i$:}
      \begin{itemize}\footnotesize
        \item \textbf{Class head}: $p_i = \mathrm{softmax}(W \mathcal{Q}_i) \in \Delta^{K+1}$
        \item \textbf{Mask head}: MLP $\to \mathcal{E}_\text{mask}^i$, then\\
          $m_i[h,w] = \sigma\bigl(\mathcal{E}_\text{mask}^{i\,\top}\,\mathcal{E}_\text{pixel}[:,h,w]\bigr)$
      \end{itemize}
    \end{column}
    \begin{column}{0.44\textwidth}
      \centering
      {\footnotesize\textbf{\textcolor{cyan!60!black}{Pixel decoder}: FPN-style upsampler}}\\[3pt]
      \begin{tikzpicture}[
        scale=0.5, every node/.style={scale=0.5},
        fmap/.style={draw, rounded corners=1pt, fill=#1, minimum height=0.42cm, font=\scriptsize, inner sep=2pt},
        up/.style={-Stealth, thick},
      ]
        % bottom-up backbone features (left)
        \node[fmap=orange!18, minimum width=2.0cm]  (c2) at (0,0)   {$C_2$};
        \node[fmap=orange!18, minimum width=1.55cm] (c3) at (0,1.0) {$C_3$};
        \node[fmap=orange!18, minimum width=1.15cm] (c4) at (0,2.0) {$C_4$};
        \node[fmap=orange!18, minimum width=0.85cm] (c5) at (0,3.0) {$C_5$};
        \draw[up, gray] (c2) -- (c3); \draw[up, gray] (c3) -- (c4); \draw[up, gray] (c4) -- (c5);
        \foreach \y/\r in {0/{\frac14}, 1/{\frac18}, 2/{\frac1{16}}, 3/{\frac1{32}}} {
          \node[font=\tiny, text=gray, anchor=east] at (-1.25,\y) {$\r$};
        }
        \node[font=\tiny, text=gray, rotate=90, align=center] at (-2.15,1.5) {backbone\\(bottom-up)};
        % FPN features (right), top-down
        \node[fmap=cyan!22, minimum width=2.0cm]  (p2) at (4.2,0)   {$P_2$};
        \node[fmap=cyan!22, minimum width=1.55cm] (p3) at (4.2,1.0) {$P_3$};
        \node[fmap=cyan!22, minimum width=1.15cm] (p4) at (4.2,2.0) {$P_4$};
        \node[fmap=cyan!22, minimum width=0.85cm] (p5) at (4.2,3.0) {$P_5$};
        \foreach \c/\p in {c2/p2,c3/p3,c4/p4,c5/p5} { \draw[up] (\c.east) -- (\p.west); }
        \node[font=\tiny] at (2.1,3.5) {$1{\times}1$ conv (lateral)};
        % top-down upsample+add
        \draw[up, red!65!black] (p5) -- (p4);
        \draw[up, red!65!black] (p4) -- (p3);
        \draw[up, red!65!black] (p3) -- (p2);
        \node[font=\tiny, text=red!65!black, rotate=-90, align=center] at (5.55,1.5) {$2\times$ upsample\\$+$ add};
        % output
        \node[fmap=cyan!8, minimum width=2.4cm] (epix) at (4.2,-1.15) {$\mathcal{E}_\text{pixel}$};
        \draw[up] (p2) -- (epix);
      \end{tikzpicture}
    \end{column}
  \end{columns}

  \vspace{0.2em}
  \begin{itemize}\small
    \item Queries are \textbf{permutation-invariant}: no fixed class-to-query slot (learned by matching).
    \item Same backbone/decoder for all tasks; only the loss target and inference change.
  \end{itemize}
\end{frame}

% -----------------------------------------------------------------------
\begin{frame}{MaskFormer: Mask Classification Loss \& Bipartite Matching}

  The $N$ predictions are an unordered \emph{set}: train with DETR-style set prediction.

  \vspace{0.3em}
  \begin{description}
    \item[Pad targets] Ground-truth set $\{(c_j^{\text{gt}}, m_j^{\text{gt}})\}_{j=1}^{N_\text{gt}}$ padded with $\varnothing$ (no object) to size $N$.
    \item[Match] Find the optimal one-to-one assignment $\sigma$ via the \textbf{Hungarian algorithm}:
      \[ \sigma^\star = \arg\min_{\sigma}\ \sum_{j=1}^{N} \mathcal{L}_\text{match}\bigl(\,\hat z_{\sigma(j)},\, (c_j^{\text{gt}}, m_j^{\text{gt}})\bigr) \]
      with $\mathcal{L}_\text{match} = -\,p_{\sigma(j)}(c_j^{\text{gt}}) + \mathcal{L}_\text{mask}\bigl(m_{\sigma(j)}, m_j^{\text{gt}}\bigr)$.
  \end{description}

  \vspace{0.2em}
  \textbf{Training loss} (over the matched pairs):
  \[
    \mathcal{L} = \sum_{j=1}^{N} \Bigl[ \underbrace{-\log p_{\sigma^\star(j)}(c_j^{\text{gt}})}_{\text{classification (CE)}}
      + \mathbbm{1}_{c_j^{\text{gt}} \neq \varnothing}\;\underbrace{\mathcal{L}_\text{mask}\bigl(m_{\sigma^\star(j)}, m_j^{\text{gt}}\bigr)}_{\lambda_\text{focal}\mathcal{L}_\text{focal} + \lambda_\text{dice}\mathcal{L}_\text{dice}} \Bigr]
  \]

  \vspace{0.2em}
  \begin{itemize}
    \item Mask loss only on \emph{non}-$\varnothing$ matches; $\varnothing$ matches supervise the class head to abstain.
    \item Class CE down-weights $\varnothing$ (weight $0.1$) to counter the many ``no object'' slots.
    \item \textbf{Deep supervision}: the same loss is added after every transformer-decoder layer.
  \end{itemize}
\end{frame}

% -----------------------------------------------------------------------
\begin{frame}{MaskFormer: The Mask Loss (Focal + Dice)}

  A mask covers few pixels among mostly background $\Rightarrow$ severe \textbf{foreground/background imbalance}. Plain per-pixel BCE is dominated by easy negatives, so MaskFormer combines two complementary terms:
  \[
    \mathcal{L}_\text{mask}(m, m^{\text{gt}}) = \lambda_\text{focal}\,\mathcal{L}_\text{focal} + \lambda_\text{dice}\,\mathcal{L}_\text{dice},
    \qquad \lambda_\text{focal}=20,\ \lambda_\text{dice}=1.
  \]

  \begin{columns}[t]
    \begin{column}{0.50\textwidth}
      \textbf{Dice loss} (overlap-based)
      \[
        \mathcal{L}_\text{dice} = 1 - \frac{2\sum_i m_i g_i + \epsilon}{\sum_i m_i + \sum_i g_i + \epsilon}
      \]
      \begin{itemize}\small
        \item \textbf{Smoothness}: numerator and denominator both sum over all pixels, so the gradient flows through every prediction -- even pixels far from the boundary.
        \item \textbf{Scale-invariant}: ratio-normalised $\Rightarrow$ a tiny 10-pixel mask and a large 10k-pixel mask contribute equally; BCE would be dominated by the large one.
        \item \textbf{Imbalance-robust}: background pixels cancel in the ratio; only the overlap region and the two mask sizes matter.
      \end{itemize}
    \end{column}
    \begin{column}{0.48\textwidth}
      \textbf{Focal loss} (hard-pixel weighting)
      \[
        \mathcal{L}_\text{focal} = -\alpha\,(1{-}p_t)^{\gamma}\,\log p_t
      \]
      \begin{itemize}\small
        \item \textbf{Easy pixels} (interior, $p_t \approx 1$): $(1-p_t)^\gamma \approx 0$ $\Rightarrow$ loss nearly zeroed out; their gradients barely affect training.
        \item \textbf{Hard pixels} (boundary, $p_t \approx 0.5$): $(1-p_t)^\gamma \approx 1$ $\Rightarrow$ loss $\approx$ standard CE; gradient concentrated here.
        \item Graph: as $\gamma$ increases, the curve flattens for high $p_t$ and stays large only for small $p_t$.
      \end{itemize}
      \vspace{0.1em}
      \begin{center}
        \begin{tikzpicture}
          \begin{axis}[
            width=5.6cm, height=3.2cm,
            xlabel={\tiny $p_t$},
            ylabel={\tiny loss},
            xmin=0, xmax=1, ymin=0, ymax=2.5,
            xtick={0,0.5,1}, ytick={0,1,2},
            tick label style={font=\tiny},
            xlabel style={font=\tiny},
            ylabel style={font=\tiny},
            axis lines=left,
            grid=major,
            grid style={dotted,gray!40},
            legend pos=north east,
            legend style={font=\tiny, inner sep=1pt, row sep=-3pt},
            clip=true,
          ]
          \addplot[black, thick, dashed, domain=0.04:1, samples=80] {-ln(x)};
          \addlegendentry{CE ($\gamma{=}0$)};
          \addplot[blue!40,   thick, domain=0.04:1, samples=80] {-(1-x)^0.5*ln(x)};
          \addlegendentry{$\gamma{=}0.5$};
          \addplot[blue!65,   thick, domain=0.04:1, samples=80] {-(1-x)^1*ln(x)};
          \addlegendentry{$\gamma{=}1$};
          \addplot[blue!85!black, thick, domain=0.04:1, samples=80] {-(1-x)^2*ln(x)};
          \addlegendentry{$\gamma{=}2$};
          \addplot[blue!100!black, thick, domain=0.04:1, samples=80] {-(1-x)^5*ln(x)};
          \addlegendentry{$\gamma{=}5$};
          \end{axis}
        \end{tikzpicture}
      \end{center}
    \end{column}
  \end{columns}
\end{frame}

% -----------------------------------------------------------------------
\begin{frame}{MaskFormer: Unified Inference}

  The \emph{same} trained model produces semantic, instance, or panoptic outputs by changing only the post-processing.

  \vspace{0.4em}
  \begin{columns}[t]
    \begin{column}{0.5\textwidth}
      \textbf{Semantic} (per-pixel labels)\\[2pt]
      Marginalize over queries: no $\varnothing$, no thresholding.
      \[
        \arg\max_{c \in \{1,\dots,K\}}\ \sum_{i=1}^{N} p_i(c)\, m_i[h,w]
      \]
      \begin{itemize}
        \item Single matrix multiply $\to$ class map.
        \item Multiple queries may vote for one class.
      \end{itemize}
    \end{column}
    \begin{column}{0.5\textwidth}
      \textbf{Panoptic / instance} (per-segment)\\[2pt]
      Keep confident non-$\varnothing$ queries, then assign each pixel:
      \[
        \arg\max_{i\,:\,c_i \neq \varnothing}\ p_i(c_i)\, m_i[h,w]
      \]
      \begin{itemize}
        \item Drop low-confidence / empty masks.
        \item Merge same-class ``stuff''; keep ``thing'' instances separate.
      \end{itemize}
    \end{column}
  \end{columns}

  \vspace{0.5em}
  \begin{itemize}
    \item Semantic uses the \emph{soft} class$\times$mask product; panoptic uses a \emph{hard} per-pixel argmax over kept segments.
    \item Panoptic quality (PQ) \cite{kirillov2019panoptic} is the standard metric for the joined task.
  \end{itemize}
\end{frame}

% -----------------------------------------------------------------------
\begin{frame}{MaskFormer: Training Details \& Results}

  \begin{columns}[t]
    \begin{column}{0.50\textwidth}
      \textbf{Optimization}
      \begin{itemize}
        \item AdamW, lr $10^{-4}$, backbone lr $\times 0.1$, wd $10^{-4}$
        \item $N = 100$ queries, $6$ decoder layers $+$ deep supervision
        \item $\lambda_\text{focal} = 20$, $\lambda_\text{dice} = 1$, class CE weight $= 1$
      \end{itemize}

      \vspace{0.3em}
      \textbf{Schedules}
      \begin{itemize}
        \item \textbf{ADE20K}: $160$k iters, batch $16$, crop $512{\times}512$
        \item \textbf{COCO panoptic}: $\sim$$300$k iters, large-scale jitter
        \item Mask loss at lower resolution to save memory
      \end{itemize}
    \end{column}
    \begin{column}{0.48\textwidth}
      \textbf{One model, two tasks} (val):
      \vspace{0.3em}
      \begin{center}
        {\footnotesize
        \begin{tabular}{lcc}
          \toprule
          Model & mIoU$\uparrow$ & PQ$\uparrow$ \\
                & {\tiny ADE20K} & {\tiny COCO} \\
          \midrule
          Per-pixel baseline (R50) & $40.7$ & --- \\
          MaskFormer (R50)         & $44.5$ & $46.5$ \\
          MaskFormer (Swin-L)      & $\mathbf{55.6}$ & $\mathbf{52.7}$ \\
          \bottomrule
        \end{tabular}}
      \end{center}
      \vspace{0.2em}
      {\scriptsize mIoU: multi-scale; PQ: panoptic quality. \textcite{cheng2021maskformer}.}

      \vspace{0.4em}
      \textbf{Limitation}: standard cross-attention is slow to converge and blurs small/overlapping masks $\Rightarrow$ motivates \textbf{Mask2Former}.
    \end{column}
  \end{columns}
\end{frame}

% =======================================================================
\subsection{Mask2Former}

% -----------------------------------------------------------------------
\begin{frame}[fragile]{Mask2Former: Masked Attention for Universal Segmentation}

  \cite{cheng2022mask2former}: Improve MaskFormer with \textbf{masked cross-attention} (restrict attention to predicted mask region) and \textbf{multi-scale} decoder features.

  \vspace{0.2em}
  \begin{center}
    \begin{tikzpicture}[scale=0.62, every node/.style={scale=0.62}, >=Stealth]
      % ===== LEFT: Standard cross-attention (arrows to every pixel) =====
      \begin{scope}
        \node[font=\small\bfseries] at (2,4.55) {Standard Cross-Attention};
        \foreach \x/\y in {1.5/1.5,2.5/1.5,1.5/2.5,2.5/2.5,1.5/3.5,2.5/3.5}{
          \fill[blue!12] (\x-0.5,\y-0.5) rectangle (\x+0.5,\y+0.5);}
        \foreach \g in {1,2,3}{\draw[gray!40] (\g,0)--(\g,4); \draw[gray!40] (0,\g)--(4,\g);}
        \draw[thick] (0,0) rectangle (4,4);
        \foreach \x in {0.5,1.5,2.5,3.5}{\foreach \y in {0.5,1.5,2.5,3.5}{
          \draw[blue!55,->,opacity=0.85,line width=0.5pt] (2,2) -- (\x,\y);
          \fill[blue!60] (\x,\y) circle (1.1pt);}}
        \fill[red] (2,2) circle (2.4pt);
        \node[font=\tiny, text=blue!60!black] at (2,-0.45) {query attends to \textbf{all} $HW$ pixels};
      \end{scope}
      % ===== RIGHT: Masked cross-attention (arrows to mask region only) =====
      \begin{scope}[shift={(6,0)}]
        \node[font=\small\bfseries] at (2,4.55) {Masked Cross-Attention};
        \fill[gray!18] (0,0) rectangle (4,4);
        \foreach \x/\y in {1.5/1.5,2.5/1.5,1.5/2.5,2.5/2.5,1.5/3.5,2.5/3.5}{
          \fill[blue!25] (\x-0.5,\y-0.5) rectangle (\x+0.5,\y+0.5);}
        \foreach \g in {1,2,3}{\draw[gray!40] (\g,0)--(\g,4); \draw[gray!40] (0,\g)--(4,\g);}
        \draw[thick] (0,0) rectangle (4,4);
        \foreach \x/\y in {0.5/0.5,1.5/0.5,2.5/0.5,3.5/0.5,0.5/1.5,3.5/1.5,0.5/2.5,3.5/2.5,0.5/3.5,3.5/3.5}{
          \node[font=\scriptsize, text=red!45] at (\x,\y) {$\times$};}
        \foreach \x/\y in {1.5/1.5,2.5/1.5,1.5/2.5,2.5/2.5,1.5/3.5,2.5/3.5}{
          \draw[blue!65,->,line width=0.6pt] (2,2) -- (\x,\y);
          \fill[blue!70] (\x,\y) circle (1.1pt);}
        \fill[red] (2,2) circle (2.4pt);
        \node[font=\tiny, text=blue!60!black] at (2,-0.45) {query attends to \textbf{predicted mask} only};
      \end{scope}
      % shared legend
      \fill[red] (10.15,3.0) circle (2pt);
      \node[font=\tiny, anchor=west] at (10.35,3.0) {query};
      \fill[blue!70] (10.18,2.5) circle (1.3pt);
      \node[font=\tiny, anchor=west] at (10.35,2.5) {attended pixel};
      \fill[blue!25] (10.02,1.92) rectangle (10.22,2.12);
      \node[font=\tiny, anchor=west, align=left] at (10.35,2.0) {predicted mask};
      \node[font=\scriptsize, text=red!55] at (10.12,1.55) {$\times$};
      \node[font=\tiny, anchor=west] at (10.35,1.55) {masked out ($-\infty$)};
    \end{tikzpicture}
  \end{center}

  \vspace{0.2em}
  \textbf{Masked cross-attention} at decoder layer $\ell$ (replaces standard cross-attention):
  \[
    X_\ell = \operatorname{softmax}\!\Bigl(\mathcal{M}_{\ell-1} + \tfrac{Q_\ell K_\ell^\top}{\sqrt{d}}\Bigr) V_\ell + X_{\ell-1},
    \qquad
    \mathcal{M}_{\ell-1}(i, h, w) = \begin{cases} 0 & \sigma\!\bigl(m_{\ell-1}(i,h,w)\bigr) \geq 0.5 \\ -\infty & \text{otherwise.}\end{cases}
  \]
\end{frame}

% -----------------------------------------------------------------------
\begin{frame}{Mask2Former: What is $X_\ell$, and how is the mask predicted?}

  \[
    X_\ell = \operatorname{softmax}\!\Bigl(\mathcal{M}_{\ell-1} + \tfrac{Q_\ell K_\ell^\top}{\sqrt{d}}\Bigr) V_\ell + X_{\ell-1},
    \qquad
    \mathcal{M}_{\ell-1}(i, h, w) = \begin{cases} 0 & \sigma\!\bigl(m_{\ell-1}(i,h,w)\bigr) \geq 0.5 \\ -\infty & \text{otherwise.}\end{cases}
  \]

  \textbf{The query embeddings $X_\ell$}
  \begin{itemize}\setlength\itemsep{2pt}
    \item $X_\ell \in \mathbb{R}^{N\times d}$: the $N$ \textbf{query} (per-segment) embeddings \emph{after} decoder layer $\ell$. $X_0$ are the learnable query inputs; one query $\to$ one predicted (class, mask).
    \item $Q_\ell = f_Q(X_{\ell-1})$ projects the queries; keys/values $K_\ell, V_\ell = f_K(F_\ell), f_V(F_\ell)$ project the image features $F_\ell$ at the layer's scale ($\tfrac{1}{32},\tfrac{1}{16},\tfrac{1}{8}$, cycled).
  \end{itemize}

  \vspace{0.4em}
  \textbf{Predicting the mask $m_{\ell-1}$} (reuses MaskFormer's mask head)
  \begin{itemize}\setlength\itemsep{2pt}
    \item Map each query to a \emph{mask embedding} with a small MLP, then dot with the per-pixel embeddings $\mathcal{E}_\text{pixel} \in \mathbb{R}^{C\times H W}$ from the pixel decoder:
      \[
        m_{\ell-1} = \operatorname{MLP}(X_{\ell-1})\;\mathcal{E}_\text{pixel}^\top \;\in\; \mathbb{R}^{N\times HW}.
      \]
    \item Row $i$ holds query $i$'s mask \emph{logits} over all pixels. Bilinearly resize to the key resolution $H_\ell W_\ell$, apply $\sigma$, threshold at $0.5$.
  \end{itemize}

  \vspace{0.4em}
  \textbf{Effect of $\mathcal{M}_{\ell-1}$}
  \begin{itemize}\setlength\itemsep{2pt}
    \item Hard binary mask: query $i$ attends \emph{only} to pixels inside its own previous-layer mask; everything else gets $-\infty$ (zero weight after softmax).
    \item \textbf{Bootstrap}: $\mathcal{M}_0$ uses a mask predicted from $X_0$ before layer $1$. Masks then sharpen layer by layer, so attention stays focused and converges far faster than DETR-style global cross-attention.
  \end{itemize}
\end{frame}

% -----------------------------------------------------------------------
\begin{frame}{Mask2Former: Multi-Scale Decoder}

  \begin{columns}
    \begin{column}{0.52\textwidth}
      \begin{center}
        \begin{tikzpicture}[
          scale=0.7, every node/.style={scale=0.7},
          layer/.style={draw, rounded corners, minimum width=2.0cm, minimum height=0.4cm, align=center, font=\scriptsize},
          arr/.style={-Stealth, thick},
        ]
          % Decoder layers in rounds
          \node[layer, fill=red!12]    (l1) at (0, 6.0) {Layer 1: $\frac{1}{32}$};
          \node[layer, fill=orange!12] (l2) at (0, 5.2) {Layer 2: $\frac{1}{16}$};
          \node[layer, fill=yellow!18] (l3) at (0, 4.4) {Layer 3: $\frac{1}{8}$};

          \node[layer, fill=red!12]    (l4) at (0, 3.4) {Layer 4: $\frac{1}{32}$};
          \node[layer, fill=orange!12] (l5) at (0, 2.6) {Layer 5: $\frac{1}{16}$};
          \node[layer, fill=yellow!18] (l6) at (0, 1.8) {Layer 6: $\frac{1}{8}$};

          \node[layer, fill=red!12]    (l7) at (0, 0.8) {Layer 7: $\frac{1}{32}$};
          \node[layer, fill=orange!12] (l8) at (0, 0.0) {Layer 8: $\frac{1}{16}$};
          \node[layer, fill=yellow!18] (l9) at (0,-0.8) {Layer 9: $\frac{1}{8}$};

          % Arrows between layers
          \foreach \i/\j in {l1/l2, l2/l3, l3/l4, l4/l5, l5/l6, l6/l7, l7/l8, l8/l9} {
            \draw[arr, thin] (\i) -- (\j);
          }

          % Round braces
          \draw[decorate, decoration={brace, amplitude=4pt, mirror}, thick]
            (1.3, 4.2) -- (1.3, 6.2) node[midway, right=4pt, font=\scriptsize] {Round 1};
          \draw[decorate, decoration={brace, amplitude=4pt, mirror}, thick]
            (1.3, 1.6) -- (1.3, 3.6) node[midway, right=4pt, font=\scriptsize] {Round 2};
          \draw[decorate, decoration={brace, amplitude=4pt, mirror}, thick]
            (1.3,-1.0) -- (1.3, 1.0) node[midway, right=4pt, font=\scriptsize] {Round 3};

          % Pixel decoder box
          \node[draw, rounded corners, fill=cyan!10, minimum width=1.5cm, minimum height=2.5cm,
                font=\scriptsize, align=center] (pixdec) at (-3.5, 2.5)
                {Pixel\\Decoder\\(MSDeform\\Attn)};

          % Arrows from pixel decoder to each round
          % Round 1
          \draw[arr, red!60,          thin] (pixdec.east) ++(0, 0.8)  -- ([xshift=-0.3cm]l1.west) -- (l1.west);
          \draw[arr, orange!70,       thin] (pixdec.east) ++(0, 0)    -- ([xshift=-0.15cm]l2.west) -- (l2.west);
          \draw[arr, yellow!70!black, thin] (pixdec.east) ++(0,-0.8)  -- ([xshift=-0.3cm]l3.west) -- (l3.west);
          % Round 2 (dashed, same features reused)
          \draw[arr, red!40,          thin, dashed] (pixdec.east) ++(0, 0.8) to[out=0,in=180] (l4.west);
          \draw[arr, orange!50,       thin, dashed] (pixdec.east) ++(0, 0)   to[out=0,in=180] (l5.west);
          \draw[arr, yellow!50!black, thin, dashed] (pixdec.east) ++(0,-0.8) to[out=0,in=180] (l6.west);
          % Round 3 (dashed)
          \draw[arr, red!40,          thin, dashed] (pixdec.east) ++(0, 0.8) to[out=0,in=180] (l7.west);
          \draw[arr, orange!50,       thin, dashed] (pixdec.east) ++(0, 0)   to[out=0,in=180] (l8.west);
          \draw[arr, yellow!50!black, thin, dashed] (pixdec.east) ++(0,-0.8) to[out=0,in=180] (l9.west);

          % Legend
          \node[layer, fill=red!12,    minimum width=0.5cm, font=\tiny] at (-3.5,-0.5) {};
          \node[font=\tiny, anchor=west] at (-3.1,-0.5) {$\tfrac{1}{32}$: coarse, global};
          \node[layer, fill=orange!12, minimum width=0.5cm, font=\tiny] at (-3.5,-1.0) {};
          \node[font=\tiny, anchor=west] at (-3.1,-1.0) {$\tfrac{1}{16}$: medium};
          \node[layer, fill=yellow!18, minimum width=0.5cm, font=\tiny] at (-3.5,-1.5) {};
          \node[font=\tiny, anchor=west] at (-3.1,-1.5) {$\tfrac{1}{8}$: fine, spatial};
        \end{tikzpicture}
      \end{center}
    \end{column}
    \begin{column}{0.46\textwidth}
      \textbf{Round-robin across scales}
      \begin{itemize}
        \item \textbf{9 decoder layers} organized as \textbf{3 rounds}, each round cycling through all three resolutions.
        \item Within each round: start coarse ($\tfrac{1}{32}$, large receptive field $\to$ global object context), then refine at $\tfrac{1}{16}$, then add spatial detail at $\tfrac{1}{8}$.
        \item The pixel decoder builds the feature pyramid \emph{once} using \textbf{multi-scale deformable attention} (as in Deformable DETR~\cite{zhu2021deformabledetr}); the same maps (solid and dashed arrows) feed all 9 layers.
      \end{itemize}

      \vspace{0.4em}
      \textbf{Iterative refinement via rounds}
      \begin{itemize}
        \item After each round, updated query features carry richer representations into the next round.
        \item Masked attention uses the predicted mask from the \emph{previous} layer $\Rightarrow$ attention stays focused and sharpens progressively.
        \item 3 rounds $\approx$ 3 steps of coarse-to-fine reasoning, \emph{not} 9 independent passes.
      \end{itemize}
    \end{column}
  \end{columns}
\end{frame}

% -----------------------------------------------------------------------
\begin{frame}{Mask2Former: Results \& Training Details}

  \begin{columns}[t]
    \begin{column}{0.52\textwidth}
      \textbf{One architecture, three tasks} (vs.\ specialized SOTA):
      \begin{center}
        \footnotesize
        \begin{tabular}{llcc}
          \toprule
          \textbf{Task / data} & \textbf{Metric} & \textbf{R50} & \textbf{Swin-L} \\
          \midrule
          COCO panoptic   & PQ   & 51.9 & 57.8 \\
          COCO instance   & AP   & 43.7 & 50.1 \\
          ADE20K semantic & mIoU & 47.2 & 57.7 \\
          Cityscapes pan. & PQ   & 62.1 & 66.6 \\
          \bottomrule
        \end{tabular}
      \end{center}
      {\scriptsize Single model, single recipe per task. Beats task-specific architectures (e.g.\ Panoptic-FPN, MaskFormer) on all three.}

      \vspace{0.5em}
      \textbf{Why it works}
      \begin{itemize}\setlength\itemsep{1pt}
        \item Masked attention $\Rightarrow$ much faster convergence, sharper small/overlapping masks
        \item Multi-scale features fix MaskFormer's weak small-object performance
      \end{itemize}
    \end{column}
    \begin{column}{0.46\textwidth}
      \textbf{Training recipe}
      \begin{description}\setlength\itemsep{1pt}
        \item[Optimizer] AdamW, LR $10^{-4}$, wd $0.05$; backbone LR $\times 0.1$; step decay at $0.9,0.95$
        \item[Schedule] COCO $50$ ep, batch $16$; ADE20K $160$k iters
        \item[Augment] large-scale jittering (scale $0.1\text{--}2.0$), crop $1024^2$
        \item[Queries] $N=100$; hidden dim $256$
      \end{description}

      \vspace{0.4em}
      \textbf{Loss}: Hungarian match, then
      \[
        \mathcal{L} = \lambda_\text{cls}\mathcal{L}_\text{cls} + \lambda_\text{ce}\mathcal{L}_\text{ce} + \lambda_\text{dice}\mathcal{L}_\text{dice}
      \]
      {\footnotesize ($\lambda_\text{ce}{=}\lambda_\text{dice}{=}5$, $\lambda_\text{cls}{=}2$); deep supervision after every decoder layer.}

      \vspace{0.4em}
      \textbf{Key efficiency trick}: compute the mask loss on $K{=}12{,}544$ \textbf{sampled points} (PointRend-style), not the full mask $\Rightarrow$ $\sim 3\times$ less training memory ($18\!\to\!6$ GB).
    \end{column}
  \end{columns}

  \vspace{0.3em}
  {\scriptsize Source: \textcite{cheng2022mask2former}; deformable pixel decoder from \textcite{zhu2021deformabledetr}.}
\end{frame}

